\documentclass[11pt,oneside,openany]{book}

% 東吳大學法研所「憲法專題研究（一）」第 9 週（2026-11-09）預習問答集：
% Alexander & Schauer, 17 Const. Comment. 455 (2000)；Kramer, 100 Calif. L. Rev. 621 (2012)；
% Ferejohn & Pasquino, in Sadurski ed., Constitutional Justice, East and West 21 (2002)。
% 底稿：notes/docs/1151/seminar/W09_1109_預習.md（87808f3），tools/md_to_tex.py 轉出一次，此後 Teile 為正本。
% 版式照 ../W06_Waldron/W06_Waldron_導讀問答.tex 與 ../../課程問答/1151_憲法專題研究_課程問答.tex。
% 編譯：./build.sh。

\def\TransLaTeXDir{../../../../../1142/Translation/LaTeX}
\input{\TransLaTeXDir/course-notes.sty}
\addbibresource{refs.bib}
\CourseNotesMeta{司法至上與憲法法院的審議：第 9 週預習}{王逸帆}
\CourseNotesChinese[部]
\CourseNotesLawIndexes[noautomatic]

\providecommand{\tightlist}{\setlength{\itemsep}{0pt}\setlength{\parskip}{0pt}}
\makeatletter\def\LTcaptype{table}\makeatother

% 刊本頁：括號內文字照錄，tools/build_sources.py 依所在之部收入出處一覽
\newcommand{\pg}[1]{（#1）}
\newcommand{\en}[1]{\textit{#1}}
% 索引標記：tools/build_index.py 依 index-terms.tsv 插入，--strip 整批移除
\newif\ifshowindex \showindextrue
\newcommand{\ix}[2]{\ifshowindex\index[#1]{#2}\fi}
\setlength{\emergencystretch}{2em}
\tikzset{
  qa/.style={font=\footnotesize,every node/.append style={execute at begin node={\setstretch{1.2}}}},
  nbox/.style={draw=figGray,line width=0.4pt,fill=white,rounded corners=1.5pt,inner xsep=3pt,inner ysep=2.6pt,align=left},
  gbox/.style={nbox,fill=figGraybg},
  kbox/.style={nbox,draw=figE,fill=figEbg,line width=0.7pt},
  lbl/.style={font=\scriptsize,text=dimink,inner sep=1pt},
  hdr/.style={font=\footnotesize,text=caseink,inner sep=1pt},
  arr/.style={figGray,line width=0.6pt,-{Stealth[length=4pt,width=3.2pt]},shorten >=1pt,shorten <=1pt},
  dsh/.style={arr,dash pattern=on 2.4pt off 1.6pt},
}
\newcolumntype{L}[1]{>{\raggedright\arraybackslash}p{#1}}

\newcommand{\documentversion}{v0.1（2026-11-09）}

\begin{document}

\frontmatter
\coursetitle{司法至上與憲法法院的審議}{第 9 週預習：Alexander 與 Schauer、Kramer、Ferejohn 與 Pasquino}%
  {東吳大學法律學系研究所　憲法專題研究（一）\\115 學年度第 1 學期　第 9 週（2026 年 11 月 9 日）}%
  {\documentversion}
\tableofcontents

\chapter*{凡例}
\addcontentsline{toc}{chapter}{凡例}
% 自 notes/docs/1151/seminar/W09_1109_預習.md（87808f3）由 tools/md_to_tex.py 於 2026-11-09 轉出一次；此後本檔為正本。
本篇為課前預習，三篇的頁碼都依刊本。Alexander 與 Schauer 所用 PDF 為明尼蘇達大學法學院 Scholarship Repository 的公開本，第 1 頁是該庫封面，刊本 455 頁起於 PDF 第 2 頁；Kramer 的 PDF 第 1 頁即刊本 621 頁；Ferejohn 與 Pasquino 截自 Sadurski 所編 \en{Constitutional Justice, East and West}（Kluwer, 2002），PDF 第 1 頁即刊本 21 頁。課綱所載 Alexander 與 Schauer 的刊年 2002、Ferejohn 與 Pasquino 的刊年 2022，與刊本不符，此處依刊本作 2000 與 2002。


\mainmatter
% 自 notes/docs/1151/seminar/W09_1109_預習.md（87808f3）由 tools/md_to_tex.py 於 2026-11-09 轉出一次；此後本檔為正本。
\chapter{導論：三篇的出處與問題}
\label{ch:intro}

\section{三篇的出處與範圍}
\label{sec:intro-1}

\begin{longtable}[]{@{}
  >{\raggedright\arraybackslash}p{(\linewidth - 8\tabcolsep) * \real{0.4100}}
  >{\raggedright\arraybackslash}p{(\linewidth - 8\tabcolsep) * \real{0.2800}}
  >{\raggedright\arraybackslash}p{(\linewidth - 8\tabcolsep) * \real{0.1100}}
  >{\raggedright\arraybackslash}p{(\linewidth - 8\tabcolsep) * \real{0.1000}}
  >{\raggedright\arraybackslash}p{(\linewidth - 8\tabcolsep) * \real{0.1000}}@{}}
\caption{第 9 週的指定文獻。}\tabularnewline
\toprule\noalign{}
\begin{minipage}[b]{\linewidth}\raggedright
篇目
\end{minipage} & \begin{minipage}[b]{\linewidth}\raggedright
原出處
\end{minipage} & \begin{minipage}[b]{\linewidth}\raggedright
課綱指定頁次
\end{minipage} & \begin{minipage}[b]{\linewidth}\raggedright
整理範圍
\end{minipage} & \begin{minipage}[b]{\linewidth}\raggedright
篇幅
\end{minipage} \\
\midrule\noalign{}
\endfirsthead
\toprule\noalign{}
\begin{minipage}[b]{\linewidth}\raggedright
篇目
\end{minipage} & \begin{minipage}[b]{\linewidth}\raggedright
原出處
\end{minipage} & \begin{minipage}[b]{\linewidth}\raggedright
課綱指定頁次
\end{minipage} & \begin{minipage}[b]{\linewidth}\raggedright
整理範圍
\end{minipage} & \begin{minipage}[b]{\linewidth}\raggedright
篇幅
\end{minipage} \\
\midrule\noalign{}
\endhead
\bottomrule\noalign{}
\endlastfoot
Larry Alexander\ix{persons}{alexander@Alexander, Larry} \& Frederick Schauer\ix{persons}{schauer@Schauer, Frederick}, Defending Judicial Supremacy: A Reply & 17 Const. Comment. 455--482 (2000) & 455--473、481--482 & 455--482 全文 & 28 頁，78 註 \\
Larry D. Kramer\ix{persons}{kramer@Kramer, Larry D.}, Judicial Supremacy and the End of Judicial Restraint & 100 Calif. L. Rev. 621--634 (2012) & 621--634 & 621--634 & 14 頁，42 註 \\
John Ferejohn\ix{persons}{ferejohn@Ferejohn, John} \& Pasquale Pasquino\ix{persons}{pasquino@Pasquino, Pasquale}, Constitutional Courts as Deliberative Institutions: Towards an Institutional Theory of Constitutional Justice & in Wojciech Sadurski ed., Constitutional Justice, East and West 21--36 (2002) & 21--36 & 21--36 & 16 頁，22 註 \\
\end{longtable}

Alexander 與 Schauer 一文課綱略去的 474 至 480 頁，是第 II 部分 B 節的後半與第 III 部分的全部，第 2 部照全文整理，各該處標明屬指定頁次以外。

\section{三篇各自處理的問題}
\label{sec:intro-2}

Alexander\ix{persons}{alexander@Alexander, Larry} 與 Schauer\ix{persons}{schauer@Schauer, Frederick} 一文是 1997 年〈On Extrajudicial Constitutional Interpretation〉（110 Harv. L. Rev. 1359）的續篇，回應該文發表後三年間的批評。\footcites{alexander2000}{alexander1997}兩人主張非司法官員應把最高法院的憲法解釋當作與憲法本身同等的權威，理由是法律的核心功能之一在於定紛止爭\ix{terms}{02@\zhlidxcat{定紛止爭論}!定紛止爭}，而憲法文本未能定紛之處，需要一個能以一致的聲音發言的解釋者。第 I 部分論這個問題不能靠憲法文本或制憲史回答，第 II 部分論定紛的價值與最高法院是否為最佳的定紛者，第 III 部分論最高法院的功能是否限於裁判個案。

Kramer\ix{persons}{kramer@Kramer, Larry D.} 一文回應同卷 Posner\ix{persons}{posner@Posner, Richard A.}〈The Rise and Fall of Judicial Self-Restraint〉。\footcites{kramer2012}{posner2012}Posner 把司法自制\ix{terms}{05@\zhlidxcat{解釋理論與自制}!司法自制}的起點歸於 Thayer\ix{persons}{thayer@Thayer, James Bradley} 1893 年的文章\footcite{thayer1893}，Kramer 則追溯到建國時期共和派\ix{terms}{01@\zhlidxcat{司法審查權限之歸屬}!共和派}與聯邦派\ix{terms}{01@\zhlidxcat{司法審查權限之歸屬}!聯邦派}的兩種司法審查\ix{terms}{01@\zhlidxcat{司法審查權限之歸屬}!司法審查}理論，主張兩派都有自制的理由，而各自的理由都以法院以外另有解釋憲法的權威為前提。依 Kramer 的敘述，Brown 案\ix{cases}{02@美國!01@聯邦最高法院!1954Brown v. Board of Education（1954）@Brown v. Board of Education（1954）}之後司法至上\ix{terms}{01@\zhlidxcat{司法審查權限之歸屬}!司法至上}逐漸成為公眾與政治上的共識，自制只剩下保全法院政治資本的審慎\ix{terms}{05@\zhlidxcat{解釋理論與自制}!審慎}考量，左右兩派因此只在自己不在意結果的領域要求法院自制。

Ferejohn\ix{persons}{ferejohn@Ferejohn, John} 與 Pasquino\ix{persons}{pasquino@Pasquino, Pasquale} 比較法國、德國、義大利、西班牙與美國的法院，問的是憲法法院如何說明裁判理由。\footcite{ferejohn2002}兩人區分議會主權\ix{terms}{06@\zhlidxcat{審查制度之類型}!議會主權體制}、權力分立\ix{terms}{06@\zhlidxcat{審查制度之類型}!權力分立體制}與 Kelsen 型\ix{terms}{06@\zhlidxcat{審查制度之類型}!Kelsen 型模式}三種審查制度，認為三者各自產生不同的審議型態：Kelsen\ix{persons}{kelsen@Kelsen, Hans（凱爾森）} 型法院以單一聲音對外說理，內部審議\ix{terms}{07@\zhlidxcat{憲法審判之審議}!內部審議}密集；美國最高法院內部審議少，但以公開的方式處於更廣的外部審議\ix{terms}{07@\zhlidxcat{憲法審判之審議}!外部審議}過程的中心。

\section{與 W08 的銜接}
\label{sec:intro-3}

W08 Tushnet\ix{persons}{tushnet@Tushnet, Mark}〈Non-Judicial Review〉\footcite[註 209]{tushnet2003}的結論以 Alexander\ix{persons}{alexander@Alexander, Larry} 與 Schauer\ix{persons}{schauer@Schauer, Frederick} 為對話對象，註 209 引的是兩人 464 頁「the empirical dimension is one that cannot be avoided」一語所在的段落，Tushnet 據此主張定紛止爭\ix{terms}{02@\zhlidxcat{定紛止爭論}!定紛止爭}論所依賴的經驗前提尚未建立。Alexander 與 Schauer 458 頁註 12 則把 Tushnet 的〈Two Versions of Judicial Supremacy〉\footcite{tushnet1998}列為六篇專門回應之一，472 頁註 48 稱其反對定紛的立場最強；482 頁註 78 承認與批評者交換意見之後，辯論的核心更清楚地呈現為經驗問題，並引 Peabody\ix{persons}{peabody@Peabody, Bruce G.} 說這種研究尚未進行。兩人在 482 頁的結語，以沒有證據顯示其他機構能與最高法院同樣一致地發言為前提，維持最高法院的解釋地位；Tushnet 的結論則以定紛止爭論的經驗前提尚未建立為由，認為非司法審查\ix{terms}{01@\zhlidxcat{司法審查權限之歸屬}!司法審查}的表現不能先被排除。Kramer\ix{persons}{kramer@Kramer, Larry D.} 只在 629 頁註 30 引 Tushnet 的 \en{Taking the Constitution Away from the Courts}\footcite{tushnet1999}，用來支持立法機關內部也有認真的憲法審議；Ferejohn\ix{persons}{ferejohn@Ferejohn, John} 與 Pasquino\ix{persons}{pasquino@Pasquino, Pasquale} 未引 Tushnet。

% 自 notes/docs/1151/seminar/W09_1109_預習.md（87808f3）由 tools/md_to_tex.py 於 2026-11-09 轉出一次；此後本檔為正本。
\chapter{Alexander 與 Schauer〈Defending Judicial Supremacy: A Reply〉\pg{455 至 482 頁}}
\label{ch:AS}

\begin{longtable}[]{@{}
  >{\raggedright\arraybackslash}p{(\linewidth - 4\tabcolsep) * \real{0.1200}}
  >{\raggedright\arraybackslash}p{(\linewidth - 4\tabcolsep) * \real{0.7200}}
  >{\raggedright\arraybackslash}p{(\linewidth - 4\tabcolsep) * \real{0.1600}}@{}}
\caption{Alexander 與 Schauer 一文的節次與頁次。}\tabularnewline
\toprule\noalign{}
\begin{minipage}[b]{\linewidth}\raggedright
節次
\end{minipage} & \begin{minipage}[b]{\linewidth}\raggedright
小標
\end{minipage} & \begin{minipage}[b]{\linewidth}\raggedright
頁次
\end{minipage} \\
\midrule\noalign{}
\endfirsthead
\toprule\noalign{}
\begin{minipage}[b]{\linewidth}\raggedright
節次
\end{minipage} & \begin{minipage}[b]{\linewidth}\raggedright
小標
\end{minipage} & \begin{minipage}[b]{\linewidth}\raggedright
頁次
\end{minipage} \\
\midrule\noalign{}
\endhead
\bottomrule\noalign{}
\endlastfoot
導言 & （無小標） & 455--458 \\
I & The Nature of the Question（A. The Preconstitutional Constitution；B. The Nature of the Inquiry） & 459--467 \\
II & The Settlement Function and Its Optimal Location（A. Is Settlement Overrated?；B. Are Courts the Best Settlers?） & 467--478 \\
III & What Is the Supreme Court For? & 478--481 \\
結論 & Conclusion & 481--482 \\
\end{longtable}

\section{導言：先前論點、批評與本文目的\pg{455 至 458 頁}}
\label{sec:AS-1}

\subsection{（一）原文的主張\pg{455 頁，註 1--4}}

Alexander\ix{persons}{alexander@Alexander, Larry} 與 Schauer\ix{persons}{schauer@Schauer, Frederick} 以 1997 年發表於 \en{Harvard Law Review} 的 On Extrajudicial Constitutional Interpretation 為前文（註 1）。本文開頭重述前文立場：最高法院對憲法的解釋，應被所有其他官員，不論司法或非司法，視為與憲法本身具有相同的權威地位\pg{455 頁}。兩人引 Hughes\ix{persons}{hughes@Hughes, Charles Evans} 的話（註 2）「the Constitution is what the judges say it is」〔憲法就是法官所說的那樣〕，承認這句話作為說明對最高法院論辯意味著什麼的嘗試，可能是糟糕的法理學，因為它無從理解（註 3 引 Dworkin\ix{persons}{dworkin@Dworkin, Ronald（德沃金）}, \en{Law\textquotesingle s Empire} 與 Hart\ix{persons}{hart@Hart, H. L. A.（哈特）}, \en{The Concept of Law}）。他們仍主張，這是非司法官員對最高法院及其見解宜有的態度，其情形與下級法院法官對最高法院及其意見宜有的態度相同，只是爭議小得多（註 4 引 Caminker 與 Schauer 論先例的著作）。

\subsection{（二）批評的來源與內容\pg{456 至 457 頁，註 5--9}}

這一立場通常與最高法院的 Cooper v. Aaron\ix{cases}{02@美國!01@聯邦最高法院!1958Cooper v. Aaron（1958）@Cooper v. Aaron（1958）}（1958；註 5）及較近的 City of Boerne\ix{cases}{02@美國!01@聯邦最高法院!1997City of Boerne v. Flores（1997）@City of Boerne v. Flores（1997）} v. Flores（1997；註 6，引 521 U.S. at 536）相連，長期受到強烈批評\pg{456 頁}。批評一部分來自學者（註 7 列 Levinson\ix{persons}{levinson@Levinson, Sanford}、Gant、McConnell、Nagel\ix{persons}{nagel@Nagel, Robert F.}、Paulsen 與 Tushnet\ix{persons}{tushnet@Tushnet, Mark} 1992 年的 The Constitution Outside the Courts），更多來自行政與立法部門的官員，後者看不出自己對憲法的解釋為何應因司法部門的先行主張而被推翻，而憲法文本並未賦予司法部門解釋上的優先地位（註 8 引 Meese\ix{persons}{meese@Meese, Edwin, III} 1987 年的 The Law of the Constitution，並指出 Lincoln\ix{persons}{lincoln@Lincoln, Abraham（林肯）} 對 Dred Scott\ix{cases}{02@美國!01@聯邦最高法院!1857Dred Scott v. Sandford（1857）@Dred Scott v. Sandford（1857）} 案的態度基於實質相同的立場）。批評者堅持，司法至上\ix{terms}{01@\zhlidxcat{司法審查權限之歸屬}!司法至上}的主張缺乏歷史、憲法體制結構與規範性的政治理論或制度設計論證的支持\pg{456 至 457 頁，註 9 引 Tribe\ix{persons}{tribe@Tribe, Laurence H.} 對 Cooper 的狹義解讀}。

註 6 另指出 City of Boerne 與本文的主要問題「orthogonal」〔正交〕：該案討論國會依增修條文第 14 條\ix{statutes}{02@美國!01@聯邦憲法!114@增修條文第 14 條}第 5 項的權力在何種情形下得偏離最高法院先例，屬實體解釋問題，兩人說這與服從問題有所不同。兩人同時表示，Kennedy\ix{persons}{kennedy@Kennedy, Anthony} 大法官所稱國會「has not just the right but the duty to make its own informed judgment on the meaning and force of the Constitution」〔不僅有權，而且有義務對憲法的意義與效力作出自己知情的判斷〕，是該意見中令人困擾的部分，他們認為那只是對國會的客氣，相當於稱讚敗訴律師的論證。

\subsection{（三）定紛止爭論的要點\pg{457 至 458 頁，註 10--11}}

最高法院在 Cooper\ix{cases}{02@美國!01@聯邦最高法院!1958Cooper v. Aaron（1958）@Cooper v. Aaron（1958）} 與 City of Boerne\ix{cases}{02@美國!01@聯邦最高法院!1997City of Boerne v. Flores（1997）@City of Boerne v. Flores（1997）} 宣告自身的至上地位，在兩人看來與問題無關，因為這種宣告預設了其他部門與公眾是否應接受該宣告的結論\pg{457 頁，註 10 引 Lawson \& Moore}。不依賴這種循環論證，兩人的前文主張可歸納為以下環節：

\begin{enumerate}
\def\labelenumi{\arabic{enumi}.}
\tightlist
\item
  法律的核心道德功能之一，是就應為之事作出定紛止爭\ix{terms}{02@\zhlidxcat{定紛止爭論}!定紛止爭}。
\item
  憲法因修改困難，比制定法與普通法裁判試圖建立更持久的定紛。
\item
  定紛的意義本身有爭議時，司法解釋的核心道德功能就在於確定該（企圖中的）定紛之意義。
\item
  司法解釋可能有誤，並不使其權威受到質疑，因為所有法律上的定紛都主張權威，即使受其拘束者認為它在道德上與法律上有誤。註 11 說明，主張權威即主張提供 content-independent\ix{terms}{02@\zhlidxcat{定紛止爭論}!與內容無關之理由} reasons〔與內容無關之理由〕，其成立與否取決於長期而言受拘束者以定紛為準，與逐案依自己的實體判斷行事相比，所犯錯誤的多寡\pg{註 11 引 Hart\ix{persons}{hart@Hart, H. L. A.（哈特）}、Raz 與 Schauer\ix{persons}{schauer@Schauer, Frederick} \en{Playing By the Rules} 128--34 頁}。
\item
  對第一階\ix{terms}{02@\zhlidxcat{定紛止爭論}!一階定紛與二階定紛}定紛的意義所作的錯誤第二階定紛，其權威的問題比第一階定紛本身的權威問題較小。因此期待人們即使在道德上不同意仍遵守憲法，便可基於同樣理由，期待他們在認為最高法院的解釋有誤時仍遵守之\pg{457 至 458 頁}。
\end{enumerate}

\subsection{（四）本文的目的\pg{458 頁，註 12--13}}

兩人說前文過去三年成為許多批評者的標的。註 12 列六篇專門回應：Devins\ix{persons}{devins@Devins, Neal} \& Fisher\ix{persons}{fisher@Fisher, Louis}、Hartnett\ix{persons}{hartnett@Hartnett, Edward A.}、Nagel\ix{persons}{nagel@Nagel, Robert F.}、Peabody\ix{persons}{peabody@Peabody, Bruce G.}、Sherwin\ix{persons}{sherwin@Sherwin, Emily} 與 Tushnet\ix{persons}{tushnet@Tushnet, Mark} 的 Two Versions of Judicial Supremacy，另附 Harrison\ix{persons}{harrison@Harrison, John} 的未刊稿；註 13 列 Tushnet 的 \en{Taking the Constitution Away from the Courts} 27--29 頁與 Buss、Goldstein、Sherry\ix{persons}{sherry@Sherry, Suzanna} 的文章。批評者視兩人為不顧歷史、為最高法院的傲慢辯護、貶抑國會與行政部門解釋憲法之能力的人。

本文說明目的不在於「critics and \textquotesingle misunderstood\textquotesingle{} authors」〔批評者與「被誤解」的作者〕之間那種針鋒相對的書信往來，兩人以這批批評作為「joining issue」〔就爭點交鋒〕的機會，延續一個涉及憲法的角色與法院在民主社會角色的辯論\pg{458 頁}。

\section{第 I 部分：問題的性質\pg{459 至 467 頁}}
\label{sec:AS-2}

\subsection{（一）先於憲法的憲法\pg{459 至 460 頁，註 14--15}}

第 I 部分 A 節處理兩人所問問題的種類：最高法院應否是憲法的最高解釋者，其解釋應否拘束他人，即使他人自己的解釋看起來更好\pg{459 頁}。他們說這個問題無從由憲法本身回答。Harrison\ix{persons}{harrison@Harrison, John}、Hartnett\ix{persons}{hartnett@Hartnett, Edward A.}、Devins\ix{persons}{devins@Devins, Neal} 與 Fisher\ix{persons}{fisher@Fisher, Louis} 都強調，兩人忽略憲法未指定最高法院為最高解釋者這項具有決定性的資料，也忽略說明並支持此點的憲法史\pg{459 頁，註 14 載三人著作的頁次}。

Alexander\ix{persons}{alexander@Alexander, Larry} 與 Schauer\ix{persons}{schauer@Schauer, Frederick} 為論證之便接受這些歷史證據：制憲者無意使最高法院成為憲法的最高解釋者，各州批准者亦無此理解，文本並未指定。他們再進一步假設憲法實際上否定這樣的角色，但仍主張其論證涉及「this Constitution in this constitutional system」〔這部憲法、這個憲法體制〕，並未停留在假想的體制\pg{459 頁}。理由是他們論證的是一項 preconstitutional\ix{terms}{03@\zhlidxcat{憲法之效力基礎}!先於憲法的憲法} norm〔先於憲法的規範〕，這項規範「determines not what was adopted then, but how what was adopted then should be regarded now」〔決定當時所採行者在今日應被如何看待〕。憲法作為根本法的權威「ultimately rests not on facts about the past, but on the Constitution\textquotesingle s acceptance as authoritative in the present」〔終究取決於它在現在被接受為有權威，與過去的事實無關〕。他們稱此為邏輯上的論點，並說整個論證以它為基礎\pg{459 至 460 頁，註 15 引 Michelman\ix{persons}{michelman@Michelman, Frank I.} 論人民的制憲者身分}。

\subsection{（二）假設的最高法院解釋條款\pg{460 至 461 頁，註 16--18}}

兩人設想一個與事實相反的情形來說明上述論點：憲法第 6 條\ix{statutes}{02@美國!01@聯邦憲法!006@第 6 條}的至上條款不僅寫「This Constitution ... shall be the Supreme Law of the Land」，還明定最高法院對憲法的解釋對其他聯邦與各州法官、行政與立法官員及各州官員「shall be final and binding」〔具有終局拘束力〕\pg{460 頁}。即使憲法文本真是如此，司法至上\ix{terms}{01@\zhlidxcat{司法審查權限之歸屬}!司法至上}的問題仍然未決，因為仍須問誰來解釋該條款，而該解釋是否拘束他人。

這可以用 Wittgenstein\ix{persons}{wittgenstein@Wittgenstein, Ludwig（維根斯坦）} 的說法表達，即規則無法決定其自身的適用（註 16 以箭頭記號為例，問人如何知道該朝箭頭尖端或尾端前進）；也可以用 Kelsen\ix{persons}{kelsen@Kelsen, Hans（凱爾森）} 的 grundnorm\ix{terms}{03@\zhlidxcat{憲法之效力基礎}!基本規範}〔基本規範〕或 Hart\ix{persons}{hart@Hart, H. L. A.（哈特）} 的 ultimate rule of recognition\ix{terms}{03@\zhlidxcat{憲法之效力基礎}!承認規則}〔終極承認規則〕表達\pg{460 至 461 頁，註 17 引 Kelsen「is grounded as valid by the presupposed basic norm」；註 18 引 Hart}。共同的要點是：法律文本的地位無法，也不能在不陷入無窮後退的情形下，由該文本自身決定。判斷文本的地位須離開文本，因此文本也無法回答關於某人對該文本的解釋，包括對其中指定權威解釋者的部分的解釋，是否有權威\pg{461 頁}。

\subsection{（三）「two」與 Conglish\pg{461 至 462 頁，註 19--23}}

兩人承認就這個虛構條款而言，上述論點在許多方面屬形式性的。同一論點也適用於第 3 條第 3 項\ix{statutes}{02@美國!01@聯邦憲法!003@第 3 條}「testimony of two Witnesses」中的 two 是否指 3：在嚴格邏輯上，文本本身無法決定它應被理解為普通英文（two 指 2），或某種 Conglish（two 指 3；註 20 稱此例為 Kripke 對 Wittgenstein\ix{persons}{wittgenstein@Wittgenstein, Ludwig（維根斯坦）} 的討論之變形）\pg{461 頁}。他們接著承認，以普通英文理解憲法是根深蒂固而被廣泛接受的先於憲法\ix{terms}{03@\zhlidxcat{憲法之效力基礎}!先於憲法的憲法}的理解，將 two 解為 3 會違反最穩固的先於憲法的規則之一。

兩人隨即指出，以普通英文解讀文本並不如設想的那樣清楚。實際憲法文本中已有明顯違反文字的解釋：第 11 修正案的「another State」被解釋為包含同一州（註 21 引 Monaco v. Mississippi\ix{cases}{02@美國!01@聯邦最高法院!1934Monaco v. Mississippi（1934）@Monaco v. Mississippi（1934）} 與 Hans v. Louisiana\ix{cases}{02@美國!01@聯邦最高法院!1890Hans v. Louisiana（1890）@Hans v. Louisiana（1890）}），以及依 Stevens 大法官的看法，禁止各州「impairing the Obligation of Contracts」〔損害契約義務〕的規定亦然（註 22 引 Keystone Bituminous Coal Ass\textquotesingle n v. DeBenedictis 的「not to be read literally」）。這些被普遍接受的對文字的違反，被理解為與其餘文本繼續有效並行不悖（註 23 說明無人把它們理解為廢止或取代整部憲法），所以對虛構條款所作的解釋，使其支持的司法解釋至上\ix{terms}{01@\zhlidxcat{司法審查權限之歸屬}!解釋至上}程度低於其字面所示，在邏輯上並非不可能\pg{461 至 462 頁}。由此兩人說：「What counts as law is determined extra-legally, and what counts as the constitution is determined extra-constitutionally」〔何者算是法律，在法律之外決定；何者算是憲法，在憲法之外決定〕，所以偏離文本表面意義者是否有效，無法以循環地援引文本回答\pg{462 頁}。

\subsection{（四）實際憲法的情形與兩種 constitution\pg{462 至 463 頁，註 24--26}}

兩人以 a fortiori〔更何況〕推及實際文本：連如此明確的司法至上\ix{terms}{01@\zhlidxcat{司法審查權限之歸屬}!司法至上}條款都使問題在邏輯上未決，則在一部既未處理司法審查\ix{terms}{01@\zhlidxcat{司法審查權限之歸屬}!司法審查}權本身（註 24 引 Harrison\ix{persons}{harrison@Harrison, John} 與 Thayer\ix{persons}{thayer@Thayer, James Bradley} 1893 年的論文），也未處理司法解釋權威的憲法上，這個問題「remains stunningly open」〔仍然令人驚訝地未決〕\pg{462 頁}。

他們用「preconstitutional\ix{terms}{03@\zhlidxcat{憲法之效力基礎}!先於憲法的憲法}」指邏輯上的先於，與時間上的先後無關\pg{462 至 463 頁，註 25 引 Kay}。憲法文本應如何解讀，在邏輯上先於文本，因為這些問題最終決定文本的意義；因為它們同時構成一個國家與一部憲法，所以先於憲法的問題在某個意義上具有構成性，因而也是憲法性的。要理解大寫 C 的 Constitution，須先理解它所依託的小寫 c 的 constitution\pg{463 頁，註 26 引 Michelman\ix{persons}{michelman@Michelman, Frank I.} 與 Schauer\ix{persons}{schauer@Schauer, Frederick}}。兩人認為誰的解釋應被視為有權威，是這一類先於憲法的問題中最重要的一個。

據此，他們認為 Devins\ix{persons}{devins@Devins, Neal}、Fisher\ix{persons}{fisher@Fisher, Louis}、Hartnett\ix{persons}{hartnett@Hartnett, Edward A.} 與 Harrison 的歷史與文本分析雖具啟發，結論性遠低於這些作者所認為的。兩人以箭頭作類比：若有人問為何朝箭頭尖端開車，回答說箭頭告訴我們朝尖端走，就是沒有理解使箭頭有意義的、箭頭之外的條件。同理，主張應終局地以憲法文本與歷史決定該文本與歷史告訴我們什麼，犯的是同一個概念錯誤\pg{463 頁}。

\subsection{（五）對「既非經驗亦非歷史」一語的修正\pg{464 頁，註 27}}

第 I 部分 B 節承認批評者在先於憲法\ix{terms}{03@\zhlidxcat{憲法之效力基礎}!先於憲法的憲法}的性質上未必同意，更實質的分歧在於用何種分析或證據回答關於解釋權威的問題\pg{464 頁}。排除求諸憲法自身的自我指涉之後，指引在何處？

前文說這項探討「neither empirical nor historical」〔既非經驗的、也非歷史的〕\pg{註 27 引前文 1369 頁}。兩人現在說這樣表述「Perhaps too casually」〔也許過於隨意〕，作為對所作所為的描述（著重制度設計的邏輯與哲學問題）雖然準確，但令許多批評者疑惑：這項探討怎能那樣不涉經驗？兩人的結論是：「the empirical dimension is one that cannot be avoided」〔經驗面向是無法避免的〕，歷史可以是這種經驗探討的構成部分之一，但只是之一。他們說這不改變其主張的性質與所用證據，但對當時的措辭表示後悔，也後悔暗示制度設計的問題可以在經驗真空中回答\pg{464 頁}。

\subsection{（六）歷史與文本的同等地位\pg{464 至 465 頁，註 28--29}}

說探討不涉歷史，指的是：以歷史的權威回答先於憲法\ix{terms}{03@\zhlidxcat{憲法之效力基礎}!先於憲法的憲法}的問題，與依賴憲法文本具有相同地位，終究是 question-begging〔乞題〕\pg{464 頁}。兩人假設發現五十五位制憲會議代表與批准各州的多數議會，都把第 3 條與\ix{statutes}{02@美國!01@聯邦憲法!003@第 3 條}第 6 條理解\ix{statutes}{02@美國!01@聯邦憲法!006@第 6 條}為指涉最高法院的解釋專屬與解釋至上\ix{terms}{01@\zhlidxcat{司法審查權限之歸屬}!解釋至上}，那時可說 Madison\ix{persons}{madison@Madison, James（麥迪遜）} 等人同意了他們的主張，但兩人仍明示不依賴 Madison 的想法決定該問題。

理由仍在自我指涉。註 28 指出這是 Schauer\ix{persons}{schauer@Schauer, Frederick} 先前用過的例子之一的變體：假設兩人寫一部以他們為最高解釋者的憲法，規定簽署並在 Constitutional Commentary 上刊登簽署的通知後生效（註 29 寫著「We signed it.」〔我們簽了〕）。這部憲法自己規定的生效條件已經成就，並依其內部條件有效，但這遠不足以使它成為美國憲法。憲法之為憲法，取決於一個政治共同體在經驗上接受其為憲法；沒有這種「acceptance as social fact」〔作為社會事實\ix{terms}{03@\zhlidxcat{憲法之效力基礎}!社會事實}的接受〕，再多的形式內部效力也不使文件成為憲法\pg{465 頁}。

據此，決定 Madison 的憲法成為美國憲法、Alexander\ix{persons}{alexander@Alexander, Larry} 與 Schauer 的憲法則否的因素，並未含在文件之中，1787 年也無從決定，現在由美國人民決定\pg{465 頁}。1787 年憲法之成為 2000 年憲法，取決於 2000 年的社會事實，則制憲者在 1787 年的意圖於 2000 年具有何種地位，也取決於 2000 年的決定。「The relevance of history is not determined historically, but by a present political and social decision」〔歷史有何相關性，取決於當下的政治與社會決定〕，所以 1787 年先於憲法的問題的判斷，必然是非歷史的問題；這個非歷史的問題完全可能以賦予歷史一定角色的方式回答\pg{465 至 466 頁}。

\subsection{（七）Cooper 規則的來源與歷史的內容價值\pg{466 頁}}

Madison\ix{persons}{madison@Madison, James（麥迪遜）} 對司法至上\ix{terms}{01@\zhlidxcat{司法審查權限之歸屬}!司法至上}的想法，若他真的想過，與文本所說的一樣不具決定性\pg{466 頁}。兩人的論證是：Cooper 規則\ix{terms}{01@\zhlidxcat{司法審查權限之歸屬}!Cooper 規則}在規範上優於他們稱為（兩人承認帶有傾向性）「institutional anarchy」〔制度上的無政府狀態\ix{terms}{02@\zhlidxcat{定紛止爭論}!制度上的無政府狀態}〕的狀態。若原始憲法沒有 Cooper\ix{cases}{02@美國!01@聯邦最高法院!1958Cooper v. Aaron（1958）@Cooper v. Aaron（1958）} 規則，則原始憲法在規範上不如含有該規則的憲法。取得該規則可以經第 5 條的修憲程序\ix{statutes}{02@美國!01@聯邦憲法!005@第 5 條}，兩人主張的途徑則是以取得原始憲法本身的方式取得它，也就是接受它為有權威。他們說：「In Cooper, the Supreme Court ran that rule up the flagpole. If we like it, all we have to do is salute.」〔Cooper 案中，最高法院把這項規則升上旗桿。如果我們喜歡，只須敬禮。〕並說多數人似乎確已敬禮，若然，Cooper 規則就與憲法本身同樣「constitutional」\pg{466 頁}。

退一步問社會是否應該敬禮時，兩人承認先前對歷史的價值否定得過快。歷史不具權威，不等於與權威無關而不具意義。制憲者聰明，與 Plato\ix{persons}{plato@Plato（柏拉圖）}、Aristotle\ix{persons}{aristotle@Aristotle（亞里斯多德）}、Confucius、Hobbes、Locke、Lincoln\ix{persons}{lincoln@Lincoln, Abraham（林肯）}、Stanton、Rawls\ix{persons}{rawls@Rawls, John（羅爾斯）}、Douglass 一樣可以學習；法國、日本、印度、巴西、南非的歷史亦然。兩人說：「what we learn from history we learn because of its content and not because of its authority」〔我們從歷史學到的東西，因其內容而學到，與其權威無關〕。Madison 說過什麼應與探討有關，但離決定問題還差得遠\pg{466 頁}。

\subsection{（八）規範性的制度設計與經驗證據\pg{466 至 467 頁}}

先於憲法\ix{terms}{03@\zhlidxcat{憲法之效力基礎}!先於憲法的憲法}的問題是規範性的制度設計問題，亦即應採行何種制度設計，不能與世界上的事實無關地探討。兩人因此說，主張探討不涉經驗「was also then and is now a mistake」〔當時與現在都是錯誤〕。正確的主張是：問題的地位，即其先於憲法的性質，本身是邏輯問題，與歷史或經驗無關；問題的地位確立之後，要判斷某種制度設計是否較佳時，無視來自外部世界的證據，包括歷史證據，是愚蠢的\pg{466 至 467 頁}。他們認為，推論某些誘因與制度設計會導致某些行為屬於適當，但這些是經驗推論，應以實際行為的證據補充與檢驗，證據可以是歷史的、社會學的、統計的、心理學的或其他性質\pg{467 頁}。

\section{第 II 部分：定紛止爭功能及其最適所在\pg{467 至 478 頁}}
\label{sec:AS-3}

\subsection{（一）定紛止爭是否被高估\pg{467 至 468 頁，註 30}}

第 II 部分 A 節重述核心主張：法律，尤其憲法，主要功能之一是為了定紛止爭\ix{terms}{02@\zhlidxcat{定紛止爭論}!定紛止爭}而就應為之事提供某種程度的協調性定紛。在道德與政治歧見並存的世界，法律可以提供一種「neither final nor conclusive, but nevertheless authoritative」〔既非終局、也非結論性，但仍有權威〕的定紛，對一階\ix{terms}{02@\zhlidxcat{定紛止爭論}!一階定紛與二階定紛}歧見提供二階的解決，使決定得以作成、行動得以協調\pg{467 頁}。

批評者分兩種。第一種以 Tushnet\ix{persons}{tushnet@Tushnet, Mark} 為最，其他人大多也屬之，在他們背後還有龐大的傳統，在公共的憲法論述中看到獨立的價值，認為這種論述不經持續的公共爭議與辯論就不會出現\pg{467 至 468 頁}。註 30 說明這個傳統涉及的文獻可舉出數以千計，並列 Ackerman \en{We the People}、Levinson\ix{persons}{levinson@Levinson, Sanford} \en{Constitutional Faith}、Sunstein \en{The Partial Constitution}、Lipkin、Michelman\ix{persons}{michelman@Michelman, Frank I.} 與 Minow。依此類主張，憲法的首要價值在於它是公共審議的場域與焦點，使犯罪控制、平權措施、性別平等、性傾向、權力下放給州、總統行為、宗教在公共生活中的角色等議題得到辯論；過度強調憲法的定紛功能，會截斷憲法傳統核心的公共辯論與審議\pg{468 頁}。

\subsection{（二）公共討論與成文憲法的關係\pg{468 頁，註 31--33}}

兩人的回應是「We demur」〔我們不同意〕。他們不否認公共論述的價值（註 31 補「At least here」，並指 Schauer\ix{persons}{schauer@Schauer, Frederick} 另有 Discourse and Its Discontents 討論其局限），但認為以為美國人需要憲法才能辯論平權措施、刑事司法、墮胎、宗教與國家、隱私、死刑，是美國憲法學者的 conceit〔自負〕。這些辯論在沒有單一成文憲法文件的國家，主要是英國、紐西蘭與以色列，同樣興盛（註 32 引 Finer 等與 Jackson \& Tushnet\ix{persons}{tushnet@Tushnet, Mark} 的憲法比較著作）；在成文憲法並未將這些特定議題憲法化的國家亦然（註 33 引 Bell 論法國憲法）。他們承認要權威地回答這個問題，須有以憲政各面向為變項的比較審議研究，而這種研究尚不存在，但依其看法，憲法是否為對重大政治與道德問題的豐富公共辯論的必要條件、充分條件，甚至顯著的因果貢獻者，「far from certain」〔遠不確定〕\pg{468 頁}。

\subsection{（三）法律的特殊功能\pg{468 至 469 頁，註 34}}

兩人說明其論證的前提：同意優於歧見、定紛\ix{terms}{02@\zhlidxcat{定紛止爭論}!定紛止爭}優於保持開放，雖然兩種世界觀都有道理，卻都不在其內；其前提是法律所履行的特殊功能，即同意與定紛似乎是法律的特有領域，也是被以近似法律的方式理解的成文憲法的特有領域\pg{469 頁}。另有以討論為焦點的公共文件，例如孔子或毛主席的著作，兩人的論證以美國憲法與這類文件之間，以及憲法與 Declaration of Independence 之間有重要差異為前提\pg{469 頁}。

註 34 指出，Tushnet\ix{persons}{tushnet@Tushnet, Mark} 在 \en{Taking the Constitution Away from the Courts} 中特別以憲法序言與獨立宣言這類非法律文件作為公共討論基本價值的焦點。兩人的問題是把重大事項的公共辯論與憲政主義兩個問題混為一談是否明智：若兩者相同，Tushnet 對而兩人錯；若不同，差別必在公共論述之外，即對該論述或其產物的限制之中。

\subsection{（四）憲法文本中已定紛的事項與 Brandeis 之語\pg{469 至 470 頁，註 35--37}}

憲法處理的事項通常比靠左或靠右行車更具道德分量，但為定紛\ix{terms}{02@\zhlidxcat{定紛止爭論}!定紛止爭}而定紛的價值並不因此在憲法中缺席。兩人列舉人們可能有嚴肅道德與政治歧見的事項：比例代表制或「first past the post」〔單一選區相對多數決〕、有無國教、民事或刑事案件是否採陪審團、刑事被告能否被迫作證（違者以藐視法庭處罰或可推論有罪）、總統可連任幾個四年任期。這些在美國憲法文本中大致已有定紛，從議程上移除，連不同意其實質內容的人也一樣\pg{469 頁}。法律被當作有權威時，為不同意定紛內容的人提供服從的 content-independent\ix{terms}{02@\zhlidxcat{定紛止爭論}!與內容無關之理由} reason\pg{469 至 470 頁}。

兩人引 Brandeis\ix{persons}{brandeis@Brandeis, Louis D.（布蘭戴斯）} 大法官在 Burnet v. Coronado Oil\ix{cases}{02@美國!01@聯邦最高法院!1932Burnet v. Coronado Oil \& Gas Co.（1932）@Burnet v. Coronado Oil \& Gas Co.（1932）} \& Gas Co.（1932）不同意見\ix{terms}{04@\zhlidxcat{法院之功能}!不同意見與協同意見}書的話：「in most matters it is more important that the applicable rule of law be settled than that it be settled right」〔多數事項上，適用的法律規則被確定，比被確定得正確更為重要〕\pg{470 頁，註 36}。他們指出，有人主張強力的先例拘束在憲法脈絡中重要性較低（註 37 引 Webster 與 Payne 兩案及 Cooper\ix{cases}{02@美國!01@聯邦最高法院!1958Cooper v. Aaron（1958）@Cooper v. Aaron（1958）}），但憲法文本本身強力地支持 Brandeis 的觀點，制憲者認為為定紛而定紛的約束，往往比面對世界變遷時更大的彈性更有價值\pg{470 頁}。

註 35 回應 Hartnett\ix{persons}{hartnett@Hartnett, Edward A.}（並言及 Nagel\ix{persons}{nagel@Nagel, Robert F.}）認為兩人反對歧見、勸人服從即是譴責歧見的主張，說不知如何理解；承認前文用語「without regard」有歧義，可能促成誤解，但以每小時 55 英里的速度行駛並公開批評任何低於 70 英里的限速，兩者可以並行，同樣，對最高法院解釋的公開異議，與服從該解釋直到法院改變看法，完全相容。

\subsection{（五）文本未定紛的爭議與最高法院意見\pg{470 至 472 頁，註 38--47}}

上述例子是 biased sample〔有偏的樣本〕，因為都是文本已相當明確的情形。美國憲法文本如此不確定，更常見的是帶有憲法面向而文本本身未提供定紛\ix{terms}{02@\zhlidxcat{定紛止爭論}!定紛止爭}的爭議\pg{470 頁}。兩人列了一串問題：

\begin{enumerate}
\def\labelenumi{\arabic{enumi}.}
\tightlist
\item
  增修條文第 1 條\ix{statutes}{02@美國!01@聯邦憲法!101@增修條文第 1 條}的禁止設立國教條款是否禁止公立學校中由教師組織的、無教派色彩的禱告（註 38 引 Wallace v. Jaffree\ix{cases}{02@美國!01@聯邦最高法院!1985Wallace v. Jaffree（1985）@Wallace v. Jaffree（1985）}、Abington Sch. Dist. v. Schempp\ix{cases}{02@美國!01@聯邦最高法院!1963Abington School District v. Schempp（1963）@Abington School District v. Schempp（1963）}、Engel v. Vitale\ix{cases}{02@美國!01@聯邦最高法院!1962Engel v. Vitale（1962）@Engel v. Vitale（1962）}），或畢業典禮上由學生主導者（註 39 引 Lee v. Weisman\ix{cases}{02@美國!01@聯邦最高法院!1992Lee v. Weisman（1992）@Lee v. Weisman（1992）}）。
\item
  平等保護條款是否禁止各州維持單一性別的大學（註 40 引 Mississippi Univ. for Women v. Hogan\ix{cases}{02@美國!01@聯邦最高法院!1982Mississippi University for Women v. Hogan（1982）@Mississippi University for Women v. Hogan（1982）} 與 United States v. Virginia\ix{cases}{02@美國!01@聯邦最高法院!1996United States v. Virginia（1996）@United States v. Virginia（1996）}）。
\item
  第 4 修正案的令狀要求是否禁止在審判中使用未經令狀取得的、具證明力的證據（註 41 引 Mapp v. Ohio\ix{cases}{02@美國!01@聯邦最高法院!1961Mapp v. Ohio（1961）@Mapp v. Ohio（1961）}）。
\item
  言論自由條款是否涵蓋明示鼓勵種族暴力的言論（註 42 引 Brandenburg v. Ohio\ix{cases}{02@美國!01@聯邦最高法院!1969Brandenburg v. Ohio（1969）@Brandenburg v. Ohio（1969）}）。
\item
  死刑是否殘酷而不尋常（註 43 引 Gregg v. Georgia\ix{cases}{02@美國!01@聯邦最高法院!1976Gregg v. Georgia（1976）@Gregg v. Georgia（1976）}）。
\item
  國會未明示行動時，各州能否對州際商業課以財務或後勤上的負擔（註 44 引 Kassel）。
\item
  現代立法實務的需要是否容許國會將立法職責委由行政機關行使（註 45 引 Mistretta 與 American Textile），或省略兩院通過與送交總統的正式程序（註 46 引 INS v. Chadha\ix{cases}{02@美國!01@聯邦最高法院!1983INS v. Chadha（1983）@INS v. Chadha（1983）}）。
\end{enumerate}

這些場合都有合理的道德與政治歧見而無清楚的文本定紛，取而代之的是最高法院試圖解決該議題的意見，其解決當然不消除根本的歧見\pg{471 頁}。兩人的核心主張是：在這種情形，為定紛而定紛的價值，使最高法院以外的機構，尤其是下級法院、州議會、國會與總統，應把該解決視為有權威，即使它們繼續不同意其實質內容。承認最高法院的權威與解釋至上\ix{terms}{01@\zhlidxcat{司法審查權限之歸屬}!解釋至上}，這些機構就如同承認憲法文本本身的權威與至上，對穩定與社會和諧作出貢獻\pg{471 頁}。

兩人以 Clinton 總統為例：若他承認服從第 22 修正案、不競選第三任期具有定紛的價值，則同理他也應承認拒絕簽署一部明顯違反最高法院近期一致見解的 Communications Decency Act 具有定紛的價值\pg{471 頁，註 47 對照 Reno v. ACLU\ix{cases}{02@美國!01@聯邦最高法院!1997Reno v. ACLU（1997）@Reno v. ACLU（1997）} 與 Sable Communications v. FCC\ix{cases}{02@美國!01@聯邦最高法院!1989Sable Communications v. FCC（1989）@Sable Communications v. FCC（1989）}}。他們進一步說：若定紛論證失敗，則憲法的大部分難以解釋；若成功，則它支持單一權威解釋者，其強度與它支持單一權威憲法相同\pg{471 至 472 頁}。

\subsection{（六）Tushnet、Peabody 與 Sherwin 的批評\pg{472 頁，註 48--50}}

兩人說 Tushnet\ix{persons}{tushnet@Tushnet, Mark} 對法律式定紛\ix{terms}{02@\zhlidxcat{定紛止爭論}!定紛止爭}作為超越性價值採取最強烈的反對立場\pg{472 頁，註 48 引 Tushnet, Two Versions of Judicial Supremacy 952--57 頁}；Peabody\ix{persons}{peabody@Peabody, Bruce G.} 同樣批評兩人過度強調定紛（註 49），雖承認其部分保留，仍主張兩人對定紛與穩定的推崇貶低了其他重要價值，包括實體與程序的，如創造性及迅速適應變動世界所需的彈性\pg{472 頁，註 50}。

註 50 另處理 Peabody 與 Sherwin\ix{persons}{sherwin@Sherwin, Emily} 的批評。Sherwin 對 Cooper\ix{cases}{02@美國!01@聯邦最高法院!1958Cooper v. Aaron（1958）@Cooper v. Aaron（1958）} 的辯護沒有異議，她的批評在於兩人如何調和該辯護與他們對 Lincoln\ix{persons}{lincoln@Lincoln, Abraham（林肯）} 有義務服從 Dred Scott\ix{cases}{02@美國!01@聯邦最高法院!1857Dred Scott v. Sandford（1857）@Dred Scott v. Sandford（1857）} 案命令所說的話。兩人主張 Lincoln 有理由，或至少可以有理由不服從法院，這項讓步不損及對 Cooper 的辯護；Sherwin 則認為那樣兩人所主張的實際上只是服從的一項考慮因素，算不上真正的規則，而唯有真正的規則才能提供穩定、可預測與協調的定紛利益。兩人回應，同意強力規則的必要性，但認為她低估了權威的不對稱\ix{terms}{02@\zhlidxcat{定紛止爭論}!權威的不對稱}：權威的規則主張絕對的服從，不是僅供考慮，要求受拘束者即使確信衡量各項理由後，包括有規則及其定紛利益的理由，不服從為宜，仍須服從；受拘束者面對權威規則時，卻不能讓渡自己的理性與道德主體性。因此權威所應命令的與主體所應為的之間，永遠有差距，說法院從其立場有理由要求 Lincoln 服從，而 Lincoln 從其立場在道德上有理由不服從，並無不一致\pg{472 頁，註 50}。

\subsection{（七）法律的職責與憲政主義的價值\pg{472 至 473 頁}}

兩人說自己的論證隱含一點，現在要明示：法律、法院或憲法不是，也不應是制度設計中一切善的寄託。他們以汽車煞車為喻：設計煞車時不應以速度、加速度或操控的容易為目標，即使承認這些目標也相關。同理，兩人依賴的前提是法律性質的制度比其他社會制度更能發揮促進穩定與定紛\ix{terms}{02@\zhlidxcat{定紛止爭論}!定紛止爭}的功能，在複雜而分化的社會中這不一定是錯誤的職能分工觀念\pg{472 至 473 頁}。主張定紛是法律的特殊職責，並不是主張所有社會制度有同樣職責，或定紛在社會價值中的地位，應與它在法律這個社會子集中同樣重要\pg{473 頁}。

兩人據此說，從定紛來看司法至上\ix{terms}{01@\zhlidxcat{司法審查權限之歸屬}!司法至上}，就是把憲法視為提供穩定與約束的功能。若穩定與約束在憲法價值中不居首位，他們不明白為何一開始要對一階\ix{terms}{02@\zhlidxcat{定紛止爭論}!一階定紛與二階定紛}政治決定施加二階的憲法約束。若穩定與約束是說明憲政主義本身的核心，則主張單一權威解釋者的論證「flows directly from the deepest values of constitutionalism」〔直接源於憲政主義最深層的價值〕\pg{473 頁}。

\subsection{（八）法院是否為最佳的定紛者\pg{473 至 474 頁，註 51}}

第 II 部分 B 節處理的批評，把除 Sherwin\ix{persons}{sherwin@Sherwin, Emily} 之外所有批評者連在一起：兩人從定紛\ix{terms}{02@\zhlidxcat{定紛止爭論}!定紛止爭}本身是好事的前提，過快地推到法院，尤其是最高法院，是履行此功能的最佳機構。這些批評者承認定紛是好事，甚至承認它可能需要單一權威解釋者，但質疑為何應交給最高法院，而非其他機構\pg{473 頁}。

論證之一強調最高法院與國會、與行政部門之間的對話。Devins\ix{persons}{devins@Devins, Neal} 與 Fisher\ix{persons}{fisher@Fisher, Louis} 為兩人稱為「institutional anarchy」的狀態的暫時形式辯護，主張法院與政治部門之間比 Cooper\ix{cases}{02@美國!01@聯邦最高法院!1958Cooper v. Aaron（1958）@Cooper v. Aaron（1958）} 所容許的複雜得多的往返是好事，長期會產生較好的憲法決定，因而帶來政治穩定的提高\pg{473 頁}。兩人說自己沒有做過回應這項命題所需的完整社會科學檢驗，因此無法直接反駁。若 Devins 與 Fisher 的意思是所有權威都脆弱地建立在受拘束者的接受之上，權威所提供的 content-independent\ix{terms}{02@\zhlidxcat{定紛止爭論}!與內容無關之理由} reasons 因深刻的弔詭而同樣脆弱，兩人同意\pg{473 至 474 頁；註 51 說明 Alexander\ix{persons}{alexander@Alexander, Larry} 在多篇文章中多次強調此點，列 Can Law Survive the Asymmetry of Authority? 等}。

以下（九）至（十三）在刊本 474 至 478 頁，與下一節丁\pg{478 至 481 頁}同屬課綱指定頁次\pg{455 至 473 頁、481 至 482 頁}以外的部分。

\subsection{（九）公眾對權威定紛的接受\pg{474 頁，註 52--54}}

若 Devins\ix{persons}{devins@Devins, Neal} 與 Fisher\ix{persons}{fisher@Fisher, Louis} 的意思是不應權威地定紛\ix{terms}{02@\zhlidxcat{定紛止爭論}!定紛止爭}憲法問題，因為這樣做冒著損害權威所賴以存在的接受之風險，兩人則表示反對。公眾可以也常常接受權威的定紛，包括憲法上的定紛，即使在道德、政策與文本解釋上強烈不同意\pg{474 頁}。兩人舉例：早期 New Deal 階段最高法院宣告為對抗大蕭條而制定的聯邦立法無效的裁判（註 52 列 United States v. Butler\ix{cases}{02@美國!01@聯邦最高法院!1936United States v. Butler（1936）@United States v. Butler（1936）} 與 A.L.A. Schechter Poultry\ix{cases}{02@美國!01@聯邦最高法院!1935A.L.A. Schechter Poultry Corp. v. United States（1935）@A.L.A. Schechter Poultry Corp. v. United States（1935）} Corp. v. United States），以及 Warren 法院宣告學校禱告違憲的裁判（註 53 指向註 39），在實質上都極不受歡迎，在解釋上也被廣泛批評為錯誤，但都沒有被廣泛認為缺乏權威。

兩人說權威可能只建立在受拘束者逐時逐刻的接受之上，但這種接受可以承受相當程度的、對政策與解釋的深刻歧見。又，若對司法角色的接受與對司法的尊重\ix{terms}{05@\zhlidxcat{解釋理論與自制}!尊重}主要取決於裁判的實質內容（註 54 引 Adamany \& Grossman、Caldeira 與 Gibson 等人的研究），則法院角色結構上的差異，對其被接受的程度與正當性可能影響不大\pg{474 頁}。

\subsection{（十）解釋分歧與實體結果分歧的區分\pg{474 至 475 頁，註 55}}

兩人主張須區分憲法歧見的種類。分歧僅在於最高法院對憲法規則的解釋，不涉實體結果，也就是人們認為最高法院的憲法規則優於制憲者採用的規則時，Cooper 規則\ix{terms}{01@\zhlidxcat{司法審查權限之歸屬}!Cooper 規則}的理由最強\pg{474 至 475 頁，註 55 補充不宜高估人民對最高法院判決所知所信}。這不是說法院不應忠實解釋憲法，即使法院認為憲法在道德或政策上不是最佳甚至有害；僅指出利害只在於對先前定紛\ix{terms}{02@\zhlidxcat{定紛止爭論}!定紛止爭}的解釋是否正確時，定紛的價值似乎重於解釋正確的價值\pg{475 頁}。

政治部門認為最高法院的解釋不僅不正確，而且不公正或政策不良時，定紛是否應讓位於實體正確？兩人從制度設計的立場，也就是貫穿其整個計畫的立場，仍答「no」\pg{475 頁}。理由是許多，或許大多數憲法議題涉及實質重要的事項，所以幾乎所有對法院的解釋性歧見都伴隨政策歧見。沒有權威的定紛，大部分憲法議題就維持開放：任何不同意最高法院憲法裁判的政府行動者，都會覺得在法律與道德上，即使不一定在政治或審慎\ix{terms}{05@\zhlidxcat{解釋理論與自制}!審慎}考量上，可以不理會該裁判。兩人舉例：不同意第 4 修正案令狀法現狀的警察，若認為不聲請令狀也能達成目的，就會略過；下級法院法官若無遵守最高法院裁判的法律義務，他達成目的的機會更高；州議會想要公立學校的有組織禱告、市政府想規制猥褻之外的不雅內容、行政官員想運用權力勸阻婦女墮胎，情形亦然\pg{475 頁}。

\subsection{（十一）批評者未主張的範圍\pg{475 至 476 頁，註 56--58}}

兩人說批評者沒有人主張下級法院無須遵循最高法院（註 56 保留 Hartnett\ix{persons}{hartnett@Hartnett, Edward A.} 為可能的例外，指向第 III 部分；註 57 引 Rogers 與 Fried），也無人主張 Meese\ix{persons}{meese@Meese, Edwin, III} 司法部長十二年前主張的立場，即各州不論在行政、立法、執法或司法的身分，均得把最高法院裁判視為只限於該案、不為其他行動者設定法律\pg{475 至 476 頁，註 58 引 Meese 與 Tulane Law Review 專題}。批評者不願強烈支持把解釋權分散給 50 個州，顯示辯論的範圍比乍看窄，兩人對國家憲法的全國解釋至上\ix{terms}{01@\zhlidxcat{司法審查權限之歸屬}!解釋至上}的論證已為人廣泛接受\pg{476 頁}。

仍然，Devins\ix{persons}{devins@Devins, Neal} 與 Fisher\ix{persons}{fisher@Fisher, Louis}、Nagel\ix{persons}{nagel@Nagel, Robert F.}、Peabody\ix{persons}{peabody@Peabody, Bruce G.} 與 Tushnet\ix{persons}{tushnet@Tushnet, Mark} 都主張，需要權威定紛\ix{terms}{02@\zhlidxcat{定紛止爭論}!定紛止爭}並不決定哪個機構應擁有該權威。民主組成的機構也能定紛憲法議題，所以對某些或許多議題，國會在履行定紛功能上可能優於法院\pg{476 頁}。

\subsection{（十二）最高法院作為定紛者的理由\pg{476 至 478 頁，註 59--61}}

兩人提出的理由之一，是法院「relative insulation from political winds」〔與政治風向的相對隔絕〕，除非認為憲法解釋就是且應是當代價值與政策的表達，否則這是明顯的優點，使最高法院優於國會或行政部門\pg{476 頁}。Nagel\ix{persons}{nagel@Nagel, Robert F.} 特別指責兩人誇大最高法院的穩定程度\pg{註 59 引 Nagel 857--60 頁}，兩人回應，穩定與否只能以兼顧變動頻率與變動幅度的函數判斷。他們舉住宅電流為例：在 108 至 112 伏之間不斷波動的電流，在多數用途上比七天中有六天在 109.9 至 110.1 伏、一天在 62 或 154 伏的電流更穩定。因此只看最高法院不穩定的次數而不比較其與其他部門的變動幅度，如 Nagel 似乎做的，是錯誤的；最高法院的不穩定若頻繁但幅度窄（註 60 引 Blasi 編的 \en{The Burger Court: The Counter-Revolution That Wasn\textquotesingle t}），就更是如此\pg{476 頁}。

兩人的懷疑（承認尚缺他們與批評者都未做過的受控研究）是：最高法院對先例的遵循，使它能比沒有類似先例拘束實踐的政治部門更持久地定紛\ix{terms}{02@\zhlidxcat{定紛止爭論}!定紛止爭}憲法議題\pg{476 至 477 頁，註 61 引 Scalia、Tribe\ix{persons}{tribe@Tribe, Laurence H.} 與 Monaghan}。即使最佳的制度設計是免除法院遵循自己先例的義務（註 62，見下一小節），法院仍優於政治部門：憲法議題經由已建立且有約束力的程序進入法院，政治部門沒有可比的程序；法院作成裁判時附具意見，政治部門不附（雖然可以）；法院只有九位成員且終身任職；法院不能選擇自己的議程，國會可以\pg{477 頁}。這些與其他特徵是否使法院在最終經驗分析上比國會或行政部門更穩定，「is more of an empirical question than we first acknowledged」〔比我們起初承認的更是經驗問題〕，但在缺乏接近解決此問題的非軼事性證據之前，兩人仍不相信國會能比最高法院提供更多為定紛而定紛的定紛\pg{477 至 478 頁}。

\subsection{（十三）最高法院本身是否受 Cooper 規則拘束\pg{477 頁，註 62}}

註 62 提出問題：最高法院自己是否受 Cooper 規則\ix{terms}{01@\zhlidxcat{司法審查權限之歸屬}!Cooper 規則}拘束，也就是它現在認為錯誤的裁判是否仍永久地定紛\ix{terms}{02@\zhlidxcat{定紛止爭論}!定紛止爭}該議題，水平先例的強度應有多大。註中列三種立場：

\begin{enumerate}
\def\labelenumi{\arabic{enumi}.}
\tightlist
\item
  Cooper\ix{cases}{02@美國!01@聯邦最高法院!1958Cooper v. Aaron（1958）@Cooper v. Aaron（1958）} 規則完全適用於法院自己。法院一旦裁判憲法事項，即永久確定；不論法院現在認為先前的裁判多麼錯誤、有害或不公正，先前裁判對法院和所有其他官員都與憲法一樣具有權威。
\item
  法院是唯一免於 Cooper 規則的機構；認為先前憲法先例錯誤時，可以自由推翻。
\item
  法院僅在先例既是錯誤的憲法解釋、又經法院認為不公正或有害時，才可推翻；憲法錯誤是必要但不充分的條件。
\end{enumerate}

兩人評論第一種立場將定紛置於一切對立價值之上，法院只有一次正確解釋特定憲法條文的機會，補救最高法院錯誤的唯一方式是漫長而困難的修憲\ix{terms}{03@\zhlidxcat{憲法之效力基礎}!修憲}。第二種則不夠重視定紛。兩人傾向第三種的某個變體，因為它給法院一道對有害憲法錯誤的安全閥。第三種也有缺點：法院須作道德與政策上的一階\ix{terms}{02@\zhlidxcat{定紛止爭論}!一階定紛與二階定紛}判斷，並犧牲部分定紛價值、容許部分憲法錯誤。故兩人對第三種立場的支持「less than wholehearted and quite tentative」〔不全心，且頗為保留〕\pg{477 頁，註 62}。

\section{第 III 部分：最高法院的功能\pg{478 至 481 頁}}
\label{sec:AS-4}

\subsection{（一）案件裁判者的圖像與 Hartnett\pg{478 頁，註 63}}

第 III 部分指出，許多批評的基底，是以案件裁判者\ix{terms}{04@\zhlidxcat{法院之功能}!案件裁判者與造法者}看待最高法院的圖像，造法者的身分不在其內\pg{478 頁}。兩人說這個圖像符合傳統的法院功能觀，不因此必然無效；它在 Hartnett\ix{persons}{hartnett@Hartnett, Edward A.} 的批評中最明確，Hartnett 認為法院所做的是裁判具體案件，其餘包括意見在內都是「gravy」〔佐料〕。對 Hartnett，司法功能的核心是裁判案件，提出意見居於其外；兩人的整個立場既然建立在意見可被視為提供某種指引之上，在 Hartnett 看來便是誤解了裁判的性質。Hartnett 並以令人印象深刻的歷史論證支持，主張最高法院自建國之初就被構想為案件裁判者，沒有被構想為造法者或意見提供者\pg{478 頁，註 63 列 Hartnett 126--36、146--58 頁，並指 Devins\ix{persons}{devins@Devins, Neal} \& Fisher\ix{persons}{fisher@Fisher, Louis} 91 頁有簡短相似主張，另列 Eisgruber、Merrill 與 Wechsler\ix{persons}{wechsler@Wechsler, Herbert}}。

\subsection{（二）每年受理案件數與法院的角色\pg{478 至 479 頁，註 64--65}}

兩人不爭論 Hartnett\ix{persons}{hartnett@Hartnett, Edward A.} 的歷史，但與對待 Devins\ix{persons}{devins@Devins, Neal} 和 Fisher\ix{persons}{fisher@Fisher, Louis} 的歷史一樣，對該歷史的地位有不同看法。把歷史視為有趣但非決定性的，則可問何種法院應做什麼的圖像，最符合兩人的定紛\ix{terms}{02@\zhlidxcat{定紛止爭論}!定紛止爭}功能觀\pg{478 頁}。他們說去年最高法院依 certiorari 程序以其自己對慎重審酌的定義，慎重審酌的案件只有 94 件，被請求慎重審酌而選擇不審的有 7015 件（註 64 引 \en{The Supreme Court, 1998 Term---The Statistics}, 113 Harv. L. Rev. 400, 407）。若最高法院主要是案件裁判者\ix{terms}{04@\zhlidxcat{法院之功能}!案件裁判者與造法者}，意見提供者至多居次，它「appears to have achieved a level of inefficiency unmatched in any other government agency we know of」〔似乎達到了我們所知任何其他政府機構都比不上的低效率〕。若法院只處理 1\% 的案件量，以求在這 1\% 上作出好的裁判，並附上依 Hartnett 等人看來不太重要的意見，則法院的角色與運作方式「the time has come for major reconsideration」〔該到全面重新考慮的時候了〕\pg{478 至 479 頁}。

註 65 引 Estreicher 與 Sexton 的 \en{A Managerial Theory of the Supreme Court\textquotesingle s Responsibilities} 主張最高法院應更把自己視為下級法院體系的管理者，兩人不反對，另提出更激進的提議：法院也應把自己視為「manager of constitutional signals」〔憲法訊號的管理者〕與「manager of a system of constitutional compliance」〔憲法遵從體系的管理者〕。

\subsection{（三）意見的數量與品質\pg{479 頁，註 66--68}}

兩人說重新考慮或許適當，但原因在於法院寫的意見不夠，或寫得不夠好，與寫了太多意見無關\pg{479 頁，註 66 引 Easterbrook 與 Schauer\ix{persons}{schauer@Schauer, Frederick} 的 Opinions as Rules}。法院作為憲法解釋者是兩人的圖像，與 Hartnett\ix{persons}{hartnett@Hartnett, Edward A.} 的法院作為爭端解決者\ix{terms}{04@\zhlidxcat{法院之功能}!爭端解決者}的圖像相對；前者看到法院需要比目前意見的數量與風格所顯示的，更關心指導、引導、協助以及命令其他機構與部門（註 67 稱此與 Post \& Salop 的看法大致一致）。

如果這是或應當是法院關心的事，兩人預期會看到更清楚的規則、更少沒有多數意見的分歧判決、大法官更關心法院以單一聲音發言，甚至勝過強調自己的論點或堅持自己對結果的看法，更重視 stare decisis，總之更多符合法院作為造法者\ix{terms}{04@\zhlidxcat{法院之功能}!案件裁判者與造法者}的行為，而法院「undeniably is, and in our judgment inevitably must be」〔無可否認就是，且依我們的判斷無可避免必然是〕造法者\pg{479 頁，註 68 引 Rogers，並指 Sherry\ix{persons}{sherry@Sherry, Suzanna} 同意兩人對 Cooper\ix{cases}{02@美國!01@聯邦最高法院!1958Cooper v. Aaron（1958）@Cooper v. Aaron（1958）} 的立場，並認為兩人的立場意味著個別大法官在許多脈絡，例如 City of Boerne\ix{cases}{02@美國!01@聯邦最高法院!1997City of Boerne v. Flores（1997）@City of Boerne v. Flores（1997）}，應如同他們期待其他官員那樣受法院裁判拘束}。

\subsection{（四）Hartnett 論證的最強處\pg{479 至 480 頁，註 69--70}}

Hartnett\ix{persons}{hartnett@Hartnett, Edward A.} 立場最強的論證，在兩人看來出自法院目前作為規則制定者的表現，他的歷史分析雖然正確，但大致無關\pg{479 至 480 頁}。若對權威規則的權威解釋是規則制定的功能，兩人相信如此，則法院的許多運作方式需要改革。兩人說，Hartnett 最有說服力的主張，即使不是他明白提出的，是：依他的圖像法院可以照現在這樣繼續，依兩人的圖像則須花更多時間思考如何對國會、行政部門、下級法院與各州給予更好的指引\pg{480 頁}。

註 69 說明兩人不接受 Hartnett 認為法院「have little business decreeing」〔不太有權宣示〕規則，也不同意意見，也就是理由，無法被遵循。理由作為對某類案件（包含眼前這件）應如何處理的陳述，本質上就是規則，可以如規則般被遵循；例如「I hold in your favor because people have a right to protest in public parks」〔我判你勝訴，因為人們有權在公園抗議〕這句話所含的理由，即構成一條公園內抗議必須獲准的規則。在兩人看來，法院每提出一個判決的理由就是制定規則，所以每寫一份意見都是；主張法院應退出這門事業才是真正激進的主張，他們的建議較溫和，只是請法院承認它不可避免地從事的規則制定，作為更好地做它現在做、且已做了數個世紀之事的第一步。

註 70 補充：若從頭開始，這也可能是要求一部不那麼晦澀的憲法文本的理由。南非 114 頁的新憲法與巴西 211 頁的新憲法也許是極端例子，但當今世界所有制憲過程中，沒有國家選擇像美國憲法那樣簡短晦澀的成文憲法。部分原因可能是美國憲法在祕密中制定，而多數現代憲法，尤其是這兩部，是在非常公開的情形下制定，政治誘因傾向全面納入各種權利、利益與價值；部分也可能是人們愈益認識到，憲法既要作為法律，便不足為奇會看起來像法律。

\subsection{（五）最高法院自身的宣稱\pg{480 至 481 頁，註 71--72}}

兩人把最高法院自己宣稱的解釋至上\ix{terms}{01@\zhlidxcat{司法審查權限之歸屬}!解釋至上}，視為與原始意圖同樣無關，因為這種自我提拔式的論證無法回答其他部門是否應尊重\ix{terms}{05@\zhlidxcat{解釋理論與自制}!尊重}該宣稱，註 71 說明此外並無來自司法以外的政治誘因支撐，例如不理會即使不受歡迎的最高法院裁判的民選官員所要承擔的政治責任，在現行政治環境中顯然不存在\pg{480 頁}。然而另一意義上，法院自己的宣稱高度相關：如果法院自稱是其他官員與部門的引導者，Cooper\ix{cases}{02@美國!01@聯邦最高法院!1958Cooper v. Aaron（1958）@Cooper v. Aaron（1958）} 與 City of Boerne\ix{cases}{02@美國!01@聯邦最高法院!1997City of Boerne v. Flores（1997）@City of Boerne v. Flores（1997）} 的陳述似乎如此主張，則法院的表現可以適當地以其自己宣布的功能衡量\pg{480 至 481 頁，註 72 引 Krotoszynski 論法院如何更有效地履行此功能}。依兩人所支持的功能定義，法院「plainly has a long way to go」〔顯然還有很長的路要走〕\pg{481 頁}。

\subsection{（六）諮詢意見禁令與裁判理由\pg{481 頁，註 73--74}}

兩人預期想像這個角色似乎與禁止諮詢意見\ix{terms}{04@\zhlidxcat{法院之功能}!諮詢意見}相抵觸（註 73 列 Dorf、Field、Frankfurter 與 Clovis \& Updegraff）。他們答覆：若諮詢意見禁令被理解為禁止對於案件結果不必要的、關於法律的有用陳述，法院將只剩宣告判決，或許，僅僅是或許，附上案件事實的陳述。法院一旦超出此範圍，不論稱之為提出意見、說明理由或附帶意見，都在從事提供諮詢意見的事業。禁止法院在個案或爭議之外這樣做，可能有良好的第 3 條理由\ix{statutes}{02@美國!01@聯邦憲法!003@第 3 條}（註 74 引 Krotoszynski 與 Lee 關於 mootness 的文章），但個案或爭議一旦存在，要消除法院所從事的全部說理實踐已為時過晚，這些實踐必然「decide」眼前案件以外的案件，因為理由在邏輯上比其所支持的結果更寬廣。法院既然已經且必然在其意見中提供建議，也就是諮詢意見，則想像它可以比至今更認真考慮如何更有效地履行這項功能，只是「the smallest of steps」〔最小的一步〕\pg{481 頁}。

\section{結論\pg{481 至 482 頁}}
\label{sec:AS-5}

\subsection{（一）Marshall 的主張與法的統一性\pg{481 至 482 頁，註 75--77}}

結論起首說，John Marshall\ix{persons}{marshall@Marshall, John（馬歇爾）} 主張司法部門的職責是「to say what the law is」〔說明法律是什麼〕（註 75 引 Marbury\ix{cases}{02@美國!01@聯邦最高法院!1803Marbury v. Madison（1803）@Marbury v. Madison（1803）} v. Madison\ix{persons}{madison@Madison, James（麥迪遜）}, 5 U.S. (1 Cranch) 137, 177 (1803)），如今流行將它視為 Marshall 式傲慢的典型，對 Marbury 一案的結果不必要，憲法未予支持，政策上也不明智\pg{481 至 482 頁，註 76 對照 Nagel\ix{persons}{nagel@Nagel, Robert F.} 稱司法至上\ix{terms}{01@\zhlidxcat{司法審查權限之歸屬}!司法至上}的主張為「self-importance」〔自重〕的表現}。

兩人的答覆是：若有多個機構說明法律是什麼，就可能有多部法律，更準確地說，對同一部法律有多種解釋\pg{482 頁}。若如 Fuller\ix{persons}{fuller@Fuller, Lon L.（富勒）} 所主張，知道法律是什麼與知道如何遵從是法治本身的必要條件（註 77 引 Fuller, \en{The Morality of Law}），則繁多的法律與繁多的解釋與法治本身相悖。他們承認意見的多樣是有價值的社會現象，但法律的存在，主要是因為意見與行動的多樣有時產生的害處大於好處；當重視協調與定紛\ix{terms}{02@\zhlidxcat{定紛止爭論}!定紛止爭}多於多樣性時，就使用法律。Marshall 的主張在兩人看來「nothing less than」〔不外乎〕後來由 Fuller 精煉的這項觀察：沒有單一而有權威的解釋者，法律與哲學書籍及建議專欄裡許多不被執行的指示之間，將差別甚微\pg{482 頁}。

\subsection{（二）讓步與結語\pg{482 頁，註 78}}

兩人承認，Marshall\ix{persons}{marshall@Marshall, John（馬歇爾）} 作為法官，從需要單一權威解釋者過快地推到司法部門是唯一能擔任此角色的機構。其他機構雖然是候選者，但沒有一個具備能使其以與最高法院同等程度的一致性發言的制度性約束。只要如此，而且沒有證據顯示並非如此，則說明法律是什麼就將、也應是最高法院的功能\pg{482 頁}。註 78 指出，與批評者交換意見最有價值的成果之一，是使辯論的核心，至少是兩人與主張由國會擔任單一權威解釋者的人之間的核心，經驗問題的面貌更為清楚；如 Peabody\ix{persons}{peabody@Peabody, Bruce G.} 所觀察，這場辯論可能成為嚴肅經驗研究的跳板，而這種研究尚未進行\pg{註 78 引 Peabody 85--90 頁}。

最後一段區分兩種位置。對向最高法院論辯、或評價其工作的人，最高法院所說與其所稱解釋的法律之間，始終存在某種概念空間。對最高法院向其發言的人，「the functions of law have been better served, and will be better served in the future, if the Constitution is what the Supreme Court says it is」〔如果憲法就是最高法院所說的那樣，法律的功能過去得到了更好的實現，將來也會得到更好的實現〕\pg{482 頁}。

% 自 notes/docs/1151/seminar/W09_1109_預習.md（87808f3）由 tools/md_to_tex.py 於 2026-11-09 轉出一次；此後本檔為正本。
\chapter{Kramer〈Judicial Supremacy and the End of Judicial Restraint〉\pg{621 至 634 頁}}
\label{ch:Kramer}

\begin{longtable}[]{@{}
  >{\raggedright\arraybackslash}p{(\linewidth - 4\tabcolsep) * \real{0.1919}}
  >{\raggedright\arraybackslash}p{(\linewidth - 4\tabcolsep) * \real{0.5455}}
  >{\raggedright\arraybackslash}p{(\linewidth - 4\tabcolsep) * \real{0.2626}}@{}}
\caption{Kramer 一文的節次與頁次。}\tabularnewline
\toprule\noalign{}
\begin{minipage}[b]{\linewidth}\raggedright
節次
\end{minipage} & \begin{minipage}[b]{\linewidth}\raggedright
小標
\end{minipage} & \begin{minipage}[b]{\linewidth}\raggedright
頁次
\end{minipage} \\
\midrule\noalign{}
\endfirsthead
\toprule\noalign{}
\begin{minipage}[b]{\linewidth}\raggedright
節次
\end{minipage} & \begin{minipage}[b]{\linewidth}\raggedright
小標
\end{minipage} & \begin{minipage}[b]{\linewidth}\raggedright
頁次
\end{minipage} \\
\midrule\noalign{}
\endhead
\bottomrule\noalign{}
\endlastfoot
甲 & 摘要與論題 & 621--622 \\
乙 & 建國時期的兩種司法審查\ix{terms}{01@\zhlidxcat{司法審查權限之歸屬}!司法審查}理論 & 622--624 \\
丙 & 解釋方法\ix{terms}{05@\zhlidxcat{解釋理論與自制}!解釋理論}與自制\ix{terms}{05@\zhlidxcat{解釋理論與自制}!司法自制} & 624--626 \\
丁 & 1800 年至新政 & 626--630 \\
戊 & Brown 案\ix{cases}{02@美國!01@聯邦最高法院!1954Brown v. Board of Education（1954）@Brown v. Board of Education（1954）}之後的自由派與保守派 & 630--632 \\
己 & 解釋理論的興起與自制的衰退 & 632--634 \\
\end{longtable}

\section{摘要與論題\pg{621 至 622 頁}}
\label{sec:Kramer-1}

\subsection{（一）首頁摘要\pg{621 頁}}

首頁摘要寫道，Thayer\ix{persons}{thayer@Thayer, James Bradley} 重申的是較早的 Jeffersonian\ix{terms}{01@\zhlidxcat{司法審查權限之歸屬}!共和派} 觀念：「primary authority to interpret the Constitution lies with the people and not with courts」〔詮釋憲法的首要權威在人民，不在法院〕。這種 popular constitutionalism\ix{terms}{01@\zhlidxcat{司法審查權限之歸屬}!人民憲政主義}〔人民憲政主義〕被現代的司法至上（judicial supremacy\ix{terms}{01@\zhlidxcat{司法審查權限之歸屬}!司法至上}）取代，後者「took hold only in the 1960s」〔遲至 1960 年代才站穩〕，使政治辯論從誰應詮釋憲法轉為憲法應如何詮釋。司法至上被接受之後，自制\ix{terms}{05@\zhlidxcat{解釋理論與自制}!司法自制}只剩一項理由，即保護司法權威，免於因行動過於頻繁或積極而招致政治反撲；摘要稱這項理由「grows weaker in direct proportion to supremacy\textquotesingle s acceptance」〔隨司法至上被接受的程度而等比減弱〕。

摘要接著引 Posner\ix{persons}{posner@Posner, Richard A.} 的觀察：司法自制今日已退化為「little more than a rhetorical tool used opportunistically by both sides」〔雙方機會主義地使用的修辭工具而已〕。Kramer\ix{persons}{kramer@Kramer, Larry D.} 對成因的看法，摘要中的原文是：「the culprit for this development is not, as Posner suggests, the development of theories of constitutional interpretation」〔此一發展的禍首，並非如 Posner 所言，是憲法解釋理論\ix{terms}{05@\zhlidxcat{解釋理論與自制}!解釋理論}的發展〕，禍首在「our embrace of the notion that deciding the ongoing meaning of a democratic Constitution is a task best left entirely to judges」〔我們接受了這個想法：決定民主憲法持續演變的意義，最宜全交給法官〕。

\subsection{（二）對 Posner 的回應\pg{621 至 622 頁}}

本文開頭以 Posner\ix{persons}{posner@Posner, Richard A.} 與 Holmes\ix{persons}{holmes@Holmes, Oliver Wendell, Jr.（霍姆斯）} 並舉，稱 Posner 之於二十世紀末與二十一世紀初，如 Holmes 之於十九世紀末與二十世紀初\pg{621 頁}。622 頁接著說，Posner 的論述有趣，不等於完全正確；Kramer\ix{persons}{kramer@Kramer, Larry D.} 認為 Posner「misunderstood important historical nuances」〔誤解了重要的歷史細微處〕，因此要講「a somewhat different story」〔一個略有不同的故事〕，提出可以並列考慮的另一組因素。被回應的文章是 Posner, \en{The Rise and Fall of Judicial Self-Restraint}, 100 Calif. L. Rev. 519 (2012)，註 8\pg{624 頁}列有出處。首頁註 * 僅載明 Kramer 為 Stanford Law School 的 Richard E. Lang Professor and Dean。

\subsection{（三）自制的興衰與起點\pg{622 頁}}

本文認為自制\ix{terms}{05@\zhlidxcat{解釋理論與自制}!司法自制}「rose but has fallen」〔興起而後衰落〕一事「seems plain enough」〔相當明顯〕。衰落的原因不能歸於左派或右派，只能說無論哪一方掌權，都似乎決意超過對方上一次的作為，因此這是「an accelerating slide of many years」〔多年來不斷加速的下滑〕。至於興起，Posner\ix{persons}{posner@Posner, Richard A.} 追溯到 Thayer\ix{persons}{thayer@Thayer, James Bradley} 1893 年的 \en{The Origin and Scope of the American Doctrine of Constitutional Law}（註 1：7 Harv. L. Rev. 129）；Kramer\ix{persons}{kramer@Kramer, Larry D.} 回答這「much too late」〔太晚了〕，要理解自制作為對政治決策者的尊重（deference\ix{terms}{05@\zhlidxcat{解釋理論與自制}!尊重}）之興衰，須回到建國時期與司法審查\ix{terms}{01@\zhlidxcat{司法審查權限之歸屬}!司法審查}的起源，Thayer 自己也是如此。司法審查如何發展，Kramer 指向自己的 \en{The People Themselves}（2004）第 9 至 144 頁，另列 Friedman\ix{persons}{friedman@Friedman, Barry} 與 Hamburger 兩部近著（註 2）。

\section{建國時期的兩種司法審查理論\pg{622 至 624 頁}}
\label{sec:Kramer-2}

\subsection{（一）共和派：人民憲政主義\pg{622 至 623 頁，註 3 至 5}}

Kramer\ix{persons}{kramer@Kramer, Larry D.} 稱建國時期出現兩套正當化司法審查\ix{terms}{01@\zhlidxcat{司法審查權限之歸屬}!司法審查}的理由，並稱此二說之出現無爭議。第一套稱為「Republican」或「Jeffersonian\ix{terms}{01@\zhlidxcat{司法審查權限之歸屬}!共和派}」，因 Jefferson\ix{persons}{jefferson@Jefferson, Thomas（傑佛遜）} 一派所倡而得名，出現於 1780 年代，前提是「primary authority to interpret, as well as to make, constitutional law rested actively in the community」〔詮釋與制定憲法的首要權威都實際落在社群手中〕，即他在別處所稱的 popular constitutionalism\ix{terms}{01@\zhlidxcat{司法審查權限之歸屬}!人民憲政主義}\pg{註 3，引 \en{The People Themselves} 第 8 頁}\pg{622 頁}。

依這種看法，憲法是人民制定來規範並約束政府的法，普通法則是政府制定來規範並約束人民的法，這個倒置又倒置了解釋與執行權限的通常分配\pg{622 頁}。政府官員是普通法的權威詮釋者；人民只能透過選舉與罷免對立法有間接控制，法律一旦制定，公民在確定其意義與施加制裁上居於官員之下，自己的詮釋只是「guesses」〔臆測〕，沒有正式法律權威\pg{622 至 623 頁}。

憲法的關係則相反。Kramer 引共和派的 William Nelson 法官在 \en{Kamper v. Hawkins\ix{cases}{02@美國!02@維吉尼亞州一般法院!1793Kamper v. Hawkins（1793）@Kamper v. Hawkins（1793）}}, 3 Va. (1 Va. Cas.) 20, 24 (1793) 的話：「{[}a{]} Constitution . . . is to the governors, or rather to the departments of government, what a law is to individuals」〔憲法之於統治者，或毋寧說之於政府各部門，如同法律之於個人〕\pg{623 頁，註 4}。憲法的對象是公職人員，官員須盡力詮釋並遵循憲法，他們的詮釋與人民對普通法的詮釋一樣不具權威，受「the People themselves」〔人民自己〕這一更高權威的直接監督與矯正\pg{623 頁，註 5}。

\subsection{（二）聯邦派：司法至上的先聲\pg{623 頁，註 6}}

第二套理論出現在 1790 年代，憲法批准之後，由 Federalist 領袖發展。Kramer\ix{persons}{kramer@Kramer, Larry D.} 把它放在更大的思想演變中：憲法開始運作，法國大革命爆發，Federalist 為了在看來民主失控的局面中確保「order」〔秩序〕，重新思考司法權的性質，形成「what amounts to the modern theory of judicial supremacy\ix{terms}{01@\zhlidxcat{司法審查權限之歸屬}!司法至上}」〔相當於現代司法至上理論的主張〕\pg{623 頁}。

這套理論把憲法視為「nothing more than a species of ordinary law」〔不過是普通法的一種〕，憲法解釋因此是普通的司法與法律行為，不授權給立法或行政官員，更不授權給全體人民。政治部門行動前仍須自行詮釋憲法，依權力分立\ix{terms}{06@\zhlidxcat{審查制度之類型}!權力分立體制}，首要的解釋權在法院，法院對憲法意義與合憲性有「special and final responsibility」〔特別且終局的責任〕。Kramer 以 James Kent 1794 年的話為例：憲法的解釋或建構「requires the exercise of the same LEGAL DISCRETION as the interpretation or construction of a Law」〔需要與解釋或建構一部法律相同的法律裁量〕，所以法官是「the proper and intended Guardians of our limited Constitutions」〔有限憲法適當而原定的守護者〕\pg{623 頁，註 6}。

\subsection{（三）羅德島議會的答覆\pg{624 頁，註 7}}

624 頁開頭，羅德島議會拒絕 Madison\ix{persons}{madison@Madison, James（麥迪遜）} 集結人民反對 Alien and Sedition Acts 之合憲性的努力，說得「still more baldly」〔更為直白〕：憲法「vests in the federal courts, exclusively, and in the Supreme Court of the United States, ultimately, the authority of deciding on the constitutionality of any act or law of the Congress of the United States」〔把決定國會任何法案或法律合憲性的權限，專屬授予聯邦法院，並以聯邦最高法院為終局〕\pg{註 7，載 Elliot 編 \en{Debates} 第 4 卷 532, 533 頁}。

\section{解釋方法與自制\pg{624 至 626 頁}}
\label{sec:Kramer-3}

\subsection{（一）雙方都未爭論如何解釋憲法\pg{624 頁}}

Kramer\ix{persons}{kramer@Kramer, Larry D.} 就兩派的差異提出兩點。第一點，雙方都沒有花時間討論如何解釋憲法，也沒有興趣發展今日意義的憲法解釋「theory」〔理論〕，因為對這個問題並無真正的分歧。大家基本上相信，憲法可以而且應當以橫跨各法領域的同一種開放的論辯程序來解釋，即相當不拘結構地運用文本、結構、歷史、先例與後果的論據，得出最有說服力的整體結論\pg{624 頁}。

既然對如何解釋看法一致，爭議就集中於誰能解釋、發生什麼效果。Kramer 補充，解釋過程鬆散而開放，不同的解釋者可能對同一問題合理地得出不同結論，所以誰有解釋權限\ix{terms}{01@\zhlidxcat{司法審查權限之歸屬}!解釋權限}才如此關鍵\pg{624 頁}。

\subsection{（二）自制見於兩套理論，Thayer 並非原創\pg{624 頁，註 8}}

第二點與本文主題更直接相關：自制\ix{terms}{05@\zhlidxcat{解釋理論與自制}!司法自制}的概念在兩套司法審查\ix{terms}{01@\zhlidxcat{司法審查權限之歸屬}!司法審查}理論裡都有位置，並且「central to both approaches」〔對兩者都居核心〕。Kramer\ix{persons}{kramer@Kramer, Larry D.} 因此不同意 Posner\ix{persons}{posner@Posner, Richard A.} 的說法：Thayer\ix{persons}{thayer@Thayer, James Bradley} 所憑藉的不只是「scattered judicial remarks」〔零散的司法評論〕，在「organizing and rationalizing」〔整理並使之合理化〕它們上也不是首創\pg{註 8，引 Posner 523 頁}。依 Kramer，Thayer 是要把一個「old and well-developed concept of self-restraint」〔古老而發展成熟的自制概念〕重新推回舞台中央，當時法院及其支持者正傾向忽視它，情形與今日相仿\pg{624 頁}。

\subsection{（三）Posner 的「無理論」前提\pg{624 至 625 頁，註 9}}

Kramer\ix{persons}{kramer@Kramer, Larry D.} 不同意 Posner\ix{persons}{posner@Posner, Richard A.} 認為自制\ix{terms}{05@\zhlidxcat{解釋理論與自制}!司法自制}的前提是「hav{[}ing{]} no theory」〔沒有理論〕、自制之死可歸因於宣稱能產生「correct」〔正確〕憲法決定的解釋理論\ix{terms}{05@\zhlidxcat{解釋理論與自制}!解釋理論}之興起\pg{624 頁，註 9，引 Posner 538 頁}。理由是即使主張只有單一正解的理論，適用上也有足夠的不確定性，使理性的人對何者為正解合理地有歧見，Kramer 說這就足以為尊重\ix{terms}{05@\zhlidxcat{解釋理論與自制}!尊重}或自制創造空間。他並指出十八、十九世紀的律師也相信自己的方法有單一正解\pg{625 頁}。

Kramer 同意 Posner 的是，今日法院與評論者都不太傾向自制。他的主張是這項變化另有來源，解釋理論的發展本身，是削弱自制的同一組發展的「by-product」〔副產品〕，後文再回到這個問題\pg{625 頁}。

\subsection{（四）共和派的自制理由\pg{625 至 626 頁，註 10、11}}

兩派都期待法院在行使司法審查\ix{terms}{01@\zhlidxcat{司法審查權限之歸屬}!司法審查}時自制\ix{terms}{05@\zhlidxcat{解釋理論與自制}!司法自制}，理由與目的不同\pg{625 頁}。對共和派\ix{terms}{01@\zhlidxcat{司法審查權限之歸屬}!共和派}，自制是保存「the People」首要且優越之權威的手段。人民既已把制定法律的責任委託給立法機關，一項法律是否合憲，首先是社群與民選代表之間透過抗議、請願、選舉與其他直接的人民行動來解決的事。少數共和派因此主張任何形式的司法審查都不正當\pg{註 10，引 11 Annals of Cong. 178--80 (1802) Breckinridge 參議員之發言，及同冊 661 頁 Randolph 眾議員之發言}。

共和派的正統立場承認法院有時必須拒絕執行被認為違憲的法律。人民既然委託法院裁判案件，法院在裁判時若不考慮合憲性，就違反了作為人民代理人的責任，與立法者不審酌法律是否合憲同理。法院宣告立法無效時，是以人民代理人的身分行動。遇有疑義，只有該法「unconstitutional beyond dispute」〔違憲而無可爭辯〕時才能宣告無效\pg{註 11，引 Iredell\ix{persons}{iredell@Iredell, James} 1787 年 8 月 26 日致 Spaight 函；並載 Thayer\ix{persons}{thayer@Thayer, James Bradley} 138--42 頁彙整之早期案例，及 Treanor 58 Stan. L. Rev. 455 (2005)}。Kramer\ix{persons}{kramer@Kramer, Larry D.} 說，代理人揣度本人（principal\ix{terms}{01@\zhlidxcat{司法審查權限之歸屬}!本人與代理人}）的指令，不可輕易為之，須待近於確定的條件，因為本人自己能夠行動，且始終保有首要責任\pg{625 至 626 頁}。

\subsection{（五）聯邦派的自制理由\pg{626 頁，註 12 至 16}}

聯邦派\ix{terms}{01@\zhlidxcat{司法審查權限之歸屬}!聯邦派}明確否認人民在詮釋憲法上有首要的、乃至任何的權威（相對於制定新的憲法法）。依 Kramer\ix{persons}{kramer@Kramer, Larry D.}，到 1790 年代中期，聯邦派採取的立場已與今日多數人相近：憲法把終局的解釋權限\ix{terms}{01@\zhlidxcat{司法審查權限之歸屬}!解釋權限}授予法院。法官解釋憲法時不是社群的代理人\ix{terms}{01@\zhlidxcat{司法審查權限之歸屬}!本人與代理人}；他們在阻止社群及其立法與行政代理人犯下 Gouverneur Morris\ix{persons}{morris@Morris, Gouverneur} 所說的「the greatest follies and absurdities」〔最大的愚行與荒謬〕（註 12）。Morris 稱讓人民解釋憲法，就是把「an inestimable treasure . . . into the hands of drunkards, madmen, and fools」〔無價之寶\ldots\ldots 交到醉漢、瘋子與傻瓜手中〕（註 13）。因此憲法刻意「vested in the judges a check」〔把一項制衡授予法官〕（註 14），照 Morris 的嘲諷，是為了「to save the people from their most dangerous enemy, to save them from themselves」〔使人民免於其最危險的敵人，即免於自己〕（註 15）。

聯邦派同樣強調，法律的違憲須屬明確才能不予執行，「in all doubtful cases . . . the Act ought to be supported」〔凡有疑義之案\ldots\ldots 該法應予維持〕\pg{註 16，引 Iredell\ix{persons}{iredell@Iredell, James} 函第 175 頁}。但在聯邦派，這種自制\ix{terms}{05@\zhlidxcat{解釋理論與自制}!司法自制}是審慎\ix{terms}{05@\zhlidxcat{解釋理論與自制}!審慎}與政治權宜的問題，目的是藉由降低越權的風險，保障並維持司法權威，原文括號標明此為 judicial (rather than popular) authority\pg{626 頁}。

\section{1800 年至新政\pg{626 至 630 頁}}
\label{sec:Kramer-4}

\subsection{（一）1800 年選舉與 Marbury\pg{626 頁，註 17、18}}

Kramer\ix{persons}{kramer@Kramer, Larry D.} 寫道，1800 年的選舉決定性地解決了兩種觀點的爭議，由共和派\ix{terms}{01@\zhlidxcat{司法審查權限之歸屬}!共和派}勝出\pg{註 17，引 \en{The People Themselves} 137--38 頁}。\footcite{kramer2004}聯邦派\ix{terms}{01@\zhlidxcat{司法審查權限之歸屬}!聯邦派}又延續幾年，在 1812 年戰爭之後完全消失；John Marshall\ix{persons}{marshall@Marshall, John（馬歇爾）} 在 \en{Marbury\ix{cases}{02@美國!01@聯邦最高法院!1803Marbury v. Madison（1803）@Marbury v. Madison（1803）} v. Madison\ix{persons}{madison@Madison, James（麥迪遜）}} 中明顯放棄聯邦派的司法審查\ix{terms}{01@\zhlidxcat{司法審查權限之歸屬}!司法審查}觀，所作意見是 Jefferson\ix{persons}{jefferson@Jefferson, Thomas（傑佛遜）} 自己可能寫出的，且大部分取自維吉尼亞共和派 Spencer Roane 與 St. George Tucker 在 \en{Kamper v. Hawkins\ix{cases}{02@美國!02@維吉尼亞州一般法院!1793Kamper v. Hawkins（1793）@Kamper v. Hawkins（1793）}} 的意見\pg{626 頁，註 18，另引 Kramer, \en{Marbury and the Retreat from Judicial Supremacy}, 20 Const. Comment. 205, 228--29 (2003)}。

\subsection{（二）聯邦派觀念的殘存\pg{626 至 627 頁，註 19、20}}

選舉很少殺死觀念，所以聯邦派\ix{terms}{01@\zhlidxcat{司法審查權限之歸屬}!聯邦派}的司法審查\ix{terms}{01@\zhlidxcat{司法審查權限之歸屬}!司法審查}概念並未完全消失，存續於一些前聯邦派人士（後來成為 National Republicans，再成為 Whigs）之間，尤其是逐漸職業化的法官與律師\pg{626 頁，註 19，引 \en{The People Themselves} 156--64 頁}。法院對憲法有特別權威、優越於人民與其他部門的想法，因此一直留在國家的政治與法律論述中；每當其倡議者推得夠用力、引起公眾注意，就被再度壓下，Kramer\ix{persons}{kramer@Kramer, Larry D.} 舉 Second Bank of the United States 之爭與 \en{Dred Scott\ix{cases}{02@美國!01@聯邦最高法院!1857Dred Scott v. Sandford（1857）@Dred Scott v. Sandford（1857）}} 引發的爭議為例\pg{627 頁，註 20}。

\subsection{（三）重建之後司法至上再度出現\pg{627 頁，註 21 至 25}}

Kramer\ix{persons}{kramer@Kramer, Larry D.} 列出重建之後司法至上\ix{terms}{01@\zhlidxcat{司法審查權限之歸屬}!司法至上}強勢重現的幾項因素：\en{Dred Scott\ix{cases}{02@美國!01@聯邦最高法院!1857Dred Scott v. Sandford（1857）@Dred Scott v. Sandford（1857）}} 的重量隨時間與法院人事更替而逐漸減輕；共和黨的優勢在 1874 年後讓位給數十年的分裂政府，使大法官較自由行動（註 21，引 Lichtman \& DeCell；Gely \& Spiller, 12 Int\textquotesingle l Rev. L. \& Econ. 45 (1992)）；內戰時期的自由勞動修辭，在保守派手中演變為新的自由放任、反重分配的正統（註 22，引 Foner；Forbath, 1985 Wis. L. Rev. 767）。法院在道德上忠於特定的財產權觀念，又懼怕勞工運動與其他社會動盪的徵兆，開始擴張自己的角色。Brewer 大法官 1893 年的演講提到「{[}m{]}agnifying, like the apostle of old, my office」〔像古代使徒一樣，頌揚我的職務〕，理由是「the salvation of the nation . . . rests upon the independence and vigor of the judiciary」〔國家的救贖\ldots\ldots 繫於司法的獨立與活力〕（註 23）。此時也再度聽到被棄的聯邦派\ix{terms}{01@\zhlidxcat{司法審查權限之歸屬}!聯邦派}舊論，如「what is now to be feared and guarded against is the despotism of the many---of the majority」〔現在該恐懼與防範的，是眾人的專制，即多數的專制〕，以及藉司法審查\ix{terms}{01@\zhlidxcat{司法審查權限之歸屬}!司法審查}「people had effectually protected themselves against themselves」〔人民已有效地保護自己免於自己〕（註 24，引 Dillon 1892 年 ABA 會長演說）。

這些發展逐漸改變法院的憲法判決。因變化緩慢，且起初主要影響已失去強烈政治支持的法律，法院過了幾年才開始受到抨擊；到 1880 年代末，法律評論者與政治領袖已開始注意，一場長期的鬥爭隨之展開，最後在 1930 年代的新政危機（New Deal crisis）告終（註 25，引 Friedman\ix{persons}{friedman@Friedman, Barry}, 76 N.Y.U. L. Rev. 1383 (2001)）。

\subsection{（四）進步時代的爭議\pg{627 至 628 頁，註 26}}

Kramer\ix{persons}{kramer@Kramer, Larry D.} 說，今日的法律人、法官與學者常以為法院的至上地位在這段期間沒有受到挑戰，這是照單全收司法意見修辭所造成的「historical blind spot」〔歷史盲點〕\pg{627 至 628 頁}。大法官關於司法在憲法體制中位置的陳述，「must be seen for what they were: partisan claims in contested territory」〔須視為其本來面目：在爭議地帶的黨派主張〕。進步時代是「flux and uncertainty」〔變動與不確定〕的時期，法院的角色並非核心議題，但法院因積極地在初現的階級衝突中選邊，使自己成為爭議來源，它為自己創造的角色是否適當，成為懸而未決的問題之一\pg{628 頁}。

\subsection{（五）Thayer 的主張\pg{628 頁，註 26 至 28}}

Thayer\ix{persons}{thayer@Thayer, James Bradley} 的文章寫於這場辯論之中，他是眾多提出類似論點者之一（註 26，引 Friedman\ix{persons}{friedman@Friedman, Barry}；T.W. Brown；W.F. Dodd；Meigs 等）。Kramer\ix{persons}{kramer@Kramer, Larry D.} 認為他的特定論點是重申並恢復 Jeffersonian\ix{terms}{01@\zhlidxcat{司法審查權限之歸屬}!共和派} 司法權威觀的優先地位。

Posner\ix{persons}{posner@Posner, Richard A.} 說 Thayer 最大的憂慮是：法院若把憲法限制執行到底，立法者將不再思考擬議立法的合憲性，只想法院會如何反應\pg{註 27，引 Posner 524 頁}。Kramer 說 Thayer 確實提過這個論點，但僅在文章最後一段帶過；把它視為 Thayer 的主要關切，是「read him anachronistically」〔以時代錯置的方式閱讀他〕。Kramer 的讀法是：從文章結構與 Thayer 重複該點的次數可見，他的主要關切是重申詮釋憲法的首要權威在法院之外，宣告法律違憲的司法權威至多是「subordinate, secondary check」〔從屬的、次要的制衡〕\pg{註 28，引 Thayer 第 136、144、148、150、152 頁}。Thayer 要恢復的是共和派較古老、歷史上居首的司法權威觀，包括其自制\ix{terms}{05@\zhlidxcat{解釋理論與自制}!司法自制}與尊重\ix{terms}{05@\zhlidxcat{解釋理論與自制}!尊重}的主張，並拒絕鍍金時代法院對憲法至上的僭越。

\subsection{（六）Thayer 對 Jeffersonian 理論的改動\pg{628 至 629 頁}}

Kramer\ix{persons}{kramer@Kramer, Larry D.} 承認 Thayer\ix{persons}{thayer@Thayer, James Bradley} 手中的 Jeffersonian\ix{terms}{01@\zhlidxcat{司法審查權限之歸屬}!共和派} 理論有所改變。最主要的是，早期共和國所假定且廣被接受的人民直接監督，已不再說得通；Thayer 改以政治部門，尤其立法機關，為人民意見的最佳表達\pg{628 頁}。「The crucial point」〔關鍵〕與 Thayer 的主要對象，是否認憲法的解釋權首先授予法院。Thayer 主張的真相正好相反，立法機關對合憲性的判斷應受尊重\ix{terms}{05@\zhlidxcat{解釋理論與自制}!尊重}，法院的角色至多是在立法機關明顯逾越理性之人所能說是允許的界限時介入，此為最低限度的角色\pg{628 至 629 頁}。

\subsection{（七）對立法機關能力的保留\pg{629 頁，註 29、30}}

Kramer\ix{persons}{kramer@Kramer, Larry D.} 另就 Posner\ix{persons}{posner@Posner, Richard A.} 把 Thayer\ix{persons}{thayer@Thayer, James Bradley} 對立法機關的信賴斥為天真，提出兩點保留，並說這個題目足以另文討論\pg{629 頁}。

\begin{enumerate}
\def\labelenumi{\arabic{enumi}.}
\tightlist
\item
  立法機關無力認真對待憲法這項假定，如同 Posner 的處理所示，比司法至上\ix{terms}{01@\zhlidxcat{司法審查權限之歸屬}!司法至上}更被視為理所當然，也更少被法律學者檢驗\pg{註 29，引 \en{The People Themselves} 238--39 頁；Whittington, 80 N.C. L. Rev. 773, 819--22 (2002)；Burgess；Morgan}。Kramer 稱幕僚與委員會花數百小時理解、辯論並解決憲法問題，所用時間比法院多、資訊比法院好；法院只花一小時聽言詞辯論，評議時間更短，實際審議甚少，多數工作由一名沒有經驗的助理在僅有黨派書狀的情形下完成。高層次的議場辯論或許已成過去，但政治學文獻顯示，認真的審議在委員會、黨團、個人之間、電子郵件與幕僚中進行（註 30，引 Tushnet\ix{persons}{tushnet@Tushnet, Mark}, \en{Taking the Constitution Away from the Courts} 65--66 (1999) 等）。
\item
  立法者會為達政治上想要的結果而行事，不能因此否定這些工作的成果；政治固然居首要地位，國會審議涉及的遠不止於政治。Kramer 也不認為典型最高法院大法官比典型國會議員更不會操縱結果，不明白何以迎合意識形態（他說這顯然驅動大法官的決定）比迎合政治更不令人憂慮，「These days, I\textquotesingle m not even sure they are different」〔如今我甚至不確定兩者有何不同〕\pg{629 頁}。
\end{enumerate}

\subsection{（八）進步派與新政\pg{629 至 630 頁，註 31 至 34}}

回到原來的司法審查\ix{terms}{01@\zhlidxcat{司法審查權限之歸屬}!司法審查}概念：它關切以自制\ix{terms}{05@\zhlidxcat{解釋理論與自制}!司法自制}保存人民對憲法的持續控制，這個概念發生了什麼？Kramer\ix{persons}{kramer@Kramer, Larry D.} 說支持者與司法至上\ix{terms}{01@\zhlidxcat{司法審查權限之歸屬}!司法至上}的支持者的鬥爭延續到二十世紀初幾十年，是 1790 年代共和派\ix{terms}{01@\zhlidxcat{司法審查權限之歸屬}!共和派}與聯邦派\ix{terms}{01@\zhlidxcat{司法審查權限之歸屬}!聯邦派}衝突的一次「close, if prolonged, replay」〔接近而冗長的重演〕\pg{629 至 630 頁}。Thayer\ix{persons}{thayer@Thayer, James Bradley} 與其他進步派（包括常被忽略的 Theodore Roosevelt\ix{persons}{roosevelt@Roosevelt, Franklin D.（羅斯福）}，註 31）大體成功地對抗法院的保守派多數，該多數有時言語大膽，偶爾依其修辭行動，但多半默許州與聯邦各層級管制的穩定成長。1920 年代法院稍為大膽，Kramer 稱 \en{Lochner\ix{cases}{02@美國!01@聯邦最高法院!1905Lochner v. New York（1905）@Lochner v. New York（1905）}} 時代最高點仍最宜描述為法院對新型國家干預與管制的「grudging acceptance」〔勉強接受〕\pg{630 頁，註 32}。

法院因 New Deal 恐慌而試圖劃線時，為時已晚，遭到對抗、威脅並被迫退讓。Kramer 引 Roosevelt 的話：「The Constitution of the United States」是「a layman\textquotesingle s document, not a lawyer\textquotesingle s contract」〔外行人的文件，不是律師的契約〕，其意義與適用應由「the people and the Congress」〔人民與國會〕決定並確定，律師與法官則不在其列\pg{630 頁，註 33，引 1937 年 9 月 17 日憲法日演說}。Roosevelt 隨後藉任命權鞏固勝利，使法院回到願意在當時重要的所有問題上尊重\ix{terms}{05@\zhlidxcat{解釋理論與自制}!尊重}政治部門的狀態\pg{630 頁}。

\section{Brown 案之後的自由派與保守派\pg{630 至 632 頁}}
\label{sec:Kramer-5}

\subsection{（一）Brown 案與司法至上的普遍接受\pg{630 至 631 頁，註 34、35}}

Kramer\ix{persons}{kramer@Kramer, Larry D.} 說轉折「no surprise here」〔不意外〕始於 \en{Brown v. Board of Education\ix{cases}{02@美國!01@聯邦最高法院!1954Brown v. Board of Education（1954）@Brown v. Board of Education（1954）}}。\en{Brown} 之後，司法至上\ix{terms}{01@\zhlidxcat{司法審查權限之歸屬}!司法至上}的觀念似乎終於逐漸獲得廣泛的公眾接受；法院的判決仍常有爭議，但自 1960 年代起，抵抗與不服從轉為抗議的形式，不再主張解釋上的平等，更不用說優越。到 1980 年代，涉及憲法的抗議多半朝向法院（directed to），不再針對法院（against），司法至上的接受成為常態\pg{630 頁，註 34}。法官任命所受注意的急遽增加（註 35，引 Steigerwalt、Wittes、Stras），Kramer 稱為在司法決定之下、任命成為重塑憲法唯一實際途徑的世界中，「a thoroughly predictable development」〔完全可預期的發展〕\pg{630 至 631 頁}。解釋這個「rather extraordinary development」〔相當特殊的發展〕不容易，也在本文範圍之外；Kramer 指出一項明顯重要的來源是 Warren Court 本身，它是美國史上第一次給進步派理由把司法視為朋友\pg{631 頁}。

\subsection{（二）自由派\pg{631 頁，註 36、37}}

自由派知識分子在實質上強烈支持法院所為，但他們的師長與英雄曾領導對抗 \en{Lochner\ix{cases}{02@美國!01@聯邦最高法院!1905Lochner v. New York（1905）@Lochner v. New York（1905）}}，許多人把自己的職業生涯投入這個想法：法院干預人民意志是不當的。\en{Brown} 與隨後其他司法創新，因此是對自由派傳統上承諾司法自制\ix{terms}{05@\zhlidxcat{解釋理論與自制}!司法自制}的「wrenching test」〔痛苦的考驗〕，起初自由派難以決定如何回應 Warren Court（註 36，引 Learned Hand, \en{The Bill of Rights}；Friedman\ix{persons}{friedman@Friedman, Barry}；Peller；Wechsler\ix{persons}{wechsler@Wechsler, Herbert}, \en{Toward Neutral Principles of Constitutional Law}, 73 Harv. L. Rev. 1 (1959)；註 37 另引 Horwitz、Kalman 等）。

Warren Court 的積極作為在 1960 年代中期達到高峰時，新一代自由派學者放棄反對法院，「turned the liberal tradition on its head」〔把自由派傳統倒轉過來〕，接受廣泛司法權威的哲學。這些年輕學者雖熟知 \en{Lochner} 的歷史，卻沒有親歷；對他們而言，當時所稱的 ``the countermajoritarian\ix{terms}{05@\zhlidxcat{解釋理論與自制}!反多數決困難} difficulty''〔反多數決困難〕看來「just did not seem so difficult」〔並不那麼困難〕。他們擱置對課責的疑慮，承認司法審查\ix{terms}{01@\zhlidxcat{司法審查權限之歸屬}!司法審查}可能與民主有緊張關係，並以法院能促進更重要的社會正義為由，證成這項取捨\pg{631 頁}。

\subsection{（三）保守派\pg{631 至 632 頁，註 38 至 40}}

保守派對司法至上\ix{terms}{01@\zhlidxcat{司法審查權限之歸屬}!司法至上}沒有問題。Kramer\ix{persons}{kramer@Kramer, Larry D.} 說這是他們自聯邦派\ix{terms}{01@\zhlidxcat{司法審查權限之歸屬}!聯邦派}時代以來一直接受的想法；他們對 Warren Court 的不滿僅在於其決定是錯的，並抗議自由派大法官以憲法為掩護處理憲法其實不涉及的事項\pg{631 頁}。從 Robert Bork 1968 年在 \en{Fortune} 雜誌對法院的攻擊起（註 38），許多保守派開始使用傳統上屬於自由派的民主與自制\ix{terms}{05@\zhlidxcat{解釋理論與自制}!司法自制}修辭\pg{632 頁}。Kramer 說，他們並未主張最高法院應就何者合憲而尊重\ix{terms}{05@\zhlidxcat{解釋理論與自制}!尊重}政治機關，因為他們支持的是司法至上，所對的是他們所理解與解釋的憲法（註 39，引 Epstein、Macedo、Powell）。他們也一直願意宣告違反原始理解的法律無效，如財產權與槍枝（註 40，引 \en{Kelo v. City of New London\ix{cases}{02@美國!01@聯邦最高法院!2005Kelo v. City of New London（2005）@Kelo v. City of New London（2005）}}, 545 U.S. 469, 505--23 (2005) (Thomas, J., dissenting)；\en{District of Columbia v. Heller\ix{cases}{02@美國!01@聯邦最高法院!2008District of Columbia v. Heller（2008）@District of Columbia v. Heller（2008）}}, 554 U.S. 570 (2008)）。

他們對憲法範圍的看法不同，認為範圍較窄，反對自由派大法官以憲法在其認為憲法未規範的領域制定政策，這些領域因此應屬一般政治。Kramer 結論是，美國史上第一次，保守派與自由派在司法至上的原則上達成一致：法官應當擔任解釋者，且其解釋應為終局且有拘束力\pg{632 頁}。

\section{解釋理論的興起與自制的衰退\pg{632 至 634 頁}}
\label{sec:Kramer-6}

\subsection{（一）爭議的焦點由誰解釋轉為如何解釋\pg{632 至 633 頁}}

此時雙方歧見改為如何解釋。Kramer\ix{persons}{kramer@Kramer, Larry D.} 回到開頭所述：在此之前，並無特別的憲法解釋理論\ix{terms}{05@\zhlidxcat{解釋理論與自制}!解釋理論}，存在的是未經檢驗的共識，即憲法可以也應當像其他法源一樣，用律師從其他脈絡熟悉的、經過檢驗的法律論證方法解釋，至多有些小變化。因為這套制度相對開放，不同解釋者可以用相同的一般方法輕易得出不同結論，所以誰有解釋權限\ix{terms}{01@\zhlidxcat{司法審查權限之歸屬}!解釋權限}很重要；爭議一向在於法官是否或應否從屬於其他解釋者，即立法者、行政首長與社群全體\pg{632 頁}。

既然誰有解釋權已定於法官，如何解釋的問題取得新的重要性，Kramer 稱這「giving rise for the first time to an active focus on interpretive methodology」〔首次產生對解釋方法的積極關注〕，並說明了關於憲法應如何建構之新理論的突然出現。左右兩派的法官與學者忽然發現，憲法須依獨特而特別的規則解釋，而各方產生的規則，似乎必然直接導向該方所偏好的結果\pg{632 至 633 頁}。

\subsection{（二）自制在兩派理論中的位置\pg{633 頁，註 41、42}}

自制\ix{terms}{05@\zhlidxcat{解釋理論與自制}!司法自制}與尊重\ix{terms}{05@\zhlidxcat{解釋理論與自制}!尊重}起初仍是自由派與保守派憲法理論的一環。對自由派，\en{Lochner\ix{cases}{02@美國!01@聯邦最高法院!1905Lochner v. New York（1905）@Lochner v. New York（1905）}} 的教訓延續為有限度審查，但這項教訓現在在實質上限於經濟管制及其他在左派關切的民權、少數者權利與平等保護之外的法律；在後者領域，積極的司法審查\ix{terms}{01@\zhlidxcat{司法審查權限之歸屬}!司法審查}成為「constitutional duty」〔憲法上的義務〕（註 41，引 Dworkin\ix{persons}{dworkin@Dworkin, Ronald（德沃金）}, \en{Taking Rights Seriously} 131--50；Ely, \en{Democracy and Distrust}；Ball, 59 Iowa L. Rev. 1059 (1974)）。對保守派情形正好相反：如前所述，自制成為限制 Warren 與 Burger 早期法院之自由派積極作為的論據，同時他們批評這些法院，使「true」Constitution〔真正的憲法〕處於「exile」〔流放〕，放棄監督國會權力並保護保守派所重視之其他權利的責任（註 42，引 Ginsburg, \en{Delegation Running Riot}, Regulation, Winter 1995, at 83, 84；Barnett）。

\subsection{（三）審慎根據是自制僅餘的理由\pg{633 頁}}

Kramer\ix{persons}{kramer@Kramer, Larry D.} 說，在兩派，推動新解釋理論\ix{terms}{05@\zhlidxcat{解釋理論與自制}!解釋理論}發展的同一個變化，即接受終局的解釋權在司法，也削弱了持久的尊重\ix{terms}{05@\zhlidxcat{解釋理論與自制}!尊重}或自制\ix{terms}{05@\zhlidxcat{解釋理論與自制}!司法自制}學說的基礎。穩固的尊重學說的基礎，在於有人可以尊重，且有好理由。接受司法至上\ix{terms}{01@\zhlidxcat{司法審查權限之歸屬}!司法至上}，消除了原本 Jeffersonian\ix{terms}{01@\zhlidxcat{司法審查權限之歸屬}!共和派} 的尊重基礎（尊重並保存美國人民首要的解釋權威），也消除了其 Thayer\ix{persons}{thayer@Thayer, James Bradley} 式的後繼者（把該權威轉給被認為最能回應人民的政治部門）。

剩下的只有審慎\ix{terms}{05@\zhlidxcat{解釋理論與自制}!審慎}的根據：不要因行動過於頻繁或積極而引發政治反撲，以確保司法權威。Kramer 的問句是：「Why should the Court defer, after all, if it is, by constitutional design, supposed to be the final, authoritative interpreter?」〔法院若依憲法設計本應是終局的、權威的解釋者，究竟為何應該尊重？〕答案是除了保護法院扮演其制度角色的能力之外沒有理由。在司法至上的世界，尊重「simply a matter of expediency」〔僅是權宜問題〕，即法院在實務上可走多遠而不致耗掉太多政治資本。據此，左右兩派都要求法院尊重，但只限於政治過程的結果是他們贊成或不太在意的領域，Kramer 稱「hardly surprising」〔不足為奇〕\pg{633 頁}。

\subsection{（四）均衡點\pg{633 至 634 頁}}

司法至上\ix{terms}{01@\zhlidxcat{司法審查權限之歸屬}!司法至上}愈深入國家的政治文化，感到需要尊重\ix{terms}{05@\zhlidxcat{解釋理論與自制}!尊重}的程度就愈低\pg{633 至 634 頁}。Kramer\ix{persons}{kramer@Kramer, Larry D.} 說，任何容許司法審查\ix{terms}{01@\zhlidxcat{司法審查權限之歸屬}!司法審查}的制度，即使採司法至上，都有一個均衡點\ix{terms}{05@\zhlidxcat{解釋理論與自制}!均衡點}，法院超過它就會損害其制度權威與行動能力。人們對司法審查的反應也受人們對法院應有多少權威之感受的影響，不單取決於是否同意法院的結果：「People will defer to a superior with whom they disagree more than a subordinate」〔人們對上位者的不同意見\ix{terms}{04@\zhlidxcat{法院之功能}!不同意見與協同意見}，比對下屬更能容忍〕，他舉的例子是院長做的事教師會接受，職員做就不會\pg{634 頁}。因此，若人們相信終局決定權屬於法院，自己與代表都不握有該權，就會讓大法官走得更遠、做更多自己不喜歡的事。

Kramer 稱「Judicial supremacy is an ideology」〔司法至上是一種意識形態〕，其全部目的與效果，是把公眾與政治接受的均衡點移向有利於司法權威的一側。當司法至上在整個政治體系中被視為理所當然，法院不必顧慮權威即可行動的能力就自然增長，尊重所起的作用也愈來愈小，直到如 Posner\ix{persons}{posner@Posner, Richard A.} 所言，成為「little more than a rhetorical tool used opportunistically by pretty much all of the Justices」〔幾乎所有大法官都機會主義地使用的修辭工具而已〕\pg{634 頁}。

\subsection{（五）結論\pg{634 頁}}

全文最後一句，Kramer\ix{persons}{kramer@Kramer, Larry D.} 說若這是壞事，而他認為是，禍首在「our embrace of the strange and counterintuitive notion that determining these answers, and so deciding the ongoing meaning of a democratic Constitution, is best left to life-tenured judges whose chief qualification is technical legal training」〔我們接受了這個奇怪而違反直覺的想法：決定這些答案，因而決定民主憲法持續演變的意義，最宜交給終身任職、主要資格是技術性法律訓練的法官〕。同一句開頭已把解釋理論\ix{terms}{05@\zhlidxcat{解釋理論與自制}!解釋理論}，以及憲法問題有正確答案這個想法，排除在禍首之外。

% 自 notes/docs/1151/seminar/W09_1109_預習.md（87808f3）由 tools/md_to_tex.py 於 2026-11-09 轉出一次；此後本檔為正本。
\chapter{Ferejohn 與 Pasquino〈Constitutional Courts as Deliberative Institutions〉\pg{21 至 36 頁}}
\label{ch:FP}

\begin{longtable}[]{@{}
  >{\raggedright\arraybackslash}p{(\linewidth - 4\tabcolsep) * \real{0.1200}}
  >{\raggedright\arraybackslash}p{(\linewidth - 4\tabcolsep) * \real{0.7500}}
  >{\raggedright\arraybackslash}p{(\linewidth - 4\tabcolsep) * \real{0.1300}}@{}}
\caption{Ferejohn 與 Pasquino 一文的節次與頁次。}\tabularnewline
\toprule\noalign{}
\begin{minipage}[b]{\linewidth}\raggedright
節次
\end{minipage} & \begin{minipage}[b]{\linewidth}\raggedright
小標
\end{minipage} & \begin{minipage}[b]{\linewidth}\raggedright
頁次
\end{minipage} \\
\midrule\noalign{}
\endfirsthead
\toprule\noalign{}
\begin{minipage}[b]{\linewidth}\raggedright
節次
\end{minipage} & \begin{minipage}[b]{\linewidth}\raggedright
小標
\end{minipage} & \begin{minipage}[b]{\linewidth}\raggedright
頁次
\end{minipage} \\
\midrule\noalign{}
\endhead
\bottomrule\noalign{}
\endlastfoot
導論 & （無小標） & 21--22 \\
1 & Deliberative Expectations & 22--25 \\
2 & Democratic Deliberation & 25--27 \\
3 & Constitutional Courts in Modern Democracy（3.1 Constitutional Adjudication in Post-Authoritarian Regimes；3.2 A Special Case: the French Conseil Constitutionel） & 28--33 \\
4 & Deliberation in Constitutional Courts & 33--36 \\
\end{longtable}

\section{導論：Kelsen 型模式的擴散與合法性負擔\pg{21 至 22 頁}}
\label{sec:FP-1}

\subsection{（一）東歐、南歐與戰後歐洲的憲法法院\pg{21 頁，註 1}}

首句即寫「The collapse of the communist empire coincided with a new wave of constitution making」〔共產帝國的瓦解與新一波制憲同時發生〕。東歐各國的新憲法都設有負責憲法審判的機關，其模式是 Ferejohn\ix{persons}{ferejohn@Ferejohn, John} 與 Pasquino\ix{persons}{pasquino@Pasquino, Pasquale} 在後文界定的 Kelsenian\ix{terms}{06@\zhlidxcat{審查制度之類型}!Kelsen 型模式} type（Kelsen\ix{persons}{kelsen@Kelsen, Hans（凱爾森）} 型）。南歐在約二十年前也有相同的現象，起因是希臘、葡萄牙與西班牙的法西斯威權政體瓦解；第二次世界大戰之後，奧地利、義大利與德國走過類似的歷程。Ferejohn 與 Pasquino 據此認為，下文對後威權\ix{terms}{06@\zhlidxcat{審查制度之類型}!後威權政體}政體憲法審判的討論，與東歐的新發展相關\pg{21 頁}。

註 1 補充希臘的情形：希臘自 19 世紀已有司法審查\ix{terms}{01@\zhlidxcat{司法審查權限之歸屬}!司法審查}，1975 年軍事政權結束時制定的憲法，才在分散審查\ix{terms}{06@\zhlidxcat{審查制度之類型}!分散審查與集中審查}制度之上設置一個特別而專門的最高法庭（Special Supreme Tribunal，註 1 所載條次為第 100 條）。

\subsection{（二）模型的共同前提\pg{21 頁，註 2}}

兩人稱這場制度變革可以用「a simple and more or less unified model」〔一個簡單而大致統一的模型〕理解\pg{21 頁}。其前提有二：歐洲各國都致力維持議會對行政與司法部門的權威；就社會學上的事實而言，各後威權\ix{terms}{06@\zhlidxcat{審查制度之類型}!後威權政體}國家對司法（後法西斯國家）、對議會（法國）或對兩者都有不信任。各國因此選擇由設在司法體系之外的專門法院，審查議會的行為。

註 2 建議把兩件事分開看：議會主權\ix{terms}{06@\zhlidxcat{審查制度之類型}!議會主權體制}理論及其階層式的權力分立\ix{terms}{06@\zhlidxcat{審查制度之類型}!權力分立體制}學說，本來就敵視任何司法審查\ix{terms}{01@\zhlidxcat{司法審查權限之歸屬}!司法審查}；制憲者對在威權政體中受教育的法官的不信任，以及對曾與法西斯或共產主義共謀的議會的不信任，則屬於社會學上的期待。兩者合起來，才說明 Kelsen 型\ix{terms}{06@\zhlidxcat{審查制度之類型}!Kelsen 型模式}模式為何成功。

\subsection{（三）合法性負擔與理由說明\pg{21 至 22 頁，註 3 至註 5}}

Ferejohn\ix{persons}{ferejohn@Ferejohn, John} 與 Pasquino\ix{persons}{pasquino@Pasquino, Pasquale} 指出，這個模式把部分立法性的政策形成交到憲法法院手中，「and not merely negative legislative authority as Kelsen\ix{persons}{kelsen@Kelsen, Hans（凱爾森）} himself was forced to admit」〔而且不只是 Kelsen 自己不得不承認的消極立法\ix{terms}{04@\zhlidxcat{法院之功能}!消極立法}權〕\pg{21 頁，註 3 指向 Kelsen 1931 年的 \en{Wer soll der Hüter der Verfassung sein?}，引義大利文譯本 \en{La giustizia costituzionale}, Milano: Giuffré 1981, p. 260}，這使法院承擔額外的 burden of legitimation\ix{terms}{07@\zhlidxcat{憲法審判之審議}!合法性負擔}〔合法性負擔〕。負擔的一部分可由兩種方式處理：使法官隔絕於政治壓力，並准許法院自訂能產生公正或公正外觀的程序，例如祕密開會、以法院名義發布意見\pg{21 至 22 頁}。但除此之外，憲法法院需要為裁判提出有理由的說明；Ferejohn 與 Pasquino 主張憲法法院面對特殊而嚴格的審議期待（deliberative expectations\ix{terms}{07@\zhlidxcat{憲法審判之審議}!審議的期待}）。

這項期待有一部分已寫入創設法院的規範：法國第五共和規範憲法委員會程序的組織法，要求憲法委員會就其裁判提出理由（註 4 引 1958 年 11 月 7 日第 58-1067 號條例第 20 條\ix{statutes}{03@法國!01@憲法委員會組織法（1958 年第 58-1067 號條例）!020@第 20 條}；義大利部分引 1953 年 3 月 11 日第 87 號法第 18 條\ix{statutes}{04@義大利!01@1953 年第 87 號法!018@第 18 條}，判決須「以義大利人民之名」作成並載明理由）。Rawls\ix{persons}{rawls@Rawls, John（羅爾斯）} 在 \en{Political Liberalism}\footcite[231--236]{rawls1993} 中把法院描述為審議機構的典範，即對於國家的強制性政策，期待並提供 reasons, explanations, and justifications〔理由、說明與正當化〕的論壇\pg{22 頁，註 5 指向該書 231 至 236 頁 ``The Supreme Court as Exemplar of Public Reason''}。

\subsection{（四）本章的比較對象\pg{22 頁}}

本章比較歐洲與美國的憲法法院作為審議論壇，對象是法國、德國、義大利、西班牙與美國的法院。兩人主張，這些法院在各政治體系中的位置很不相同：被不同的政治行動者問不同種類的問題，面對不同的期待、歷史與文化及政治上的限制，所以即使處理相當類似的議題，也會採取不同的審議實踐。Ferejohn\ix{persons}{ferejohn@Ferejohn, John} 與 Pasquino\ix{persons}{pasquino@Pasquino, Pasquale} 同時認為，這五個法院都保有 Rawls\ix{persons}{rawls@Rawls, John（羅爾斯）} 所描述的典範性審議特徵\pg{22 頁}。

\section{審議的期待\pg{22 至 25 頁}}
\label{sec:FP-2}

\subsection{（一）審議的界定\pg{22 至 23 頁}}

Ferejohn\ix{persons}{ferejohn@Ferejohn, John} 與 Pasquino\ix{persons}{pasquino@Pasquino, Pasquale} 從 Aristotle\ix{persons}{aristotle@Aristotle（亞里斯多德）} 的審議概念起論：該概念的重點在於審議的目的是批判性地評估、或許改變目標與偏好\pg{22 至 23 頁}。兩人自己的界定放寬了這項要求：「Rather than asking, like Aristotle, whether minds are changed, we would ask only whether public action is conducted through reasoning」〔我們不像 Aristotle 那樣問心意有沒有改變，只問公共行動是否透過說理而進行〕\pg{23 頁}。審議即參與關於公共行動的說理過程，包括對理由保持開放，也就是在有說服力的理由出現時願意改變偏好、信念或行動，並願意提出理由說服他人，排除威脅強制或欺瞞。

Ferejohn 與 Pasquino 承認美國最高法院的成員很少改變心意，至少在深層的政治承諾上如此。但法院的裁判作用於真實或假設的案件，目的在決定具體情境中可以或應該做什麼，所以對基本原則有歧見的法官，可以在較具體的層次上取得一致：「While abstract ideas may not be open to revision, ideas about particular cases may be much more flexible and pragmatic」〔抽象的觀念或許無從修正，關於個案的想法卻可以靈活得多、務實得多〕\pg{23 頁}。

\subsection{（二）理由說明的兩層意義\pg{23 頁，註 6}}

Ferejohn\ix{persons}{ferejohn@Ferejohn, John} 與 Pasquino\ix{persons}{pasquino@Pasquino, Pasquale} 指出，高等法院不同於其他政治機構，不只公布命令或裁判，也被期待公布其主文的 plausible rationales〔可信的理據〕，即他人不論自己的利益是否得到伸張，都可以期待加以尊重\ix{terms}{05@\zhlidxcat{解釋理論與自制}!尊重}與接受的論證\pg{23 頁}。這項期待在兩個意義上屬於審議的期待\ix{terms}{07@\zhlidxcat{憲法審判之審議}!審議的期待}。

\begin{enumerate}
\def\labelenumi{\arabic{enumi}.}
\tightlist
\item
  法院所給的理由可以被理解並接受為在某種規範意義上「我們自己的」行動理由；用 Rawls\ix{persons}{rawls@Rawls, John（羅爾斯）} 的語彙說，法院提供 public reasons〔公共理由\ix{terms}{07@\zhlidxcat{憲法審判之審議}!公共理由}〕。
\item
  法院是合議制機構，這些理由經由受該法院的裁判規範所引導的內部審議\ix{terms}{07@\zhlidxcat{憲法審判之審議}!內部審議}過程取得。這個過程是否受法院將公布單一意見或同時公布多數意見的共同期待規範，視情形而定。法院所公布的理由，在這個意義上是在一個「ranging from consensus seeking to majoritarian」〔從追求共識到多數決〕的規範框架內協商出來的。
\end{enumerate}

註 6 以美國最高法院愈來愈常公布不同意見\ix{terms}{04@\zhlidxcat{法院之功能}!不同意見與協同意見}書與協同意見書為例，說明這些裁判規範具有歷史上的彈性；註中另提出一種可能的解釋：在日益多元的文化中，對主文提出多種理由說明\ix{terms}{07@\zhlidxcat{憲法審判之審議}!理由說明}或有好的理由。

\subsection{（三）理由說明對法院權威的兩項功能\pg{24 至 25 頁}}

兩人認為理由說明\ix{terms}{07@\zhlidxcat{憲法審判之審議}!理由說明}對法院權威特別重要，分兩方面\pg{24 頁}。第一方面稱為民主的：由於（多數）法官不受選舉問責，裁判中所提出的理由，特別是其中一些理由植根於選舉問責機關的實定行為時，可以為公共行動提供間接的民主正當化。即使司法審議無法追溯到立法行為，只要植根於民主政府所由以建立、或民主的人民所預設的憲法或道德原則，司法理由仍可視為間接或遞移地民主。Ferejohn\ix{persons}{ferejohn@Ferejohn, John} 與 Pasquino\ix{persons}{pasquino@Pasquino, Pasquale} 因此認為，審議的法官可能增強民選官員的權力，替他們提供更有彈性且更聰明的制度，以執行並精練自己的立法；若這是立法者對法官的期待，民主的立法者就應該樂於擴張司法權。Ferejohn 與 Pasquino 另提出相反的可能：依憲法或道德原則行事的審議法官，可能增強人民（作為集體或個人）相對於其民主代表的權威。Rawls\ix{persons}{rawls@Rawls, John（羅爾斯）} 的公共理由\ix{terms}{07@\zhlidxcat{憲法審判之審議}!公共理由}說，在這個意義上強調了法官的民主性格，因為公共理由是不論各人對善的整全觀點為何，都可期待其認同的理由。司法審議對立法者也有風險：法官依理由行事，可以修改或廢棄民主的命令，並要求立法模式具有某種一致性或合理性，禁止民選官員制定某些政策，不論這些政策在政治上多麼有吸引力。這些風險是贊成民主或反民主，兩人認為取決於立法機關與政治共同體之間的連結等因素。

第二方面是司法所提供的理由使國家官員、其他法官、律師與一般公民，得以預期本次裁判對日後案件的意涵\pg{24 頁}。這項功能 ``efficiency enhancing''〔提升效率〕，發揮預測與協調的作用，不具正當化的作用；它使法院得以引導或協調大量私人與公共行為，而不致有太多徒勞的訴訟\pg{24 頁}。Ferejohn 與 Pasquino 接著說，效率也可以構成正當化，但「this kind of justification has nothing to do with democracy and would be available to courts in any legal system」〔這種正當化與民主無關，在任何法律體系中法院都可以援用〕\pg{24 至 25 頁}。

審議與說理是任何法律體系中法官的合法性來源，但在民主政體中分量特別重：政治權威須被視為直接或間接來自人民，民主機構須回應民意，這兩項期待使法官處於「a peculiarly weak position」〔特別脆弱的地位〕，這也是對司法審查\ix{terms}{01@\zhlidxcat{司法審查權限之歸屬}!司法審查}的 countermajoritarian\ix{terms}{05@\zhlidxcat{解釋理論與自制}!反多數決困難}〔反多數決〕面向感到憂慮的根源\pg{25 頁}。

\section{民主制度中的審議\pg{25 至 27 頁}}
\label{sec:FP-3}

\subsection{（一）行政、立法與選舉所受的審議期待\pg{25 至 26 頁，註 7}}

Ferejohn\ix{persons}{ferejohn@Ferejohn, John} 與 Pasquino\ix{persons}{pasquino@Pasquino, Pasquale} 拿法院所受的期待與其他機構對照：「In contrast to what we expect from courts, we do not demand as much by way of reason-giving from our other, more evidently democratic, institutions」〔相對於我們對法院的期待，我們對其他更明顯屬於民主的機構，沒有要求那麼多的理由說明\ix{terms}{07@\zhlidxcat{憲法審判之審議}!理由說明}〕\pg{25 頁}。三類機構的情形如下表。

\begin{longtable}[]{@{}
  >{\raggedright\arraybackslash}p{(\linewidth - 4\tabcolsep) * \real{0.1000}}
  >{\raggedright\arraybackslash}p{(\linewidth - 4\tabcolsep) * \real{0.8300}}
  >{\raggedright\arraybackslash}p{(\linewidth - 4\tabcolsep) * \real{0.0700}}@{}}
\caption{行政、立法與選舉所受的審議期待。}\tabularnewline
\toprule\noalign{}
\begin{minipage}[b]{\linewidth}\raggedright
機構
\end{minipage} & \begin{minipage}[b]{\linewidth}\raggedright
兩人所述的審議期待
\end{minipage} & \begin{minipage}[b]{\linewidth}\raggedright
頁次
\end{minipage} \\
\midrule\noalign{}
\endhead
\bottomrule\noalign{}
\endlastfoot
行政機關 & 可期待未經選舉的官僚在其決定與立法機關的決定之間提供某種理性連結；在美國及多數西方民主國家，不要求行政行為的理由既清楚又有說服力，只要所用手段與法律所定目的大致成比例；對個別主張作判斷時，正當程序的考量使行政機關承受類似法院的審議期待，制定一般或抽象政策時，這項期待改由選舉程序承擔（註 7） & 25 \\
立法機關 & 更少要求立法機關在任何法律論壇中說明理由；政黨綱領遠不足以為具體立法提案提供嚴謹的理由；法律文本的前言很少被認真對待，也幾乎不曾實質改變法律本身的理解；容許立法者互相護航（logroll）、拼湊式立法（checkerboard statutes）、聖誕樹式法案（Christmas tree bills），法律上只限制終局的立法行為與憲法價值或規範相衝突的情形 & 25--26 \\
選舉 & 公共規範近乎禁止說理，例如祕密投票、禁止在投票所附近競選、一人一票的民主觀，以及對競選捐款與支出的限制；法院不願追查選舉結果（含公民創制與公民投票）背後不被允許的意圖 & 26 \\
\end{longtable}

Ferejohn 與 Pasquino 把投票所稱為「a reason free zone\ix{terms}{07@\zhlidxcat{憲法審判之審議}!無須理由的領域}」〔無須理由的領域〕，行使投票權不以選民能提出或擁有好的理由為條件；理由若要反映在選舉程序中，必須透過個別選民，說服她在兩組考量之間擇一投票。Ferejohn 與 Pasquino 的結論是：「it is numbers and not reasons that count」〔計算的只有票數，理由不在其內〕\pg{26 頁}。

\subsection{（二）對公共審議可被腐化的疑慮\pg{26 至 27 頁}}

上述限制的理由，兩人認為相當明顯：自 Plato\ix{persons}{plato@Plato（柏拉圖）} 以來的理論家，擔心公共審議可被腐化，普通公民太容易被誘人的修辭或賄賂左右，或被公職人員與私人利益所脅迫\pg{26 頁}。歷史上這些疑慮引起對政黨、利益團體等可能壓倒普通公民的組織的深刻懷疑。Rousseau\ix{persons}{rousseau@Rousseau, Jean-Jacques（盧梭）} 與 Madison\ix{persons}{madison@Madison, James（麥迪遜）} 是兩個顯著的例子：Rousseau 與 Madison 都認為公共領域可被腐化，因而竭力尋求能引導這些強大傾向的制度改革。Rousseau 在《社會契約論》中主張，若可能，應禁止政黨與其他派系的形成，並禁止立法提案以私人個人與團體為對象；他雖然惋惜羅馬人必須採用祕密投票，仍承認共和國的腐化已太深，不容公開行使投票權\pg{26 至 27 頁}。Madison 對直接限制或規制私人利益持較懷疑的態度，提議以增加利益的數量來限制其在公共論壇的影響力。Ferejohn\ix{persons}{ferejohn@Ferejohn, John} 與 Pasquino\ix{persons}{pasquino@Pasquino, Pasquale} 記載，美國人又過了一百年才採用祕密投票，以限制選舉權的腐化\pg{27 頁}。

\subsection{（三）期待落差的三種解釋\pg{27 頁}}

Ferejohn\ix{persons}{ferejohn@Ferejohn, John} 與 Pasquino\ix{persons}{pasquino@Pasquino, Pasquale} 概述民主制度中各機構所受的審議期待：最民主、即最貼近人民的機構，完全不被期待審議，反而被各種妨礙審議的條件包圍；最不民主的決策機構，則被期待成為近乎純粹的公共理由\ix{terms}{07@\zhlidxcat{憲法審判之審議}!公共理由}論壇\pg{27 頁}。兩人提出三種可能的解釋。

\begin{enumerate}
\def\labelenumi{\arabic{enumi}.}
\tightlist
\item
  審議或說理可視為較直接的民主正當性的替代品：有票數就不必說服他人跟隨，因為他們選你就已承諾跟隨；沒有票數，說服與說理就是唯一的辦法。審議與說理因此是非民主機構取得人們服從的方式，是法院與行政機關為作成權威性決定而必須承擔的特別負擔。
\item
  西方民主傳統深處有一種懷疑，認為審議要求在某種程度上反民主或菁英主義。告訴人民或立法機關只有多數不夠，還必須有好的理由才能行動，等於弱化民主生活的一項要素。這個立場未必反對審議，但對在多數行動上加諸審議要求持懷疑態度；理由若是好的，自會影響多數並反映於票數。Ferejohn 與 Pasquino 寫道：「Reasons, in a democracy, should not get counted twice」〔在民主之下，理由不應被計算兩次〕。
\item
  在憲法民主之下，公共政策的選擇發生在規範的位階之中：道德與憲法規範適當地限制民主的選擇。規範這些規範與民主選擇之間的界線時，說理與審議特別重要，因為所設想的是告訴人民或其代表，某項政策不被允許。法院與行政機關（當其行為如法院時），因此在憲法價值可能涉及的領域承擔審議期待；立法機關在作為憲法解釋者時，會面臨類似的期待，但在制定明確處於憲法可接受範圍之內的政策時則不然。
\end{enumerate}

\section{憲法法院在現代民主中的位置\pg{28 至 33 頁}}
\label{sec:FP-4}

\subsection{（一）議會主權體制的審查\pg{28 頁，註 8、9}}

Ferejohn\ix{persons}{ferejohn@Ferejohn, John} 與 Pasquino\ix{persons}{pasquino@Pasquino, Pasquale} 說：「Roughly speaking, we may describe the institutional position of constitutional courts as situated along a single dimension」〔大致而言，可以把憲法法院的制度位置描述為落在單一向度上〕，左端是純粹的議會主權\ix{terms}{06@\zhlidxcat{審查制度之類型}!議會主權體制}體制，例如英國與法國第三、第四共和，右端是 Montesquieu\ix{persons}{montesquieu@Montesquieu（孟德斯鳩）} 式的權力分立\ix{terms}{06@\zhlidxcat{審查制度之類型}!權力分立體制}體制，中間是混合立法優越與其他較獨立機構的各種憲法體制\pg{28 頁}。議會主權體制依定義把行政與法院置於立法機關之下，其角色在執行並實施立法者的命令；司法審查\ix{terms}{01@\zhlidxcat{司法審查權限之歸屬}!司法審查}的觀念因此與這類法制完全格格不入。

圖~\ref{fig:FP-placement} 依第 3 節並列三種體制中各機關的配置、憲法審查權之所在與審查介入之時點。

\begin{qafigure}[三種體制中憲法審查權之所在]{三種體制中憲法審查權之所在。Ferejohn 與 Pasquino，28 至 33 頁。\label{fig:FP-placement}}
% Ferejohn 與 Pasquino 第 3 節（28 至 32 頁）：三種體制中憲法審查權之所在。
% 欄序依兩人所說的單一向度，左端議會主權、右端權力分立（28 頁）；Kelsen 型居中係依 30 頁「departed from the
% hierarchical parliamentary model」之描述排列。垂直、水平二詞出自 29 頁「separated horizontally rather than vertically」。
% 色調：審查權所在之機關 accent（terracotta），其餘機關 slate；卡片淡底加左側色條（skill diagram-design「動手前」第 2 條）。
% 各格文字逐一回刊本核過頁次（2026-11-09）。
\definecolor{figRev}{HTML}{B08060}\definecolor{figRevbg}{HTML}{F5ECEB}
\definecolor{figLine}{HTML}{E8D3D1}
% 卡片：\placecard{色}{左}{上}{右}{下}{內容}
\providecommand{\placecard}[6]{%
  \fill[#1bg] (#2,#3) rectangle (#4,#5);
  \fill[#1] (#2,#3) rectangle ($(#2,#5)+(0.6,0)$);
  \node[t] at ($(#2,#3)+(2.2,-1.4)$) {#6};}
\begin{tikzpicture}[qa,x=1mm,y=1mm,
    t/.style={anchor=north west,inner sep=0pt,align=left,execute at begin node={\raggedright\setstretch{1.2}}},
    cell/.style={t,font=\scriptsize,text width=40mm},
    rowlbl/.style={t,font=\scriptsize,text=dimink,text width=17mm},
    link/.style={figGray,line width=0.5pt},
    arr2/.style={figGray,line width=0.6pt,-{Stealth[length=4pt,width=3.2pt]}}]
  % 單一向度
  \draw[figGray,line width=0.5pt,{Stealth[length=3.5pt,width=3pt]}-{Stealth[length=3.5pt,width=3pt]}] (19,6) -- (146,6);
  \node[font=\scriptsize,text=dimink,anchor=south] at (82.5,6.6) {制度位置的單一向度（28 頁）};
  % 欄頭
  \foreach \x/\h/\s in {
      19/議會主權體制/{垂直：行政與法院\\從屬於立法機關（28 頁）},
      62.5/{Kelsen 型模式}/{憲法法院置於司法與\\其他部門之外（30 頁）},
      106/權力分立體制/{水平：三權分立（29 頁）}}{
    \node[t,font=\footnotesize] at (\x,2.5) {\h};
    \node[t,font=\scriptsize,text=dimink] at (\x,-1.8) {\s};
  }
  % 列名與列間淡線
  \node[rowlbl] at (0,-13) {機關之配置};
  \node[rowlbl] at (0,-50) {審查之時點\\與方式};
  \node[rowlbl] at (0,-61.5) {違憲之效果};
  \node[rowlbl] at (0,-73) {例};
  \foreach \y in {-10,-47.5,-59,-70.5}{\draw[figLine,line width=0.4pt] (0,\y) -- (146,\y);}
  % 議會主權：審查在立法機關內部，行政與法院居其下
  \begin{scope}[shift={(19,0)}]
    \placecard{figGray}{0}{-13}{40}{-34}{立法機關}
    \placecard{figRev}{3.5}{-19}{38}{-32}{憲法審查在其內部\\{\scriptsize 上議院、司法委員會、}\\{\scriptsize 法國憲法委員會（28 頁）}}
    \placecard{figGray}{0}{-38.5}{18}{-44.5}{行政}
    \placecard{figGray}{22}{-38.5}{40}{-44.5}{法院}
    \draw[link] (9,-34) -- (9,-38.5);
    \draw[link] (31,-34) -- (31,-38.5);
  \end{scope}
  % Kelsen 型：憲法法院居各部門之外、之上；普通法院停止訴訟並移送
  \begin{scope}[shift={(62.5,0)}]
    \placecard{figRev}{0}{-13}{40}{-29}{憲法法院\\{\scriptsize 審查權集中於此，}\\{\scriptsize 居各機關之上（31 頁）}}
    \placecard{figGray}{0}{-38.5}{10}{-44.5}{立法}
    \placecard{figGray}{12}{-38.5}{22}{-44.5}{行政}
    \placecard{figGray}{24}{-38.5}{40}{-44.5}{普通法院}
    \draw[arr2] (32,-38.5) -- (32,-29);
    \node[t,font=\scriptsize,text=dimink,anchor=north east,align=right] at (30.5,-30.6) {疑法律違憲者\\停止訴訟並移送};
  \end{scope}
  % 權力分立：三權並列，審查權散布於各級法院
  \begin{scope}[shift={(106,0)}]
    \placecard{figGray}{0}{-13}{10}{-19}{立法}
    \placecard{figGray}{12}{-13}{22}{-19}{行政}
    \placecard{figGray}{24}{-13}{40}{-19}{司法}
    \placecard{figRev}{24}{-25.5}{40}{-31.5}{最高法院}
    \placecard{figRev}{24}{-38.5}{40}{-44.5}{下級法院}
    \draw[link] (32,-19) -- (32,-25.5);
    \draw[link] (32,-31.5) -- (32,-38.5);
    \node[t,font=\scriptsize,text=dimink,anchor=north east,align=right] at (22.5,-27.6) {任何法院皆得審查，\\最高法院居其首};
  \end{scope}
  % 各列內容
  \foreach \x/\a/\b/\c in {
      19/{事前、抽象：比較法律\\與憲法之文本（28 頁）}/{法國：提案不得公布；\\公布後不受審查（32 頁）}/{英國、芬蘭、瑞典；法國\\第三至第五共和（28 頁）},
      62.5/{抽象審查之外，兼有\\事後之具體審查（31 頁）}/{廢止法律並重建法律狀態，\\兼為消極與積極立法（30 頁）}/{奧地利、德國、義大利、\\西班牙、葡萄牙（30 頁）；\\東歐（31 頁）},
      106/{於具體案件中為之，\\非抽象、非事前（29 頁）}/{法律不廢止，僅於該案\\拒絕適用（29 頁）}/{美國、挪威（29 頁，註 11）}}{
    \node[cell] at (\x,-50) {\a};
    \node[cell] at (\x,-61.5) {\b};
    \node[cell] at (\x,-73) {\c};
  }
\end{tikzpicture}

\end{qafigure}

兩人推論，議會主權體制必須把憲法審查設在立法機關之內，否則議會就不是主權者。因此憲法審查只能是 a priori\ix{terms}{06@\zhlidxcat{審查制度之類型}!抽象審查與具體審查}（先於立法提案公布為法律）而且是 abstract（僅涉及法律與憲法文本的比較）；援引憲法原則的權限，須集中於立法機關內部少數人手中。Ferejohn 與 Pasquino 舉英國、法國、芬蘭與瑞典（直到最近）為例，憲法審查全部發生在立法機關之內。傳統上這項審查可集中於上議院（可能有某種消極的立法權，如停止性否決權）、某種司法委員會，或獨立的機構如法國的國務院（Conseil d'Etat），後者可就憲法問題向立法機關提供意見。法國第五共和更設計出「a third legislative body」〔第三個立法機構〕，即 Conseil Constitutionel，可以修改或否決立法提案（註 8 引 Alec Stone, \en{The Birth of Judicial Politics in France}, Oxford 1992，與 John Bell, \en{French Constitutional Law}, Oxford 1992）。這個機構原由戴高樂派（Gaullists）創設，用來使議會合理化或加以控制，後來演變為允許對政府的立法提案作一般性憲法審查。無論審查或意見來自何處，把憲法原則適用於立法的權限仍在立法機關自身\pg{28 頁}。

註 9 記載這項轉變的兩個關鍵時刻：1971 年憲法委員會主張有權依廣泛而未成文的憲法原則（所謂 bloc de constitutionnalité\ix{terms}{06@\zhlidxcat{審查制度之類型}!憲法整體}〔憲法整體〕）廢棄政府的立法，在該案中從「共和國基本原則」辨識出結社自由的權利，並把 1789 年人權宣言與 1946 年憲法前言提升為可用以推翻立法提案的憲法規範；1974 年，政府擔心選舉失敗，成功促成一項修憲\ix{terms}{03@\zhlidxcat{憲法之效力基礎}!修憲}，使任何 60 名議員得將法律提交憲法委員會，此後憲法委員會憑 1971 年創造的強力憲法工具，愈來愈集中審查\ix{terms}{06@\zhlidxcat{審查制度之類型}!分散審查與集中審查}政府的提案\pg{28 至 29 頁}。

\subsection{（二）分散而具體的司法審查\pg{29 至 30 頁，註 10、11}}

權力分立\ix{terms}{06@\zhlidxcat{審查制度之類型}!權力分立體制}體制依定義把立法、行政、司法三權視為水平分立，並傾向把憲法權限賦予行使各項權力的機構\pg{29 頁}。Ferejohn\ix{persons}{ferejohn@Ferejohn, John} 與 Pasquino\ix{persons}{pasquino@Pasquino, Pasquale} 以 Montesquieu\ix{persons}{montesquieu@Montesquieu（孟德斯鳩）} 為例：他把無權力分立的狀態定義為專制，並認為法院適用法律於個案時，不應有任何重要意義上的立法作用，司法權因此是 pouvoir nul。兩人補充，這不表示法官（或更精確地說，陪審團）在適用法律時全無裁量；法院可以因適用某法律會產生不正義而拒絕在特定案件中適用，該法律仍然存在，其效力僅在該案中被排除。

這種解釋法律與事實以便適用的權力，產生一種特色鮮明的司法審查\ix{terms}{01@\zhlidxcat{司法審查權限之歸屬}!司法審查}，與議會主權\ix{terms}{06@\zhlidxcat{審查制度之類型}!議會主權體制}體制不同\pg{29 頁}。Ferejohn 與 Pasquino 列出三點：審查權散布於整個司法體系，因為任何法院為適用法律都須解釋憲法或制定法；審查在具體案件中行使，既非抽象也非 a priori\ix{terms}{06@\zhlidxcat{審查制度之類型}!抽象審查與具體審查}；審查不具立法性，法律並未被廢止，只是對該案（視法制而定，也及於類似案件）拒絕適用。Ferejohn 與 Pasquino 以美國憲法第三條為例，該條使最高法院居於聯邦司法部門之首，只准許憲法審查產生於真正的案件或爭議，並准許任何法院進行這類審查\pg{29 頁}。兩人寫道，在這方面美國的憲法審判制度與法國及多數歐洲模式「completely different」〔完全不同〕；法院適用制定法時必須常以憲法讀之，絕不以違反憲法保障的方式適用，所以制定法常被賦予使其配合憲法架構的解釋，偶爾被完全拒絕適用。這些作為也許有類似立法的效果，但不像法國制度那樣在形式上改變或消除制定法文本；違憲的制定法並未被廢止，只是不適用於特定爭議。Ferejohn 與 Pasquino 說明，這樣的說法凸顯司法審查的司法面向\pg{29 至 30 頁}。

註 10 指出，司法審查的非立法性質是理論上的觀念；實行分散審查\ix{terms}{06@\zhlidxcat{審查制度之類型}!分散審查與集中審查}的司法體系，通常會設計對下級法院的階層控制，從歧異的裁判中形成一致的規則或法理，這使司法行為更具立法性。註 11 提到挪威的憲法審判制度在許多方面與美國相似：審查發生於具體案件，且分散於普通法院\pg{29 頁}。

\subsection{（三）Kelsen 型模式的立法功能論\pg{30 至 31 頁，註 12、13}}

Ferejohn\ix{persons}{ferejohn@Ferejohn, John} 與 Pasquino\ix{persons}{pasquino@Pasquino, Pasquale} 寫道，歐洲多數法律體系有議會主權\ix{terms}{06@\zhlidxcat{審查制度之類型}!議會主權體制}的傳統，即行政與司法皆隸屬於議會，但 20 世紀的發展，尤其兩次世界大戰之後，侵蝕了這些承諾；因此一次大戰後的奧地利，二次大戰後的德國與義大利（以及法西斯政權瓦解後的西班牙與葡萄牙），採用了在重要方面偏離階層式議會模型的憲法\pg{30 頁}。兩人為方便起見把這些體制統稱 Kelsenian\ix{terms}{06@\zhlidxcat{審查制度之類型}!Kelsen 型模式}，取名自發明其獨特制度形式的奧地利法學家。

Ferejohn 與 Pasquino 歸納 Kelsen\ix{persons}{kelsen@Kelsen, Hans（凱爾森）} 的核心觀念為：憲法審判與其說是司法功能，不如說是立法功能。憲法法院宣告制定法違憲，不只是在消極意義上立法而廢止一部法律，由於必須重建該制定法之前的法律狀態，也在積極意義上立法\pg{30 頁，註 13}。註 13 說明，憲法法院依憲法價值解釋制定法時更是如此；Kelsen 本人對憲法法院角色的界定曾經猶豫，起初稱為「negative legislation」〔消極立法\ix{terms}{04@\zhlidxcat{法院之功能}!消極立法}〕，後來回應 C. Schmitt\ix{persons}{schmitt@Schmitt, Carl（施密特）} 時，在 \en{Wer soll der Hüter der Verfassung sein?}（1931）中承認法院有積極的立法功能（Ferejohn 與 Pasquino 另引 P. Pasquino, ``Gardien de la constitution ou justice constitutionnelle? C. Schmitt et H. Kelsen'', in \en{1789 et l'invention de la constitution}, Paris: Bruylant, 1994, pp. 141--152）。註 12 記載奧地利於 1920 年成為歐洲大陸第一個設憲法法院的國家，Hans Kelsen 在其設立中扮演關鍵角色，其歷史前身是奧匈帝國與神聖日耳曼帝國（Reichskammergericht）中裁判中央政府與各邦（Länder）衝突的法院。

依兩人的描述，Kelsen 對憲法審判的理解著重其立法面向，把它設想為制定法與憲法文本的比較；這種 abstract review 產生於文本的 a priori\ix{terms}{06@\zhlidxcat{審查制度之類型}!抽象審查與具體審查} 抽象比較，與有真實（受害）當事人主張權利的具體個案無關\pg{30 頁}。這樣看待的憲法審判「seems inherently political」〔似乎本質上具有政治性〕，因為憲法法院必須在規範社會行為的各種替代規範規則之間審議並選擇。Kelsen 因此認為憲法法院應置於司法以及其他政府部門之外，其權力由經政治任命的法官行使，這些法官通常來自特別能作文本抽象比較、並有能力就規範審議且說明裁判的人，未必來自有司法經驗者\pg{30 至 31 頁}。

\subsection{（四）後威權政體的憲法審判\pg{31 至 32 頁，註 14 至 16}}

在 3.1 中，Ferejohn\ix{persons}{ferejohn@Ferejohn, John} 與 Pasquino\ix{persons}{pasquino@Pasquino, Pasquale} 寫道 Kelsen\ix{persons}{kelsen@Kelsen, Hans（凱爾森）} 的觀念對後威權\ix{terms}{06@\zhlidxcat{審查制度之類型}!後威權政體}政體特別有吸引力：奧地利、德國、義大利與西班牙都採行，1989 年之後也在東歐生根\pg{31 頁}。各後威權憲法所設的憲法審判機構，可審查立法行為，有時並審查行政與司法行為。但各國都正處於從舊的、被不信任的政權過渡到新政權的過程，法官因而被特別懷疑，不宜作為憲法審查權的持有者（註 14 說明這些國家的司法體系基本上是封閉的職業階層，特別不受外界影響，這種極端的隔離使它們在威權崩潰之後更不可能被賦予重要的額外權力）。所以有強有力的政治理由，把憲法審判置於司法之外：審查機構實際上位於其他各機構之上，得從憲法觀點審查任何政府行為，並交給經政治任命的專門機構。

Ferejohn 與 Pasquino 指出，Kelsen 強調抽象的憲法審查，但所有現代後威權憲法法院也都有具體的 a posteriori 審查權\pg{31 頁}。這些法院的管道，不只由政府與政治少數控制，也由個案中的普通訴訟當事人，或如義大利由普通法院控制。因此法院可能被要求以 Kelsen 式的方式抽象比較憲法與制定法文本（由直接提交憲法問題），或面對下級法院在進行中的案件所生的憲法問題，或面對整個已裁判的案件（德國與西班牙均有，註 15 說明對設於 Strasbourg 的歐洲人權法院的申訴，傳統上也如此，須待訴訟當事人用盡國內救濟後才審查）。在任何這些情形，廢棄或修改立法的實際權限，一般集中於憲法法院，不分散於整個司法體系。普通法院若懷疑法律合憲性，必須停止其面前的程序，並將問題提交憲法法院決定\pg{31 頁；註 16 說明適用歐洲共同體法的法院遵循相同的做法}。

憲法法院的組成也特別\pg{31 至 32 頁}。權力分立\ix{terms}{06@\zhlidxcat{審查制度之類型}!權力分立體制}體制傾向分散而具體的審查，普通法官可望發展出憲法審判的能力；在審查集中而且常為抽象的後威權體制，普通法官並無特殊的權威主張，加上威權的過去，法官至少在初期被懷疑能否擔任憲法與民主價值的仲裁者。因此在所有後威權體制中，法學教授往往與若干法官共同占據法院的多數席位\pg{32 頁}。

\subsection{（五）法國憲法委員會的特殊位置\pg{32 至 33 頁，註 17}}

3.2 把法國第五共和的憲法審查獨立討論。兩人說法國的共和傳統穩固屬於議會主權\ix{terms}{06@\zhlidxcat{審查制度之類型}!議會主權體制}傳統，自大革命以來即敵視憲法審判；但戴高樂及其支持者堅持對議會設置制度性的限制，憲法委員會即其一\pg{32 頁；註 17 記載 1795 年後雅各賓的國民公會一致否決了 Sieyès 提出的 jury constitutionnaire〔憲法陪審〕提案}。在議會主權體制內，這只能把法院實際上置於立法機關之內，因此比 Kelsen\ix{persons}{kelsen@Kelsen, Hans（凱爾森）} 所建議的更強調憲法委員會的立法功能。

Ferejohn\ix{persons}{ferejohn@Ferejohn, John} 與 Pasquino\ix{persons}{pasquino@Pasquino, Pasquale} 列出該委員會審查權的範圍與限制：

\begin{enumerate}
\def\labelenumi{\arabic{enumi}.}
\tightlist
\item
  只能在制定法公布之前審查，且只在政府或（自 1974 年起）議員或參議員的重要政治少數提請時。
\item
  立法提案若被憲法委員會宣告違憲，便不能成為法律。
\item
  制定法一旦成為法律，憲法委員會無權審查；普通法院也無權在一般訴訟中進行憲法審查。
\item
  依法國的見解，制定法是主權者，不受任何外部檢驗，更不受法官或憲法法院檢驗；憲法法院作為立法機關自身的一部分，在保存法律主權的觀念上扮演關鍵角色。
\end{enumerate}

Ferejohn 與 Pasquino 這樣概括：「Legislative action cannot occur in the presence of an objection by the Conseil, but once it takes place, is unchecked by constitutional mechanisms」〔有憲法委員會的異議，立法行為不能成立；一旦成立，就不受憲法機制檢驗〕\pg{32 頁}。

因為法國的立法提案在憲法委員會拒絕時不能成為法律，審查（若被請求，且幾乎總是由政治少數針對政府所提立法提出）必須在立法行為之後立刻進行，面對的是其提案受到挑戰的現任政府。兩人認為，這表示審查必須迅速，而且鑑於法國政治制度與政治文化的多數決性質，處在可能高度政治化的情況\pg{32 頁}。Ferejohn 與 Pasquino 對照美國：憲法對制定法的挑戰須等到真正提出該議題的案件，往往在制定該法的立法機關消失之後很久；在德國、西班牙與義大利等後威權\ix{terms}{06@\zhlidxcat{審查制度之類型}!後威權政體}體制，憲法法院不能阻止法律生效，即使審查具有抽象性質，政府正在執行受爭議之法律，並沒有太大的政治壓力要迅速解決憲法問題（Ferejohn 與 Pasquino 並說政治壓力或許反向作用），抽象審查\ix{terms}{06@\zhlidxcat{審查制度之類型}!抽象審查與具體審查}因此可在政治戰火稍歇之後進行\pg{32 至 33 頁}。

關於組成，兩人認為或許因為位在立法機關之內，憲法委員會的組成相當獨特：成員只是偶爾有實質的法律或司法背景，法學教授也少；憲法委員會由任命權者的政治盟友主導。Ferejohn 與 Pasquino 補充這並不表示成員無法理解或提出法律論證，而是該機構看重的是政治判斷，這在他們必須迅速作成決定的急迫性之下尤為關鍵\pg{33 頁}。

\section{憲法法院的審議\pg{33 至 36 頁}}
\label{sec:FP-5}

\subsection{（一）三種制度與三種審議型態\pg{33 頁，註 18}}

第 4 節論證三種不同的憲法審查制度，產生三種可以區別的審議型態\pg{33 頁}。

\begin{longtable}[]{@{}
  >{\raggedright\arraybackslash}p{(\linewidth - 6\tabcolsep) * \real{0.1400}}
  >{\raggedright\arraybackslash}p{(\linewidth - 6\tabcolsep) * \real{0.4400}}
  >{\raggedright\arraybackslash}p{(\linewidth - 6\tabcolsep) * \real{0.3400}}
  >{\raggedright\arraybackslash}p{(\linewidth - 6\tabcolsep) * \real{0.0800}}@{}}
\caption{三種審查制度與其審議型態。}\tabularnewline
\toprule\noalign{}
\begin{minipage}[b]{\linewidth}\raggedright
制度
\end{minipage} & \begin{minipage}[b]{\linewidth}\raggedright
審查的所在與方式
\end{minipage} & \begin{minipage}[b]{\linewidth}\raggedright
審議型態
\end{minipage} & \begin{minipage}[b]{\linewidth}\raggedright
頁次
\end{minipage} \\
\midrule\noalign{}
\endhead
\bottomrule\noalign{}
\endlastfoot
權力分立\ix{terms}{06@\zhlidxcat{審查制度之類型}!權力分立體制}體制 & 憲法審查在普通法院進行；法官之間需要高度協調與溝通 & 傾向發展出公開的審議形式：公開聽審、公布意見與表決、先例理論 & 33 \\
Kelsen 型\ix{terms}{06@\zhlidxcat{審查制度之類型}!Kelsen 型模式}模式 & 憲法審查權集中於專門法庭；普通法院須能適用其裁判，故裁判須明確 & 傾向以單一聲音發言並闡明理由；很少公開聽審或公布不同意見\ix{terms}{04@\zhlidxcat{法院之功能}!不同意見與協同意見}書 & 33 \\
議會主權\ix{terms}{06@\zhlidxcat{審查制度之類型}!議會主權體制}體制 & 憲法審查是立法程序的一部分，目的在確保立法合於憲法要求；如何進行，取決於審查如何納入立法程序 & 依審查納入立法程序的方式而定 & 33 \\
\end{longtable}

Ferejohn\ix{persons}{ferejohn@Ferejohn, John} 與 Pasquino\ix{persons}{pasquino@Pasquino, Pasquale} 就議會主權體制中審查制度最為制度化的法國說明：憲法委員會對列入其審查事項的立法提案，在祕密會議中審議，沒有公開聽審或陳述\pg{33 頁}。法國憲法審查的重要特徵是，經議會多數通過的制定法在憲法委員會同意之前不能公布，所以審查程序承受巨大的時間壓力：延遲通常為一個月，「this period shall be reduced to eight days where the Government declares the matter urgent」〔政府宣告案件緊急時，此期間縮短為八日〕（註 18 引 1958 年 11 月第 58-1067 號條例第 25 條\ix{statutes}{03@法國!01@憲法委員會組織法（1958 年第 58-1067 號條例）!025@第 25 條}）。

\subsection{（二）法國憲法委員會的意見與解釋保留\pg{33 至 35 頁，註 19 至 22}}

憲法委員會以法院的名義發布意見，沒有紀錄表決，也沒有公布的不同意見\ix{terms}{04@\zhlidxcat{法院之功能}!不同意見與協同意見}書；意見通常非常簡短，不詳述對立法的憲法異議\pg{33 至 34 頁}。委員會可以報告提案全部或一部在憲法上有疑義，所依據的憲法淵源有數個：1789 年人權宣言、1946 年憲法前言，以及「Fundamental Principles underlying the laws of the {[}Third{]} Republic」〔第三共和國法律所依據的基本原則〕\pg{34 頁}。兩人稱最具影響力的許多裁判，涉及對行政機關過寬或不受限制的授權（註 20 舉車輛搜索案件與國有化裁判；1971 年發現結社自由是憲法原則的裁判也屬此型）。面對這類疑慮，政府通常可以重新起草立法，納入若干規範日後行政行為的標準\pg{34 至 35 頁}。

註 19 以 1971 年 7 月 16 日的意見為代表，說明該意見規定結社權並大幅修改委員會的角色，而且很短，並未真正闡明把 bloc de constitutionnalité\ix{terms}{06@\zhlidxcat{審查制度之類型}!憲法整體} 列為憲法審判依據的理由；註中並錄 J. Bell 的英譯（\en{French Constitutional Law}, Oxford: Clarendon Press, 1992, pp. 272--73），其中的理由段落稱：「among the fundamental principles recognized by the laws of the Republic and solemnly reaffirmed by the Constitution, is to be found the freedom of association」〔在共和國法律所承認並為憲法所莊嚴重申的基本原則中，可以找到結社自由〕\pg{34 頁}。

Ferejohn\ix{persons}{ferejohn@Ferejohn, John} 與 Pasquino\ix{persons}{pasquino@Pasquino, Pasquale} 接著列舉三點\pg{35 頁}：

\begin{enumerate}
\def\labelenumi{\arabic{enumi}.}
\tightlist
\item
  法國憲法委員會宣告立法提案無效的頻率不低，但更常認定提案的一部分與憲法原則相悖（註 21 引 A. Stone 的資料：1981 至 1993 年，約 50\% 被提請審查的法律被委員會指為有違憲之處）。
\item
  委員會愈來愈常發布意見，指出法律只有在以特定方式解釋時才在憲法上有效（註 22 引 Alexandre Viala, \en{Les réserves d'interprétation\ix{terms}{07@\zhlidxcat{憲法審判之審議}!解釋保留} dans la jurisprudence du Conseil constitutionnel}, Paris: L.G.D.J., 1999）。
\item
  Ferejohn 與 Pasquino 認為最後這項策略特別成問題，因憲法委員會與司法體系分離：「How can the Conseil offer any assurance that a promulgated law will be interpreted in a constitutional fashion?」〔憲法委員會如何能保證已公布的法律將以合憲的方式被解釋？〕沒有 stare decisis 之類的正式機制或法理可確保此事，唯一的保證是委員會向普通法院、最高法院（Cour de Cassation）或行政機關所提理由的說服力。兩人引述的跡象是：行政機關為尋求來自一個可能製造麻煩的法庭的明確指引，傾向把委員會的裁判用於制定對屬員的內部指引；普通法院為尋求法律穩定，常把委員會的解釋當作權威接受\pg{35 頁}。
\end{enumerate}

\subsection{（三）Kelsen 型與後威權法院的內部審議\pg{35 頁}}

Ferejohn\ix{persons}{ferejohn@Ferejohn, John} 與 Pasquino\ix{persons}{pasquino@Pasquino, Pasquale} 概括 Kelsen 型\ix{terms}{06@\zhlidxcat{審查制度之類型}!Kelsen 型模式}或後威權\ix{terms}{06@\zhlidxcat{審查制度之類型}!後威權政體}體制的法院：通常祕密審議，很少舉行公開聽審，意見以全體法院的名義發布，沒有記載表決，少有制度允許公布不同意見\ix{terms}{04@\zhlidxcat{法院之功能}!不同意見與協同意見}書\pg{35 頁}。直接開放訴訟當事人聲請的法院（如西班牙與德國）通常設有內部分工，以小組處理這類案件；有些法院准許或鼓勵就標的事項專業化（義大利常見，德國則反對）。Ferejohn 與 Pasquino 稱所有這些法院都是「internally deliberative」〔內部審議\ix{terms}{07@\zhlidxcat{憲法審判之審議}!內部審議}的〕，意即其工作多以合議與面對面的方式進行，真實地嘗試說服他人並作成集體說理的裁判；義大利尤其如此，法院以共識的裁判為目標而且通常達成（兩人註明這是從與法官及其他觀察者的討論所能得知的範圍）。

\subsection{（四）美國最高法院的外部審議\pg{35 至 36 頁}}

Ferejohn\ix{persons}{ferejohn@Ferejohn, John} 與 Pasquino\ix{persons}{pasquino@Pasquino, Pasquale} 寫道，美國最高法院的審議實踐與上述情形相當不同，可稱為 externally deliberative〔外部審議\ix{terms}{07@\zhlidxcat{憲法審判之審議}!外部審議}〕，與內部審議\ix{terms}{07@\zhlidxcat{憲法審判之審議}!內部審議}相對\pg{35 頁}。一部分原因在於法院對憲法解釋並無壟斷，主要作用在規範其他法院適用憲法的過程；這種協調或規制的角色，迫使法院以公開而透明的方式工作，使其他行動者，即法官、行政機關與律師，能清楚並（多半）成功地預期法院對尚未提到面前的案件會如何裁判。制定法可能在適用於個案時引起憲法問題，這使法院有時准許或鼓勵下級法院實驗，而不尋求立即解決限於文本本身的制定法的合憲性\pg{35 至 36 頁}。Ferejohn 與 Pasquino 說，這項策略只存在於准許具體審查\ix{terms}{06@\zhlidxcat{審查制度之類型}!抽象審查與具體審查}的體制，Kelsen 型\ix{terms}{06@\zhlidxcat{審查制度之類型}!Kelsen 型模式}體制僅在處理憲法訴願時可能發生（兩人列舉 amparo 與 Verfassungsbeschwerde）。

Ferejohn 與 Pasquino 列舉最高法院公開行事的幾個意義\pg{36 頁}：

\begin{enumerate}
\def\labelenumi{\arabic{enumi}.}
\tightlist
\item
  向一般訴訟當事人開放的案源，並公開其是否受理案件的決定；偶爾有成員甚至公開宣布個人對此的看法。
\item
  公開聽審是常態（至少對法院案件中最重要的部分），表決有紀錄，並公布不同意見\ix{terms}{04@\zhlidxcat{法院之功能}!不同意見與協同意見}書與協同意見書。
\item
  法院很少試圖以一個聲音發言，似乎寧可讓衝突與意見不一發酵（原文印作 ``ferment''）。
\end{enumerate}

就內部審議而言，Ferejohn 與 Pasquino 說「On most accounts the actual level of internal deliberation on the court may be low」〔依多數說法，法院內部審議的實際程度可能偏低〕：大法官在評議中不花多少時間審議案件中的具體爭點，觀察者通常報告大法官與自己的助理互動最多，與其他大法官較少，也很少以真正改變心意或票數為目的。內部審議多半發生於評議時的初步表決之後，即大法官決定能否全部或部分同意為法院所撰意見時；這種審議大多以書面進行，並非面對面的互動。兩人認為正因如此，審議過程對局外人比面對面的審議透明得多：「The public manner of conducting its business puts the Court at the center of wider deliberative processes within the judiciary, the legal, and political communities」〔公開辦理事務的方式，使法院處於司法、法律與政治各社群中更廣的審議過程的中心〕\pg{36 頁}。

圖~\ref{fig:FP-deliberation} 依三種審查制度，並列審議在法院之內進行或對著法院之外的行動者進行。

\begin{qafigure}[審議的所在與對象]{審議的所在與對象。Ferejohn 與 Pasquino，33 至 36 頁。\label{fig:FP-deliberation}}
% Ferejohn 與 Pasquino 第 4 節（33 至 36 頁）：審議發生在法院之內或對著法院之外。
% Kelsen 型與後威權法院：祕密審議、意見以全體法院名義發布（35 頁）；法官面對面說服、產生集體說理的裁判（35 頁）；
%   裁判須明確，普通法院始能適用（33 頁）；很少公開聽審或公布不同意見書（33 頁）。
% 法國憲法委員會：祕密會議、無公開聽審（33 頁）；意見簡短、無表決紀錄（33 至 34 頁）；與司法體系分離，
%   無 stare decisis 之類機制，唯一的保證是向普通法院、Cour de Cassation、行政機關所提理由的說服力（35 頁）。
% 美國最高法院：公開聽審、記名表決、公布協同與不同意見書（36 頁）；內部審議多在初步表決後以書面進行（36 頁）；
%   公開透明使法官、行政機關、律師得預期法院之裁判（35 頁）；法院處於司法、法律、政治社群更廣的審議過程中心（36 頁）。
% 色調：外部審議之法院與其流向用 accent（terracotta），其餘 slate；卡片淡底加左側色條。
\definecolor{figRev}{HTML}{B08060}\definecolor{figRevbg}{HTML}{F5ECEB}
\definecolor{figLine}{HTML}{E8D3D1}
\providecommand{\placecard}[6]{%
  \fill[#1bg] (#2,#3) rectangle (#4,#5);
  \fill[#1] (#2,#3) rectangle ($(#2,#5)+(0.7,0)$);
  \node[t] at ($(#2,#3)+(2.4,-1.9)$) {#6};}
\begin{tikzpicture}[qa,x=1mm,y=1mm,
    t/.style={anchor=north west,inner sep=0pt,align=left,execute at begin node={\raggedright\setstretch{1.2}}},
    note/.style={font=\scriptsize,text=dimink,inner sep=0pt,align=left,anchor=west,execute at begin node={\setstretch{1.15}}},
    link/.style={figGray,line width=0.6pt,-{Stealth[length=4pt,width=3.2pt]}},
    weak/.style={figGray,line width=0.6pt,dash pattern=on 2pt off 1.6pt,-{Stealth[length=4pt,width=3.2pt]}},
    flow/.style={figRev,line width=1.1pt,-{Stealth[length=4.5pt,width=3.6pt]}},
    both/.style={figGray,line width=0.5pt,{Stealth[length=3pt,width=2.6pt]}-{Stealth[length=3pt,width=2.6pt]}}]
  \foreach \x in {47.5,98.5}{\draw[figLine,line width=0.4pt] (\x,5) -- (\x,-96);}
  % 欄頭
  \foreach \x/\h/\s in {
      0/{Kelsen 型與後威權法院}/{內部審議（35 頁）},
      51/法國憲法委員會/{議會主權體制（33 頁）},
      102/美國最高法院/{外部審議（35 頁）}}{
    \node[t,font=\footnotesize] at (\x,4) {\h};
    \node[t,font=\scriptsize,text=dimink] at (\x,-0.6) {\s};
  }
  % Kelsen 型：審議在法官之間
  \begin{scope}[shift={(0,0)}]
    \placecard{figGray}{0}{-12}{44}{-46}{憲法法院\\{\scriptsize 祕密審議（35 頁）}\\{\scriptsize 很少公開聽審（33 頁）}\\{\scriptsize 意見以全院名義發布（35 頁）}}
    \foreach \p in {10,22,34}{\fill[figGray] (\p,-35) circle (1.3);}
    \draw[both] (11.8,-35) -- (20.2,-35);
    \draw[both] (23.8,-35) -- (32.2,-35);
    \node[note,anchor=north west,text=caseink] at (2.4,-37.4) {面對面說服，集體說理（35 頁）};
    \draw[link] (8,-46) -- (8,-68);
    \node[note] at (11,-57) {裁判須明確，普通法院\\始能適用（33 頁）};
    \placecard{figGray}{0}{-68}{20}{-80}{{\scriptsize 普通法院}}
  \end{scope}
  % 法國：祕密審議，意見向外傳遞但無正式拘束機制
  \begin{scope}[shift={(51,0)}]
    \placecard{figGray}{0}{-12}{44}{-46}{憲法委員會\\{\scriptsize 祕密會議，無公開聽審（33 頁）}\\{\scriptsize 意見簡短（33 至 34 頁）}\\{\scriptsize 與司法體系分離（35 頁）}}
    \draw[figGray,line width=0.6pt,dash pattern=on 2pt off 1.6pt] (7,-46) -- (7,-59);
    \draw[figGray,line width=0.6pt,dash pattern=on 2pt off 1.6pt] (7,-59) -- (37,-59);
    \draw[weak] (7,-59) -- (7,-68);
    \draw[weak] (22,-59) -- (22,-68);
    \draw[weak] (37,-59) -- (37,-68);
    \node[note] at (10,-52.5) {無正式機制，唯憑理由\\的說服力（35 頁）};
    \placecard{figGray}{0}{-68}{14}{-80}{{\scriptsize 普通法院}}
    \placecard{figGray}{15}{-68}{29}{-80}{{\scriptsize Cour de}\\{\scriptsize Cassation}}
    \placecard{figGray}{30}{-68}{44}{-80}{{\scriptsize 行政機關}}
  \end{scope}
  % 美國：公開行事，對著法院之外的行動者與更廣的社群
  \begin{scope}[shift={(102,0)}]
    \placecard{figRev}{0}{-12}{44}{-46}{最高法院\\{\scriptsize 公開聽審，記名表決（36 頁）}\\{\scriptsize 公布協同與不同意見書（36 頁）}\\{\scriptsize 內部審議多在初步表決後}\\{\scriptsize 以書面進行（36 頁）}}
    \draw[flow] (7,-46) -- (7,-59);
    \draw[figRev,line width=1.1pt] (7,-59) -- (37,-59);
    \draw[flow] (7,-59) -- (7,-68);
    \draw[flow] (22,-59) -- (22,-68);
    \draw[flow] (37,-59) -- (37,-68);
    \node[note] at (10,-52.5) {公開行事，使其得預期\\裁判（35 頁）};
    \placecard{figGray}{0}{-68}{14}{-80}{{\scriptsize 下級法院}}
    \placecard{figGray}{15}{-68}{29}{-80}{{\scriptsize 行政機關}}
    \placecard{figGray}{30}{-68}{44}{-80}{{\scriptsize 律師}}
    \placecard{figGray}{0}{-84}{44}{-96}{{\scriptsize 法院處於司法、法律、政治社群}\\{\scriptsize 更廣的審議過程之中（36 頁）}}
  \end{scope}
\end{tikzpicture}

\end{qafigure}

% 自 notes/docs/1151/seminar/W09_1109_預習.md（87808f3）由 tools/md_to_tex.py 於 2026-11-09 轉出一次；此後本檔為正本。
\chapter{與前後幾週的關係及課前問題}
\label{ch:relations}

\section{與 W04 至 W08、W10 的對照}
\label{sec:relations-1}

\begin{longtable}[]{@{}
  >{\raggedright\arraybackslash}p{(\linewidth - 4\tabcolsep) * \real{0.1300}}
  >{\raggedright\arraybackslash}p{(\linewidth - 4\tabcolsep) * \real{0.2400}}
  >{\raggedright\arraybackslash}p{(\linewidth - 4\tabcolsep) * \real{0.6300}}@{}}
\caption{本週三篇與 W04 至 W08、W10 指定文獻的對照。}\tabularnewline
\toprule\noalign{}
\begin{minipage}[b]{\linewidth}\raggedright
文獻
\end{minipage} & \begin{minipage}[b]{\linewidth}\raggedright
相關論點
\end{minipage} & \begin{minipage}[b]{\linewidth}\raggedright
與本週三篇的關係
\end{minipage} \\
\midrule\noalign{}
\endhead
\bottomrule\noalign{}
\endlastfoot
W04 Fallon、Bickel\ix{persons}{bickel@Bickel, Alexander M.} & 反多數決困難；法院在憲法意義不確定時的裁判權威 & Kramer\ix{persons}{kramer@Kramer, Larry D.} 631 頁說 1960 年代的自由派學者覺得「the countermajoritarian\ix{terms}{05@\zhlidxcat{解釋理論與自制}!反多數決困難} difficulty」並不那麼困難；Ferejohn\ix{persons}{ferejohn@Ferejohn, John} 與 Pasquino\ix{persons}{pasquino@Pasquino, Pasquale} 25 頁提到司法審查\ix{terms}{01@\zhlidxcat{司法審查權限之歸屬}!司法審查}的反多數決面向；兩篇都未引 Bickel \\
W05 Kelsen\ix{persons}{kelsen@Kelsen, Hans（凱爾森）}、Schmitt\ix{persons}{schmitt@Schmitt, Carl（施密特）} & 憲法的守護者應為何者；消極立法\ix{terms}{04@\zhlidxcat{法院之功能}!消極立法}者 & Ferejohn 與 Pasquino 21 頁註 3 引 Kelsen 1931 年〈Wer soll der Hüter der Verfassung sein?〉的義大利文譯本，30 頁註 13 說 Kelsen 回應 Schmitt 時由消極立法轉為承認憲法法院的積極立法功能；Alexander\ix{persons}{alexander@Alexander, Larry} 與 Schauer\ix{persons}{schauer@Schauer, Frederick} 460 至 461 頁註 17 引 Kelsen 的基本規範\ix{terms}{03@\zhlidxcat{憲法之效力基礎}!基本規範}，說明憲法的地位須在文本之外決定 \\
W06 Waldron\ix{persons}{waldron@Waldron, Jeremy} & 反對違憲審查的核心論證；1354 頁註 22 界定司法至上\ix{terms}{01@\zhlidxcat{司法審查權限之歸屬}!司法至上} & 三篇都未引 Waldron；Waldron 1350 頁註 11\footcite[1350]{waldron2006} 把 Kramer 的著作與 Tushnet\ix{persons}{tushnet@Tushnet, Mark} 並列為近年攻擊違憲審查的著作；W06 講義稱 Kramer 接受違憲審查而反對司法至上，與 Kramer 622 至 623 頁共和派\ix{terms}{01@\zhlidxcat{司法審查權限之歸屬}!共和派}的人民憲政主義\ix{terms}{01@\zhlidxcat{司法審查權限之歸屬}!人民憲政主義}相合 \\
W08 Tushnet & 非司法審查的表現；定紛止爭\ix{terms}{02@\zhlidxcat{定紛止爭論}!定紛止爭}論的經驗前提 & Tushnet 結論回應 Alexander 與 Schauer 的 Reply\pg{註 209 引 464 頁}；Alexander 與 Schauer 458 頁註 12、472 頁註 48 回應 Tushnet；Kramer 629 頁註 30 引 Tushnet 支持立法機關也能認真審議憲法 \\
W10 Law & 司法權與司法審查的理論 & 待下週預習時補 \\
\end{longtable}

\section{課前可以先想的問題}
\label{sec:relations-2}

\begin{enumerate}
\def\labelenumi{\arabic{enumi}.}
\tightlist
\item
  Alexander\ix{persons}{alexander@Alexander, Larry} 與 Schauer\ix{persons}{schauer@Schauer, Frederick} 457 頁註 11 把權威界定為提供與內容無關\ix{terms}{02@\zhlidxcat{定紛止爭論}!與內容無關之理由}之理由，其成立取決於受拘束者長期以定紛\ix{terms}{02@\zhlidxcat{定紛止爭論}!定紛止爭}為準時所犯的錯誤是否較少。Kramer\ix{persons}{kramer@Kramer, Larry D.} 634 頁稱司法至上\ix{terms}{01@\zhlidxcat{司法審查權限之歸屬}!司法至上}是一種意識形態，作用在把公眾接受的均衡點\ix{terms}{05@\zhlidxcat{解釋理論與自制}!均衡點}移向司法權威。兩篇談的都是其他機關與公眾對最高法院的接受，前者從規範上證成，後者從歷史上描述。Kramer 的敘述若成立，是否動搖 Alexander 與 Schauer 的論證？
\item
  Alexander 與 Schauer 457 頁說，最高法院在 Cooper\ix{cases}{02@美國!01@聯邦最高法院!1958Cooper v. Aaron（1958）@Cooper v. Aaron（1958）} 與 City of Boerne\ix{cases}{02@美國!01@聯邦最高法院!1997City of Boerne v. Flores（1997）@City of Boerne v. Flores（1997）} 宣告自身的至上地位，預設了其他部門應接受該宣告的結論，屬於循環論證。釋字第 185 號\ix{cases}{01@中華民國!01@司法院大法官!0185@釋字第 185 號}解釋自行宣示「其所為之解釋，自有拘束全國各機關及人民之效力」\footcite{jy185}，2022 年施行的憲法訴訟法第 38 條\ix{statutes}{01@中華民國!02@憲法訴訟法!038@第 38 條}第 1 項\footcite[第 38 條第 1 項]{csfa}則以法律明定判決拘束各機關及人民。依兩人 459 至 463 頁先於憲法\ix{terms}{03@\zhlidxcat{憲法之效力基礎}!先於憲法的憲法}的憲法之說，中華民國憲法法庭的解釋權威，其根據在憲法第 78 條\ix{statutes}{01@中華民國!01@中華民國憲法!078@第 78 條}\footcite[第 78 條]{roc-constitution}、釋字第 185 號解釋、憲法訴訟法，或各機關與人民實際的接受？
\item
  Kramer 633 頁說，在司法至上之下，自制\ix{terms}{05@\zhlidxcat{解釋理論與自制}!司法自制}只剩不致引發政治反撲的審慎\ix{terms}{05@\zhlidxcat{解釋理論與自制}!審慎}理由。釋字第 328 號\ix{cases}{01@中華民國!01@司法院大法官!0328@釋字第 328 號}解釋\footcite{jy328}以固有疆域的界定為「重大之政治問題」而不予解釋\ix{terms}{06@\zhlidxcat{審查制度之類型}!不予解釋}。這種不予解釋，屬於 Kramer 所說共和派\ix{terms}{01@\zhlidxcat{司法審查權限之歸屬}!共和派}或 Thayer\ix{persons}{thayer@Thayer, James Bradley} 式以其他機關另有解釋權威為前提的自制，或屬於審慎的權宜？
\item
  Alexander 與 Schauer 482 頁註 78 與 W08 Tushnet\ix{persons}{tushnet@Tushnet, Mark} 的結論都說，最高法院與國會何者較能定紛，需要尚未進行的經驗研究。若要做這種研究，應比較哪些指標：裁判被遵守的程度、見解的穩定、立法的修正、公眾的信任？證據不足時，應由主張司法至上者或主張分權解釋者負擔論證？
\item
  Ferejohn\ix{persons}{ferejohn@Ferejohn, John} 與 Pasquino\ix{persons}{pasquino@Pasquino, Pasquale} 33 頁說 Kelsen 型\ix{terms}{06@\zhlidxcat{審查制度之類型}!Kelsen 型模式}法院傾向以單一聲音發言，很少公開聽審或公布不同意見\ix{terms}{04@\zhlidxcat{法院之功能}!不同意見與協同意見}書。大法官解釋與憲法法庭判決都附公布協同意見書與不同意見書。依兩人的三種類型，中華民國憲法法庭的審議較接近 Kelsen\ix{persons}{kelsen@Kelsen, Hans（凱爾森）} 型或美國最高法院？公布意見書對 Ferejohn 與 Pasquino 所說的合法性負擔\ix{terms}{07@\zhlidxcat{憲法審判之審議}!合法性負擔}有何影響？
\item
  Ferejohn 與 Pasquino 21 頁把 Kelsen 型法院的擴散歸因於後威權\ix{terms}{06@\zhlidxcat{審查制度之類型}!後威權政體}國家對法官或議會的不信任。中華民國的大法官設於 1948 年行憲之時，在 1990 年代民主化期間作成釋字第 261 號\ix{cases}{01@中華民國!01@司法院大法官!0261@釋字第 261 號}解釋\footcite{jy261}等裁判。兩人對後威權憲法審判的分析\pg{31 至 32 頁}，能否適用於中華民國大法官在民主轉型時期的角色？
\end{enumerate}

% 2026-11-09 站主問本週三篇對財政法、稅捐法的意義與美國法上的聯想；底稿同日記入 notes/docs/1151/seminar/問答.md。
% 美國判決照既有體例行內引，不入 refs.bib。案號、卷頁、日期 2026-11-09 經 agent 核對（LII、聯邦憲法法院、憲法委員會、憲法法庭網站）；
% Loper Bright、Mayo、NFIB 之卷頁與 Golsen 僅經 Wikipedia 核；IRS nonacquiescence 之效果未取得 IRS 官方說明；BVerfGE 93, 121 主文未取原文。
\chapter{與財政法、稅捐法的關係}
\label{ch:55_fiscal}

\begin{qaquestion}
Alexander\ix{persons}{alexander@Alexander, Larry} 與 Schauer\ix{persons}{schauer@Schauer, Frederick}、Kramer\ix{persons}{kramer@Kramer, Larry D.}、Ferejohn\ix{persons}{ferejohn@Ferejohn, John} 與 Pasquino\ix{persons}{pasquino@Pasquino, Pasquale} 三篇，對財政法與稅捐法有何意義？美國法上有哪些可以對照的案件與制度？
\end{qaquestion}

三篇皆未以租稅為例；以下所引租稅案件與制度，均為三篇以外之對照。

\section{一、Alexander 與 Schauer：解釋權威與定紛止爭}
\label{sec:55_fiscal-1}

\begin{enumerate}
\tightlist
\item
  兩人主張其他部門應服從最高法院對憲法的解釋，理由在於權威性解釋所提供的定紛止爭（settlement\ix{terms}{02@\zhlidxcat{定紛止爭論}!定紛止爭}）利益：穩定、可預測與協調（457 至 463 頁）。批評者無人主張下級法院無須遵循最高法院，也無人支持 Meese\ix{persons}{meese@Meese, Edwin, III} 主張的各州得將最高法院裁判視為只拘束該案當事人的立場（476 頁）。\footcite[457--463, 476]{alexander2000}
\item
  美國聯邦稅務局（IRS）對其敗訴的稅務法院（Tax Court）判決，於 Internal Revenue Bulletin 公告 acquiescence（接受）或 nonacquiescence（不接受）；公告不接受者，該局在其他納稅義務人的案件中仍維持原有見解。稅務法院自身則依 Golsen v. Commissioner\ix{cases}{02@美國!03@稅務法院!1970Golsen v. Commissioner（1970）@Golsen v. Commissioner（1970）}, 54 T.C. 742 (1970) 遵循該案上訴所屬巡迴上訴法院的先例。兩項制度都以下級法院的法律解釋為對象，與 Cooper 規則\ix{terms}{01@\zhlidxcat{司法審查權限之歸屬}!Cooper 規則}所涉最高法院的憲法解釋不在同一層次；兩人所稱的定紛止爭利益，在租稅行政上須與全國課稅之一致性相權衡，兩人文中未處理此項權衡。
\item
  Loper Bright Enterprises v. Raimondo\ix{cases}{02@美國!01@聯邦最高法院!2024Loper Bright Enterprises v. Raimondo（2024）@Loper Bright Enterprises v. Raimondo（2024）}, 603 U.S. 369 (2024) 廢棄 Chevron 原則，法律之解釋改由法院獨立判斷，不再尊重\ix{terms}{05@\zhlidxcat{解釋理論與自制}!尊重}行政機關之合理解釋。Mayo Foundation for Medical Education \& Research v. United States, 562 U.S. 44 (2011) 曾明示 Chevron 原則適用於財政部之稅法規則。Loper Bright 所涉為法律解釋，兩人所論為憲法解釋；兩人對解釋權威歸屬的論證能否延伸至法律解釋，是可以提出的問題。
\item
  中華民國：釋字第 217 號\ix{cases}{01@中華民國!01@司法院大法官!0217@釋字第 217 號}解釋以憲法第 19 條\ix{statutes}{01@中華民國!01@中華民國憲法!019@第 19 條}\footcite[第 19 條]{roc-constitution}為據，界定人民僅依法律所定之納稅主體、稅目、稅率、納稅方法及納稅期間等項負納稅義務。\footcite{jy217}大法官解釋與憲法法庭判決拘束各機關及人民（釋字第 185 號\ix{cases}{01@中華民國!01@司法院大法官!0185@釋字第 185 號}解釋；憲法訴訟法第 38 條\ix{statutes}{01@中華民國!02@憲法訴訟法!038@第 38 條}第 1 項）。依兩人 457 頁的循環論證之說，此一拘束力以法律明定之後，是否仍須另以定紛止爭利益說明其正當性，可列為課前問題。
\end{enumerate}

\section{二、Kramer：自制的兩種理由與司法至上的興起}
\label{sec:55_fiscal-2}

\begin{enumerate}
\tightlist
\item
  Kramer\ix{persons}{kramer@Kramer, Larry D.} 626 至 630 頁敘述重建之後司法至上\ix{terms}{01@\zhlidxcat{司法審查權限之歸屬}!司法至上}強勢重現，舉 Brewer 大法官 1893 年的演說，並稱進步派與法院的衝突延續至新政危機，法院在 1930 年代對新政立法「tried, too late, to draw a line」（630 頁）。\footcite[626--630]{kramer2012}
\item
  Pollock\ix{cases}{02@美國!01@聯邦最高法院!1895Pollock v. Farmers' Loan \& Trust Co.（1895）@Pollock v. Farmers' Loan \& Trust Co.（1895）} v. Farmers' Loan \& Trust Co., 157 U.S. 429, 158 U.S. 601 (1895) 宣告 1894 年所得稅法違憲，理由為所得稅屬直接稅而未依各州人口分配。國會與各州其後以第十六修正案（1913 年生效）准許聯邦課徵所得稅而不依人口分配。Kramer 未引 Pollock，該案作成於 Kramer 所述 1893 年以後法院與進步派衝突的期間。推翻 Pollock 的途徑是修憲\ix{terms}{03@\zhlidxcat{憲法之效力基礎}!修憲}，各部門並未拒絕執行該判決；Alexander\ix{persons}{alexander@Alexander, Larry} 與 Schauer\ix{persons}{schauer@Schauer, Frederick} 466 頁所提第 5 條修憲程序，即此一途徑。
\item
  新政時期，United States v. Butler\ix{cases}{02@美國!01@聯邦最高法院!1936United States v. Butler（1936）@United States v. Butler（1936）}, 297 U.S. 1 (1936) 宣告農業調整法以加工稅支應之生產調整計畫違憲；Steward Machine Co.~v. Davis, 301 U.S. 548 (1937) 維持社會安全法的失業保險稅。兩案相隔一年，正落在 Kramer 630 頁所述法院對新政劃線而後退讓的時期。
\item
  National Federation of Independent Business v. Sebelius\ix{cases}{02@美國!01@聯邦最高法院!2012NFIB v. Sebelius（2012）@NFIB v. Sebelius（2012）}, 567 U.S. 519 (2012) 以課稅權維持健保法個人強制納保之規定，與 Kramer 一文同年。Kramer 633 頁說，在司法至上之下，自制\ix{terms}{05@\zhlidxcat{解釋理論與自制}!司法自制}只剩不致引發政治反撲的審慎\ix{terms}{05@\zhlidxcat{解釋理論與自制}!審慎}理由。該案以課稅權為依據而維持法律，屬於 Kramer 所說共和派\ix{terms}{01@\zhlidxcat{司法審查權限之歸屬}!共和派}以其他機關另有解釋權威為前提的自制，或屬審慎的權宜，可依 Kramer 的分類提問。
\end{enumerate}

\section{三、Ferejohn 與 Pasquino：審議型態與立法者角色}
\label{sec:55_fiscal-3}

\begin{enumerate}
\tightlist
\item
  兩人說明法國憲法委員會於法律公布前審查，延遲通常為一個月，政府宣告緊急時縮短為八日，審查程序承受時間壓力（33 頁）。\footcite[33]{ferejohn2002}財政法案另受第五共和憲法第 47 條所定議決期限之限制。
\item
  兩人歸納 Kelsen\ix{persons}{kelsen@Kelsen, Hans（凱爾森）} 的觀念：憲法法院宣告法律違憲，須重建該法律之前的法律狀態，因而兼有積極立法的作用（30 頁，註 13）。德國聯邦憲法法院於租稅案件有宣告法律與基本法不符而不宣告無效、並定期限命立法者修正之例，如財產稅裁判（BVerfGE 93, 121，1995 年）與不動產稅（Grundsteuer）裁判（BVerfGE 148, 147，2018 年）。租稅法律一經宣告無效即生稅收短缺，立法者須另定替代規定，法院之裁判因此兼及立法上的安排。
\end{enumerate}


\appendix
\IfFileExists{Teile/A1_timeline.tex}{% 時間線：圖在 figures/timeline.tex；事件之年份與所在節次見各章。
\chapter{相關事件年表}
\label{app:timeline}

圖~\ref{fig:timeline}依年份列出三篇本文所引之美國判決與著作、法國與義大利之立法，以及第五部所涉之中華民國司法院大法官解釋與憲法訴訟法，各列其法域。各事件之刊本頁次見各節，出處見附錄~\ref{app:sources}。

\begin{qafigure}[相關事件年表]{三篇本文所涉之事件，依年份與法域。\label{fig:timeline}}
% 三篇本文所涉之事件：直列單軌，版型照 ../../W06_Waldron/figures/timeline.tex，另加法域一欄，排在類別之後（2026-11-09 站主：「哪個國啊，不標清楚嗎」）。
% 年份逐一回 Teile 正文核過（2026-11-09）；正文未載年份者（Dred Scott 1857、Brown 1954）依案號通行年份，
% 釋字三號依 refs.bib 之解釋日期。著作之法域依作者任教之國；Ferejohn 與 Pasquino 為比較法著作。本文三則用強調色。
\begin{tikzpicture}[qa,
    yr/.style={font=\footnotesize,text=#1,anchor=north east,inner sep=0pt},
    cat/.style={font=\scriptsize,text=dimink,anchor=north west,inner sep=0pt},
    ev/.style={font=\footnotesize,text=#1,anchor=north west,inner sep=0pt,text width=98mm,
      execute at begin node={\raggedright\setstretch{1.2}}},
    dot/.style={circle,fill=#1,draw=white,line width=0.8pt,inner sep=0pt,minimum size=5.5pt}]
  \coordinate (prev) at (0,0.25);
  \node[cat,anchor=south east] at (14mm,0.15) {年份};
  \node[cat,anchor=south west] at (19mm,0.15) {類別};
  \node[cat,anchor=south west] at (29mm,0.15) {法域};
  \node[cat,anchor=south west] at (44mm,0.15) {事件};
  \foreach \k/\d/\j/\c/\t [count=\i] in {
      n/1787/美國/函件/{Iredell 致 Spaight 函（北卡羅來納州）：違憲而無可爭辯者始得宣告無效},
      n/1793/美國/判決/{維吉尼亞州 General Court，Kamper v.\ Hawkins：Nelson 法官論憲法之於政府部門},
      n/1800/美國/選舉/{總統選舉，共和派勝出},
      n/1803/美國/判決/{聯邦最高法院，Marbury v.\ Madison},
      n/1857/美國/判決/{聯邦最高法院，Dred Scott v.\ Sandford},
      n/1893/美國/著作/{Thayer, The Origin and Scope of the American Doctrine of Constitutional Law},
      n/1948/中華民國/機關/{司法院大法官設立},
      n/1953/義大利/立法/{1953 年 3 月 11 日第 87 號法：憲法法院判決應載明理由},
      n/1954/美國/判決/{聯邦最高法院，Brown v.\ Board of Education},
      n/1958/美國/判決/{聯邦最高法院，Cooper v.\ Aaron},
      n/1958/法國/立法/{1958 年 11 月 7 日第 58-1067 號條例：憲法委員會之組織法},
      n/1984/中華民國/解釋/{司法院大法官釋字第 185 號：解釋拘束全國各機關及人民},
      n/1990/中華民國/解釋/{司法院大法官釋字第 261 號},
      n/1993/美國/著作/{Rawls, Political Liberalism},
      n/1993/中華民國/解釋/{司法院大法官釋字第 328 號：固有疆域之界定屬重大之政治問題},
      n/1997/美國/判決/{聯邦最高法院，City of Boerne v.\ Flores},
      n/1997/美國/著作/{Alexander 與 Schauer, On Extrajudicial Constitutional Interpretation},
      n/1998/美國/著作/{Tushnet, Two Versions of Judicial Supremacy},
      n/1999/美國/著作/{Tushnet, Taking the Constitution Away from the Courts},
      k/2000/美國/本文/{Alexander 與 Schauer, Defending Judicial Supremacy: A Reply},
      k/2002/比較法/本文/{Ferejohn 與 Pasquino, Constitutional Courts as Deliberative Institutions：法國、德國、義大利、西班牙、美國},
      n/2003/美國/著作/{Tushnet, Non-Judicial Review},
      n/2004/美國/著作/{Kramer, The People Themselves},
      n/2012/美國/著作/{Posner, The Rise and Fall of Judicial Self-Restraint},
      k/2012/美國/本文/{Kramer, Judicial Supremacy and the End of Judicial Restraint},
      n/2022/中華民國/立法/{憲法訴訟法施行：憲法法庭判決拘束各機關及人民}}{
    \if k\k\def\ink{figE}\else\def\ink{caseink}\fi
    \path (prev) ++(0,-0.25) coordinate (top);
    \node[ev=\ink] (e\i) at (44mm,0 |- top) {\t};
    \node[cat] at (19mm,0 |- top) {\c};
    \node[cat] at (29mm,0 |- top) {\j};
    \node[yr=\ink] at (14mm,0 |- top) {\d};
    \node[dot=\ink] at ($(16.5mm,0 |- top)+(0,-1.4mm)$) {};
    \coordinate (prev) at (0,0 |- e\i.south);
    \xdef\last{\i}
  }
  \begin{scope}[on background layer]
    \draw[ruleink,line width=0.8pt] (16.5mm,-0.6mm) -- (16.5mm,0 |- e\last.south);
  \end{scope}
\end{tikzpicture}

\end{qafigure}
}{}
\IfFileExists{Teile/A2_sources.tex}{% 本檔由 tools/build_sources.py 自各章產生，每次建置整份重寫，不要手改。
\chapter{各節出處一覽}
\label{app:sources}

表~\ref{tab:sources}依各部之序，列出各節所據之出處：第二至第四部各節先列該部所讀之本文及該節所述之刊本頁次，其次為該節另引之文獻，依第一次引用之先後排列，頁碼照該節所引。完整書目見參考文獻。

\begin{small}\setstretch{1.3}
\setlength{\LTpre}{6pt}
\def\LTcaptype{table}
\begin{longtable}{@{}>{\raggedright\arraybackslash}p{10mm}>{\raggedright\arraybackslash}p{62mm}>{\raggedright\arraybackslash}p{\dimexpr\linewidth-72mm-4\tabcolsep\relax}@{}}
\caption{各節之出處}\label{tab:sources}\\
\toprule
節 & 標題 & 出處 \\
\midrule
\endfirsthead
\toprule
節 & 標題 & 出處 \\
\midrule
\endhead
\bottomrule
\endlastfoot
\multicolumn{3}{@{}l}{導論：三篇的出處與問題} \\* \addlinespace[2pt]
\ref{sec:intro-1} & 三篇的出處與範圍 & --- \\
\ref{sec:intro-2} & 三篇各自處理的問題 & \citeauthor{alexander2000}，\citetitle{alexander2000} \newline \citeauthor{alexander1997}，\citetitle{alexander1997} \newline \citeauthor{kramer2012}，\citetitle{kramer2012} \newline \citeauthor{posner2012}，\citetitle{posner2012} \newline \citeauthor{thayer1893}，\citetitle{thayer1893} \newline \citeauthor{ferejohn2002}，\citetitle{ferejohn2002} \\
\ref{sec:intro-3} & 與 W08 的銜接 & \citeauthor{tushnet2003}，\citetitle{tushnet2003}，註 209 \newline \citeauthor{tushnet1998}，\citetitle{tushnet1998} \newline \citeauthor{tushnet1999}，\citetitle{tushnet1999} \\
\addlinespace[6pt]
\multicolumn{3}{@{}l}{Alexander 與 Schauer〈Defending Judicial Supremacy: A Reply〉} \\* \addlinespace[2pt]
\ref{sec:AS-1} & 導言：先前論點、批評與本文目的 & \citeauthor{alexander2000}，\citetitle{alexander2000}，455--458 \\
\ref{sec:AS-2} & 第 I 部分：問題的性質 & \citeauthor{alexander2000}，\citetitle{alexander2000}，459--467 \\
\ref{sec:AS-3} & 第 II 部分：定紛止爭功能及其最適所在 & \citeauthor{alexander2000}，\citetitle{alexander2000}，467--478 \\
\ref{sec:AS-4} & 第 III 部分：最高法院的功能 & \citeauthor{alexander2000}，\citetitle{alexander2000}，478--481 \\
\ref{sec:AS-5} & 結論 & \citeauthor{alexander2000}，\citetitle{alexander2000}，481--482 \\
\addlinespace[6pt]
\multicolumn{3}{@{}l}{Kramer〈Judicial Supremacy and the End of Judicial Restraint〉} \\* \addlinespace[2pt]
\ref{sec:Kramer-1} & 摘要與論題 & \citeauthor{kramer2012}，\citetitle{kramer2012}，621--622 \\
\ref{sec:Kramer-2} & 建國時期的兩種司法審查理論 & \citeauthor{kramer2012}，\citetitle{kramer2012}，622--624 \\
\ref{sec:Kramer-3} & 解釋方法與自制 & \citeauthor{kramer2012}，\citetitle{kramer2012}，624--626 \\
\ref{sec:Kramer-4} & 1800 年至新政 & \citeauthor{kramer2012}，\citetitle{kramer2012}，626--630 \newline \citeauthor{kramer2004}，\citetitle{kramer2004} \\
\ref{sec:Kramer-5} & Brown 案之後的自由派與保守派 & \citeauthor{kramer2012}，\citetitle{kramer2012}，630--632 \\
\ref{sec:Kramer-6} & 解釋理論的興起與自制的衰退 & \citeauthor{kramer2012}，\citetitle{kramer2012}，632--634 \\
\addlinespace[6pt]
\multicolumn{3}{@{}l}{Ferejohn 與 Pasquino〈Constitutional Courts as Deliberative Institutions〉} \\* \addlinespace[2pt]
\ref{sec:FP-1} & 導論：Kelsen 型模式的擴散與合法性負擔 & \citeauthor{ferejohn2002}，\citetitle{ferejohn2002}，21--22 \newline \citeauthor{rawls1993}，\citetitle{rawls1993}，231--236 \\
\ref{sec:FP-2} & 審議的期待 & \citeauthor{ferejohn2002}，\citetitle{ferejohn2002}，22--25 \\
\ref{sec:FP-3} & 民主制度中的審議 & \citeauthor{ferejohn2002}，\citetitle{ferejohn2002}，25--27 \\
\ref{sec:FP-4} & 憲法法院在現代民主中的位置 & \citeauthor{ferejohn2002}，\citetitle{ferejohn2002}，28--33 \\
\ref{sec:FP-5} & 憲法法院的審議 & \citeauthor{ferejohn2002}，\citetitle{ferejohn2002}，33--36 \\
\addlinespace[6pt]
\multicolumn{3}{@{}l}{與前後幾週的關係及課前問題} \\* \addlinespace[2pt]
\ref{sec:relations-1} & 與 W04 至 W08、W10 的對照 & \citeauthor{waldron2006}，\citetitle{waldron2006}，1350 \\
\ref{sec:relations-2} & 課前可以先想的問題 & \citefield{jy185}{title} \newline \citefield{csfa}{title}，第 38 條第 1 項 \newline \citefield{roc-constitution}{title}，第 78 條 \newline \citefield{jy328}{title} \newline \citefield{jy261}{title} \\
\addlinespace[6pt]
\multicolumn{3}{@{}l}{與財政法、稅捐法的關係} \\* \addlinespace[2pt]
\ref{sec:55_fiscal-1} & 一、Alexander 與 Schauer：解釋權威與定紛止爭 & \citeauthor{alexander2000}，\citetitle{alexander2000}，457--463, 476 \newline \citefield{roc-constitution}{title}，第 19 條 \newline \citefield{jy217}{title} \\
\ref{sec:55_fiscal-2} & 二、Kramer：自制的兩種理由與司法至上的興起 & \citeauthor{kramer2012}，\citetitle{kramer2012}，626--630 \\
\ref{sec:55_fiscal-3} & 三、Ferejohn 與 Pasquino：審議型態與立法者角色 & \citeauthor{ferejohn2002}，\citetitle{ferejohn2002}，33 \\
\addlinespace[6pt]
\end{longtable}
\end{small}
}{}
% 自 notes/docs/1151/seminar/W09_1109_預習.md（87808f3）由 tools/md_to_tex.py 於 2026-11-09 轉出一次；此後本檔為正本。
\chapter{主要用語對照}

\begin{longtable}[]{@{}
  >{\raggedright\arraybackslash}p{(\linewidth - 6\tabcolsep) * \real{0.3700}}
  >{\raggedright\arraybackslash}p{(\linewidth - 6\tabcolsep) * \real{0.2700}}
  >{\raggedright\arraybackslash}p{(\linewidth - 6\tabcolsep) * \real{0.1300}}
  >{\raggedright\arraybackslash}p{(\linewidth - 6\tabcolsep) * \real{0.2300}}@{}}
\caption{三篇的主要用語。篇目欄 AS 為 Alexander 與 Schauer，FP 為 Ferejohn 與 Pasquino。}\tabularnewline
\toprule\noalign{}
\begin{minipage}[b]{\linewidth}\raggedright
原文
\end{minipage} & \begin{minipage}[b]{\linewidth}\raggedright
中譯
\end{minipage} & \begin{minipage}[b]{\linewidth}\raggedright
篇目
\end{minipage} & \begin{minipage}[b]{\linewidth}\raggedright
出處（頁次）
\end{minipage} \\
\midrule\noalign{}
\endfirsthead
\toprule\noalign{}
\begin{minipage}[b]{\linewidth}\raggedright
原文
\end{minipage} & \begin{minipage}[b]{\linewidth}\raggedright
中譯
\end{minipage} & \begin{minipage}[b]{\linewidth}\raggedright
篇目
\end{minipage} & \begin{minipage}[b]{\linewidth}\raggedright
出處（頁次）
\end{minipage} \\
\midrule\noalign{}
\endhead
\bottomrule\noalign{}
\endlastfoot
judicial supremacy & 司法至上 & AS、Kramer & 標題、456；Kramer 621、623、633 \\
interpretive supremacy & 解釋至上 & AS & 457、459 \\
settlement（for settlement\textquotesingle s sake） & 定紛止爭（為定紛而定紛） & AS & 467、469 \\
first-order / second-order settlement & 一階定紛／二階定紛 & AS & 457、467 \\
content-independent reasons & 與內容無關之理由 & AS & 457 註 11、469 \\
preconstitutional & 先於憲法的 & AS & 459--463 \\
grundnorm / rule of recognition & 基本規範／承認規則 & AS & 460 註 17--18 \\
social fact（acceptance as） & 社會事實（作為社會事實的接受） & AS & 465 \\
Cooper rule & Cooper 規則 & AS & 466 \\
institutional anarchy & 制度上的無政府狀態 & AS & 466、473 \\
asymmetry of authority & 權威的不對稱 & AS & 472 註 50、474 註 51 \\
case-decider / law-maker & 案件裁判者／造法者 & AS & 478--479 \\
dispute settler / constitutional interpreter & 爭端解決者／憲法解釋者 & AS & 479 \\
advisory opinion & 諮詢意見 & AS & 481 \\
stare decisis & 遵循先例 & AS & 477、479 \\
judicial (self-)restraint／deference & 司法自制／尊重 & Kramer & 621、624--625 \\
popular constitutionalism & 人民憲政主義 & Kramer & 622（註 3） \\
Republican／Jeffersonian view & 共和派（Jeffersonian）觀點 & Kramer & 622 \\
Federalist & 聯邦派 & Kramer & 623--626 \\
interpretive authority & 解釋權限（詮釋權威） & Kramer & 622、623、632 \\
prudence／expediency & 審慎／權宜 & Kramer & 626、633 \\
countermajoritarian difficulty & 反多數決困難 & Kramer & 631 \\
interpretive theory／methodology & 解釋理論／解釋方法 & Kramer & 621、632 \\
equilibrium point & 均衡點 & Kramer & 634 \\
principal／agent & 本人／代理人 & Kramer & 625--626 \\
\en{Lochner} era & \en{Lochner} 時代 & Kramer & 630、631 \\
Kelsenian type / Kelsenian model & Kelsen 型（模式） & FP & 21；30；33 \\
burden of legitimation & 合法性負擔 & FP & 21 \\
deliberative expectations & 審議的期待 & FP & 22 \\
public reasons & 公共理由 & FP & 23（Rawls 的語彙） \\
internal deliberation / internally deliberative & 內部審議 & FP & 23；35；36 \\
externally (deliberative) & 外部審議 & FP & 35 \\
parliamentary sovereignty regimes & 議會主權體制 & FP & 28 \\
division of power / separation of powers regimes & 權力分立體制 & FP & 29；33 \\
a priori review & 事前審查（先於公布） & FP & 28；30 \\
abstract review & 抽象審查 & FP & 28；30 \\
concrete (a posteriori) review & 具體（事後）審查 & FP & 31 \\
dispersed review / concentrated review & 分散審查／集中審查 & FP & 29；31；32 \\
bloc de constitutionnalité & 憲法整體 & FP & 28 註 9；34 註 19 \\
negative legislation & 消極立法 & FP & 30 註 13 \\
reason free zone & 無須理由的領域 & FP & 26 \\
countermajoritarian & 反多數決 & FP & 25 \\
\end{longtable}


\backmatter
\printbibliography[heading=bibintoc]
\IfFileExists{Teile/90_index_xref.tex}{% 由 tools/build_index.py 自 index-terms.tsv 產生，勿手改。
\ix{terms}{01@\zhlidxcat{司法審查權限之歸屬}!解釋至上|seealso{司法至上}}
\ix{terms}{01@\zhlidxcat{司法審查權限之歸屬}!共和派|seealso{人民憲政主義}}
\ix{terms}{01@\zhlidxcat{司法審查權限之歸屬}!Cooper 規則|seealso{司法至上}}
\ix{terms}{05@\zhlidxcat{解釋理論與自制}!尊重|seealso{司法自制}}
\ix{terms}{07@\zhlidxcat{憲法審判之審議}!內部審議|seealso{外部審議}}
\ix{terms}{02@\zhlidxcat{定紛止爭論}!settlement|see{定紛止爭}}
\ix{terms}{01@\zhlidxcat{司法審查權限之歸屬}!違憲審查|see{司法審查}}
}{}
\printindex[terms]
\printindex[statutes]
\printindex[cases]
\printindex[persons]

\end{document}
